\documentclass[english,12pt]{article}
%%%%%%%%%%%%%%%%%%%%%%%%%%%%%%%%%%%%%%%%%%%%%%%%%%%%%%%%%%%%%%%%%%%%%%%%%%%%%%%%%%%%%%%%%%%%%%%%%%%%%%%%%%%%%%%%%%%%%%%%%%%%%%%%%%%%%%%%%%%%%%%%%%%%%%%%%%%%%%%%%%%%%%%%%%%%%%%%%%%%%%%%%%%%%%%%%%%%%%%%%%%%%%%%%%%%%%%%%%%%%%%%%%%%%%%%%%%%%%%%%%%%%%%%%%%%
\usepackage[letterpaper,tmargin=1in,bmargin=1in,lmargin=1in,rmargin=1in]{geometry}
\usepackage[dvips]{graphicx}
\usepackage{multirow}
\usepackage{setspace}
\usepackage{sectsty}
\usepackage{natbib}
\usepackage[T1]{fontenc}
\usepackage{lmodern}
\usepackage{lscape}
\usepackage{marginnote}
\usepackage{amssymb,amsfonts,amsmath,latexsym,setspace,amsthm}
\usepackage{microtype}
\usepackage{bm}
\usepackage[multiple, stable, bottom]{footmisc}
\usepackage{color}
\usepackage{xcolor}
\usepackage[normalem]{ulem}
\usepackage[colorlinks,breaklinks,               bookmarks=false,
    pdfstartview=Fit,      pdfview=Fit,           linkcolor = ChadBlue,
    urlcolor  = ChadGreen,
    citecolor = ChadGreen,
    hyperfootnotes=false
    ]{hyperref}
\usepackage[figure,table]{hypcap}
\usepackage{booktabs}
\usepackage{titlesec}
\usepackage{subcaption}
\usepackage[export]{adjustbox}
\usepackage[capposition=top]{floatrow}
\usepackage[format=hang,font=normalsize,labelfont=bf]{caption}
\usepackage[para,online,flushleft]{threeparttable}
\usepackage{array}
\usepackage{babel}
\usepackage{enumitem}
\usepackage{tabularx}
\usepackage{soul}
\usepackage[T1]{fontenc}
\usepackage[latin9]{inputenc}
\usepackage{geometry}
\usepackage{color}
\usepackage{babel}
\usepackage{verbatim}
\usepackage{prettyref}
\usepackage{float}
\usepackage{amsthm}
\usepackage{amsmath,amsbsy}
\usepackage{amssymb}
\usepackage{bm}
\usepackage{graphicx}
\usepackage{setspace}
\usepackage{esint}
\usepackage{titletoc}
\usepackage{pdfpages}
\usepackage{xcolor}
\usepackage{colortbl}
\usepackage{xargs}
\usepackage[colorinlistoftodos,prependcaption,textsize=tiny]{todonotes}
\usepackage{epstopdf}
\usepackage{lmodern}
\usepackage{amsfonts}
\usepackage{mathabx}
\usepackage{bbm}
\usepackage{booktabs}
\usepackage{appendix}
\usepackage{caption}
\usepackage{subcaption}
\usepackage[flushleft]{threeparttable}
\usepackage{sgame}
\usepackage{multirow,array}
\usepackage{bookmark}
\usepackage{nicefrac,xfrac}
\usepackage{enumitem}

\setcounter{MaxMatrixCols}{10}
%TCIDATA{OutputFilter=LATEX.DLL}
%TCIDATA{Version=5.50.0.2960}
%TCIDATA{<META NAME="SaveForMode" CONTENT="1">}
%TCIDATA{BibliographyScheme=BibTeX}
%TCIDATA{LastRevised=Sunday, December 28, 2025 12:01:53}
%TCIDATA{<META NAME="GraphicsSave" CONTENT="32">}

\theoremstyle{definition}
\newtheorem{definition}{Definition}[section]
\newtheorem{assumption}{Assumption}[section]
\newtheorem{condition}{Condition}[section]
\newtheorem{remark}{Remark}[section]
\definecolor{ChadDarkBlue}{rgb}{.1,0,.2}
    \definecolor{ChadBlue}{rgb}{.1,.1,.5}
    \definecolor{ChadRoyal}{rgb}{.2,.2,.8}
    \definecolor{ChadGreen}{rgb}{0,.4,0}
    \definecolor{ChadRed}{rgb}{.5,0,.5}
\newcommand{\ra}[1]{\renewcommand{\arraystretch}{#1}}
\renewcommand{\abstractname}{\color{ChadBlue}Abstract}
\renewcommand{\figurename}{\color{ChadBlue}Figure}
\renewcommand{\tablename}{\color{ChadBlue}Table}
\titleformat*{\section}{\normalfont\Large\bfseries\color{ChadBlue}}
\titleformat*{\subsection}{\normalfont\large\bfseries\color{ChadBlue}}
\titleformat*{\subsubsection}{\normalfont\normalsize\bfseries\color{ChadBlue}}
\newcommand{\Prob}[1]{\mbox{\textnormal{Pr}} \, [#1] }
\newcommand{\SubSubSection}[1]{\subsubsection{#1} \baselineskip=18pt}
\newcommand{\problem}[2]{\hspace{.1in} {{\color{ChadBlue}{\bf #1:} #2}}}
\newcommand{\proptitle}[1]{\color{ChadBlue} \textnormal{(#1):}}
\newtheorem{proposition}{\color{ChadBlue} Proposition}
\newtheorem{fact}{\color{ChadBlue} Fact}
\newtheorem{technicalassumption}{\color{ChadBlue} Technical Assumption}
\newtheorem{conjecture}{\color{ChadBlue} Conjecture}
\newtheorem{lemma}{\color{ChadBlue} Lemma}
\newcommand{\assume}[1]{{\bf{Assumption #1.} }}
\newcommand{\Ex}[1]{\mbox{ }\mathbb{E}}
\newcommand{\Proof}[2]{{\hspace{-\parindent} {\color{ChadBlue}\bf Proof of Proposition}~\ref{#1}.}
{\color{ChadBlue} #2} \vspace{.1in}}
\let\oldref\ref
\renewcommand{\ref}[1]{(\oldref{#1})}
\onehalfspacing
\geometry{verbose,tmargin=2.33cm,bmargin=2.33cm,lmargin=2.33cm,rmargin=2.33cm}
\definecolor{ceruleanblue}{rgb}{0.16, 0.32, 0.75}
\definecolor{cerise}{rgb}{0.87, 0.19, 0.39}
\definecolor{ao(english)}{rgb}{0.0, 0.5, 0.0}
\definecolor{antiquebrass}{rgb}{0.8, 0.58, 0.46}
\definecolor{arylideyellow}{rgb}{0.91, 0.84, 0.42}
\newcommand\rr[1]{\textcolor{red}{#1}}
\newcommandx{\info}[2][1=]{\todo[linecolor=ao(english),backgroundcolor=ao(english)!25,bordercolor=ao(english),#1]{#2}}
\newcommandx{\comments}[2][1=]{\todo[linecolor=cerise,backgroundcolor=cerise!25,bordercolor=cerise,#1]{#2}}
\makeatletter
\providecommand{\tabularnewline}{\\}
\hypersetup{
    linkcolor=cerise,              citecolor=ceruleanblue,            filecolor=magenta,          urlcolor=ceruleanblue           }
\newrefformat{ass}{Assumption \ref{#1}}
\newrefformat{proof}{Appendix \ref{#1}}
\newrefformat{sec}{Section \ref{#1}}
\newrefformat{sub}{Section \ref{#1}}
\newrefformat{fact}{Fact \ref{#1}}
\newrefformat{technicalassumption}{Technical Assumption \ref{#1}}
    \makeatletter
    \pgfdeclarepatternformonly[\hatchdistance,\hatchthickness,\tikz@pattern@color]{flexible hatch}
    {\pgfqpoint{0pt}{0pt}}
    {\pgfqpoint{\hatchdistance}{\hatchdistance}}
    {\pgfpoint{\hatchdistance-1pt}{\hatchdistance-1pt}}
    {
        \pgfsetcolor{\tikz@pattern@color}
        \pgfsetlinewidth{\hatchthickness}
        \pgfpathmoveto{\pgfqpoint{0pt}{0pt}}
        \pgfpathlineto{\pgfqpoint{\hatchdistance}{\hatchdistance}}
        \pgfusepath{stroke}
    }
    \makeatother
\newlist{stages}{enumerate}
{1}
\setlist[stages, 1]{label = t\arabic*.}
\tolerance=1
\emergencystretch=\maxdimen
\hyphenpenalty=10000
\hbadness=10000
\makeatother
  \providecommand{\propositionname}{Proposition}

\usepackage{array}
\newcolumntype{L}[1]{>{\raggedright\let\newline\\\arraybackslash\hspace{0pt}}m{#1}}
\newcolumntype{C}[1]{>{\centering\let\newline\\\arraybackslash\hspace{0pt}}m{#1}}
\newcolumntype{R}[1]{>{\raggedleft\let\newline\\\arraybackslash\hspace{0pt}}m{#1}}
\usepackage[commandnameprefix=always]{changes}


\begin{document}

\title{\color{ChadBlue}{
\textbf{The Remains of the Trade: \linebreak The U.S.--China Trade War and its Aftermath}}\thanks{
A preliminary version of this paper was prepared for the 2nd Annual AlUla Conference for Emerging Market Economies, held in AlUla, Saudi Arabia on February 8--9, 2026. We are grateful to Yipei Zhang for research assistance, and to Roberto Cardarelli, Shushanik Hakobyan and Neil Meads for feedback. The views expressed here are those of the authors and should not be attributed to the International Monetary Fund, its Executive Board, or management.}}


\author{%
\vspace*{15pt}\\[-1em]
\begin{tabular}{@{}C{5cm}C{5cm}C{6cm}@{}}
Pol Antr{\`a}s & Adrian Kulesza & Andrea F. Presbitero  \\
\small{Harvard University and NBER} & \small{Harvard University} & \small{International Monetary Fund and CEPR} \\
\end{tabular}
}

\date{\vspace{15pt} \normalsize June 2026}
\maketitle

\begin{abstract} \noindent
We study how the U.S.--China trade conflict has reshaped global trade patterns with significant spillovers 
to third-countries. Using recent tariff policy changes and monthly international trade flow data, we document 
a redirection of Chinese exports away from the U.S. market and toward third--country economies. To interpret 
the impact of these reallocations on third countries, we introduce a taxonomy that separates four exposure 
margins -- export competition, import competition, consumer opportunities, and sourcing opportunities -- and 
construct theoretically-derived indices that map countries into these channels. 
\end{abstract}

\newpage
\medskip

\setcounter{page}{1}

\section{Introduction}

Talk of `deglobalization' has become a fixture of policy debates since the
mid-2010s, and it has intensified in the wake of the U.S.--China trade war
and the COVID-19 shock. Yet, as many authors have emphasized %
\citep{Antras2021DeglobalisationECB,goldberg2023global}, there is little
systematic evidence that, at the global level, international trade has become 
a declining share of world consumption or income. Figure \ref{fig: world_trade} 
plots the ratio of world trade (exports plus imports) over world GDP in the 
period 1970-2024. It is true that after having grown enormously during the late 
1980s, 1990s and early 2000s, this ratio has remained flat since the Great 
Financial Crisis (GFC). Nevertheless, the world is not deglobalizing, it has 
just ceased to globalize \emph{more}.

\begin{figure}[htbp]
\caption{World Trade Over World GDP (1970-2024)}
\label{fig: world_trade}\centering
\par
    \centering
    \includegraphics[scale=0.7]{Figures/trade_GDP_income.png}
\bigskip
\par
\begin{minipage}{.95\textwidth}
\footnotesize 
\textbf{Notes:} The figure plots world trade as a percentage of GDP for the period 1970-2024.
It decomposes world trade between advanced economies, China, and emerging and developing countries
(ex. China). For each region, trade is scaled by world GDP and taken as a
backward-looking 3-year moving average. \par
  \smallskip
\textbf{Sources:} World Bank's World Development Indicators.
\end{minipage}
\end{figure}

At the same time, the recent stability in the world trade-to-GDP ratio masks a notable
reshuffling of cross-border linkages. As panel (b) of Figure \ref{fig:
world_trade} shows, since the GFC, trade openness has fallen markedly more
in advanced economies than in both emerging and developing economies, whether
China is included or excluded. Even more notably, direct trade flows between
the U.S. and China have tumbled in recent years. As illustrated in the left 
panel of Figure \ref{fig: US_China_shares}, the share of U.S. imports accounted for 
by China has fallen drastically from 22\% in 2017 to about 8\% in 2026, a level
not recorded since China entered the WTO in 2001. Similarly, China's share 
of exports destined to the U.S. market has fallen from around 20\% in 2017 to 
about 10\% in 2026. Given the sheer size of Chinese exports ($\$3.77$ trillion in 2025),
this entails a redirection of Chinese exports from the U.S. market
to other markets comparable in size to the GDP of certain advanced economies 
such as Chile, Portugal or Greece.
% Trade has also changed in a more
% profound sense, as powerful countries are increasingly flexing their
% geopolitical leverage through unilateral measures that sit uneasily with a
% strictly rules-based, WTO-centered regime.


\begin{figure}[htbp]
\caption{U.S.-China Decoupling}
\label{fig: US_China_shares}\centering
\par
\begin{subfigure}[t]{0.5\linewidth}
    \centering
    \includegraphics[scale=0.4]{Figures/China_market_shares_US.png}
    \caption{China's market share in U.S. imports}
    \label{fig: Trade World}
  \end{subfigure}\hfill 
\begin{subfigure}[t]{0.5\linewidth}
    \centering
  \includegraphics[scale=0.4]{Figures/US_market_shares_China.png}
    \caption{U.S. market share in China's exports}
    \label{fig: Trade Regions}
  \end{subfigure}
\par
\bigskip
\par
\begin{minipage}{.95\textwidth}
\footnotesize 
\textbf{Notes:} Panel (a) plots the share of total U.S. goods imports from China, from US Census. Panel (b) plots the share of total China's goods exports to the U.S. The series is from the CEPII-BACI dataset until 2024 and from TDM for 2025 and 2026. For both panels, data for 2026 refer to January and February. \par
  \smallskip
\textbf{Sources:} US Census, CEPII-BACI dataset \citep{Gaulier-Zignago_2010}, and Trade Data Monitor (TDM).
\end{minipage}
\end{figure}
    

Seen in this light, the key question raised by recent events is not whether
globalization is ending, but how trade and production relationships are
being reorganized across countries, regions and geopolitical blocs, especially
since the onset of the U.S.-China trade war. The defining feature of the post-2018 
period is not an indiscriminate retreat
from trade, but a `Great Reallocation' %
\citep{AlfaroChor2023JacksonHoleGreatReallocation} of trade flows across
partners and products. Much as in Kazuo Ishiguro's \emph{The Remains of the
Day}, the drama here lies less in rupture than in reorganization: continuity
in the aggregate, but a profound reshuffling in who trades what with whom,
and on what terms.


In this paper, we make three contributions that speak to this post-2018
reorganization. First, we offer a fresh new look at the U.S.--China trade
war, shedding new light on how it has developed in recent months --
especially since ``Liberation Day'' -- by drawing on up-to-date data to
document the latest shifts in trade patterns and trade policy intensity. Second,
we zoom in on the effects of the trade war on third countries: we empirically
characterize the spillovers of U.S.--China decoupling for third countries
(especially emerging markets) rather than merely treating them as residual
beneficiaries of diversion. More specifically, we offer theoretically-motivated
diagnostics that separate four exposure margins -- export competition, import 
competition, consumer opportunities, and sourcing opportunities -- to organize the evidence. 
Third, we provide preliminary evidence suggestive that our indices have empirical
bite in predicting how the ongoing U.S.-China trade tensions have spilled over 
to third countries. textbf{[PERHAPS INSTEAD SAY SOMETHING ABOUT OPENING THE 
BLACK BOX OF THE DIAGNOSTICS]}


Our starting point is the influential work of \cite%
{AlfaroChor2023JacksonHoleGreatReallocation} documenting systematic changes
in U.S. sourcing patterns in recent years, involving a reduction in direct exposure to China
and rising shares for alternative suppliers. In section \ref{Sec: Great
Reallocation}, we document some novel aspects of this reallocation with a
specific focus on the implications of the U.S.--China trade war in terms of (i)
a reallocation of U.S. import sources; (ii) a reorientation of Chinese export
destinations; and (iii) broader shifts in global trade patterns. Among other new findings, we find 
that the initial phases of the U.S.--China trade war saw a relatively slow reallocation of 
U.S. imports away from China because of the sluggish adjustment associated with products with 
longer production cycles \citep[cf.][]{antras2025measuring}. 

More importantly, our analysis also shows that the pace of reallocation has 
accelerated since the beginning of the 
second Trump administration and especially since Liberation Day, arguably because
that day's announcements were interpreted as a permanent shock to the U.S.--China trade 
relationship. Conversely, we find less evidence that other countries facing higher 
tariffs after Liberation Day saw their exports to the U.S. significantly reduced,
likely because those other tariffs were surrounded by reversals and uncertainty. The sharp 
reduction of Chinese exports to the U.S. was mirrored by sizable increases in Chinese exports 
to other markets, which generated an intensification of competitive pressures in third markets 
where China and other exporters compete head-to-head.

Section \ref{Sec: Third Countries} turns to a more systematic empirical exploration of the
implications of the U.S.--China trade war for the countries `caught in the
middle'. We develop a practical taxonomy to measure third-country exposure
to U.S.--China decoupling along four (non-mutually-exclusive) channels: (i) 
\emph{export competition}, whereby China redirects sales from the U.S.
toward third destinations and erodes other exporters' market shares; (ii) 
\emph{import competition}, whereby redirected Chinese goods raise import
penetration and competitive pressure in domestic markets; (iii) \emph{consumer
opportunities}, whereby cheaper imported goods from China reduce prices for 
consumers in destination markets; and (iv) \emph{sourcing opportunities}, 
whereby cheaper Chinese intermediates reduce input
costs and strengthen downstream competitiveness. 

We operationalize these channels with transparent and theoretically-grounded
exposure indices, and we also compare them to alternative exposure indices
based on export and import similarity measures %
\citep[cf.][]{FingerKreinin1979ExportSimilarity}. We then use these
diagnostics to classify countries in terms of their exposure, while allowing
for multi-label exposure. \textbf{[MORE]}

% Having done this, we next turn to a second, equally salient feature of the current global
% environment: the growing entanglement of trade policy with geopolitics. This
% entanglement is visible both in official rhetoric and in policy practice.
% Across major economies, trade and investment policy is increasingly framed
% through the language of `economic security', `de-risking,' and strategic
% resilience, with governments explicitly seeking to reduce dependencies in
% sectors viewed as critical \citep{alfaro_trade_2025} or as involving
% dual-use products \citep{alekseev2025trade}. The WTO itself emphasizes that
% the multilateral trading system is being reshaped by rising security
% concerns, as wars and geopolitical rivalry broaden the scope of what states
% treat as `security' in economic policy \citep{wto2023world}. Consistent with
% this shift, the IMF warns that policy-driven geoeconomic fragmentation
% (rather than purely economic forces) is now a central risk to global
% integration \citep{Aiyar2023Geo}.\footnote{%
% These concerns have also translated into actions that sit uneasily with a
% strictly rules-based, WTO-centered regime. The U.S.--China conflict is the
% most salient example: the United States has recently imposed large-scale
% Section 301 tariffs without first securing approval from the WTO. In
% parallel, the WTO's dispute-settlement system has remained impaired as the
% Appellate Body has been effectively inoperative since late 2019, with
% repeated failures to restart appointments. On China's side, export controls
% on strategic inputs and technologies -- such as licensing requirements for
% critical minerals used in semiconductors -- have become an overt instrument
% of statecraft.}

% Our empirical analysis in section \ref{sec:geoeconomics} underscores two
% broad patterns tied to current events. First, trade flows have become
% increasingly shaped by geopolitics: countries that are more closely aligned
% tend to trade disproportionately more with each other. Second, the policy
% environment has also evolved in ways that do not fit a simple `targeted
% decoupling' narrative: in the most recent phase, U.S. tariff increases
% appear broad-based rather than concentrated on geopolitical adversaries,
% consistent with a more unilateral and transactional use of tariffs as a
% general instrument of leverage. Finally, these changes have taken place
% against the backdrop of shifting geopolitical and `geoeconomic' power,
% associated with a very pronounced increase in Chinese power and an equally
% notable decline in the relative power of the United States.

% When governments weaponize trade policy by treating tariffs and related
% measures as instruments of strategic rivalry, the sustainability of
% cooperative trade arrangements may weaken even if the economic logic of
% specialization remains intact. The analytical framework developed in section %
% \ref{Sec:Theory} of this paper formalizes this intuition by embedding
% geopolitical pressures into a standard trade model in which governments are
% tempted to resort to trade policy to manipulate their terms of trade in
% non-cooperative settings \citep{johnson1953optimum,bagwell1999economic}. Our
% framework yields a clear implication: cooperative outcomes with low tariffs
% are more difficult to sustain when economic size and geopolitical power are
% tightly aligned, because large and influential countries can often secure
% relatively favorable outcomes even absent binding commitments.

% Empirically, we construct proxies for geoeconomic influence and show that
% they are not only positively related to economic size, but that the
% size--power alignment has strengthened markedly since the mid-1990s. Through
% this lens, the recent strain on the rules-based trading system is not a
% paradox but a predictable consequence of a world in which the countries with
% the greatest market power also possess the greatest geopolitical leverage.
% When those two dimensions coincide, large states can obtain comparatively
% favorable outcomes outside cooperative constraints, raising the temptation
% to defect and making low-tariff cooperation harder to sustain.

% Taken together, our findings point to a future in which globalization
% persists, but increasingly takes the form of \emph{fragmented interdependence%
% }: continuity in aggregate openness alongside a sharper fragmentation in who
% trades with whom and through which routes. As the U.S.--China relationship
% is rewired, third-country ``connectors'' and geopolitically aligned partners
% become pivotal nodes, so diversion and rerouting are likely to be recurring
% margins of adjustment rather than transitory responses to a one-off shock.
% More broadly, because the same rise in geoeconomic rivalry that drives
% reorganization also strains the incentives to cooperate, the next phase of
% globalization is likely to be governed less by universal rule-making and
% more by a patchwork of contested, selectively enforced arrangements.

Our paper relates to several strands of work. First, it speaks to the debate
on `deglobalization' and geoeconomic fragmentation, which emphasizes that
global integration has not uniformly reversed but is increasingly shaped by
security concerns and policy-driven shifts in the organization of trade %
\citep{Antras2021DeglobalisationECB,goldberg2023global,AiyarPresbiteroRuta2023GEFBook}%
. Second, it connects to the emerging evidence on the `great reallocation'
in U.S. sourcing and global supply chains following the U.S.--China tariff
war %
\citep{AlfaroChor2023JacksonHoleGreatReallocation,AlfaroChor2025AnatomyGreatReallocation}%
, including work highlighting reallocation, rerouting, and `exports in disguise' through
third countries \citep{MaleskyEtAl2024ExportsInDisguise,gopinath2025changing,Schulze2025}. Third, it complements quantitative assessments of the incidence and
welfare consequences of recent U.S. tariff actions and retaliation %
\citep{AmitiReddingWeinstein2019JEP,FajgelbaumEtAl2020Return,CavalloEtAl2021TariffPassThrough,FlaaenHortacsuTintelnot2020WashingMachines}%
, and uses up-to-date tariff information to situate the post-2018 policy
regime \citep{Bown2025USChinaTariffsChart}. \textbf{MORE - RECENT LITERATURE ON CHINA 2.0}
% Finally, our mechanism builds on
% the classic terms-of-trade logic of unilateral tariffs and the cooperative
% role of trade agreements \citep{johnson1953optimum,bagwell1999economic},
% while embedding the idea that trade policy is increasingly intertwined with
% geopolitical leverage and economic coercion %
% \citep{hirschman1945national,Baldwin1985EconomicStatecraft,FarrellNewman2019WeaponizedInterdependence,ClaytonMaggioriSchreger2024NBERCoercionFragmentation,broner2025hegemonic,alekseev2025trade,BeckoOConnor2025StrategicDisIntegration,BeckoGrossmanHelpman2025OptimalTariffsGeopoliticalAlignment}%
% . To discipline these geopolitical forces empirically, we draw on measures
% of national power and trade leverage %
% \citep{moyer2024relative,liu2025international} and on approaches to
% measuring alignment and influence in international relations %
% \citep{Bailey-etal_2017,kuziemko2006how,antras2011foreign,antras2025exporting,bubeckmarinov2017,bubeckmarinov2019}%
.

%%% ADDING HERE SOME NOTES, too long for now and some parts can be moved elsewhenre (sec 2) %%%
A growing number of papers and analyses---especially in policy circles---point to an increasing impact of the expansion of Chinese exports to other countries, and especially Europe.\footnote{\ The discussion of the effects of the China shock 2.0 is now at the forefront of the policy debate, as testified b its extensive news coverage---see Financial Times [add REFS], among others.} While the increase of Chinese export globally predates the 2025 tariffs and is the continuation of a secular trend---partly due to weak exchange rate, industrial policy, and subdued domestic demand, consistent with the ``vent for surplus'' mechanism \citep{almunia2021}---the U.S. April 2025 tariffs are a large enough shock to trigger a significant further reallocation of trade flows. 
\citet{Emlinger2026} study the initial response to the April 2025 U.S. tariffs and do not find robust evidence of a deflection of Chinese exports to Europe until October 2025. Extending the sample to December 2025, \citet{Schulte2026} find that export diversion to Europe has been concentrated in a small tail of highly exposed products, suggesting limited macroeconomic effects. By contrast, Chinese exports seem to have been mostly redirected to emerging and developing countries. 

Different from the first China shock, this time Chinese export growth is increasingly concentrated in capital- and technology-intensive sectors \citep{deSoyres2026}. This shift translates into an increasing export similarity between China and advanced economies, while this is not the case for several emerging markets which are less aligned in export structures and may face less competitive pressures. A recent analysis by \citet{kiel2026} support this hypothesis, as increasing Chinese exports displace those of large advanced economies in third markets, especially in more complex products and technologically-intensive industries. However, \citet{chatterjee2026} argue that developing economies are not insulated from the China shock 2.0. On the opposite, rising Chinese competition---partly still concentrated on traditional, low-skilled goods---does not only imply foregone exports in global markets but also a squeeze for domestic producers. As a result, the China shock 2.0 risks slowing down the industrialization of developing countries.

Consistent with our proposed framework, \citet{aprigliano2026} look at the effects of the China shock 2.0 in the euro area and consider three main margins of adjustment. On the negative side, increasing Chinese exports erode euro area market shares abroad. In addition, through the import competition channels, domestic manufacturing investment declines. However, on the positive side the shock leads to a decline of non-energy goods prices by 1 percent over three years.



The remainder of the paper is organized as follows. Section \ref{Sec: Great
Reallocation} documents the key empirical patterns of the Great Reallocation
associated with U.S.--China decoupling. Section \ref{Sec: Third Countries}
proposes a taxonomy to characterize third-country spillovers.  
Section \ref{Sec: Conclusion} concludes.

\section{The U.S.--China Trade War and the Great Reallocation\label{Sec:
Great Reallocation}}

The U.S.--China trade war has been one of the most significant economic
events in recent times. Early U.S. trade actions included Section 201
safeguard tariffs on solar products and washing machines and Section 232
tariffs on steel and aluminum, but the dispute quickly broadened into
large-scale, China-specific Section 301 tariffs. Over 2018--19, successive
Section 301 rounds expanded from targeted product lists to cover most
bilateral trade, pushing the trade-weighted average U.S. tariff on imports
from China from a low pre-war baseline into double digits 
\citep[see
Figure~\ref{fig:us_china_effective_tariff}, based
on][]{Bown2025USChinaTariffsChart}. China followed suit and retaliated in a
commensurate manner. The January 2020 ``Phase One'' agreement halted further
escalation and reduced some tariff rates, but most duties remained in place
and bilateral trade barriers stayed far above pre-2018 levels, with a
trade-weighted average of roughly 20\% in both countries.

\begin{figure}[htbp]
\caption{The U.S.--China Trade War: Trade-Weighted Average Tariffs 2016-2025}
\label{fig:us_china_effective_tariff}
\begin{center}
{\includegraphics[width=.7\textwidth]{Figures/tariffs.png}}
\end{center}
\par
\begin{minipage}{.8\textwidth}
{\footnotesize  \textbf{Notes:} The series plot trade-weighted average applied tariffs. The solid blue line is the U.S. tariff on imports from China; the dashed red line is China's tariff on imports from the United States. The PIIE series begins in 2018; 2016-2017 are shown as a flat extension at the January 2018 level for display. The series excludes antidumping/countervailing duties. 
\par
\textbf{Source:} \cite{Bown2025USChinaTariffsChart}.}
\end{minipage}
\end{figure}


\begin{figure}[htbp]
\caption{The Great Reallocation: U.S. import shares (2017-2025)}
\label{Fig:Mkt_Shares}%
\begin{subfigure}[t]{0.47\linewidth}
    \centering
    \includegraphics[width=\linewidth]{Figures/US_import_shares_top3.png}
    \caption{U.S. import shares for large exporters}
 
  \end{subfigure}\hfill 
\begin{subfigure}[t]{0.49\linewidth}
    \centering
    \includegraphics[width=\linewidth]{Figures/US_import_shares_winners.png}
    \caption{U.S. import shares for some smaller exporters}

  \end{subfigure}
\par
\bigskip 
\begin{minipage}{.98\textwidth}
\footnotesize 
\textbf{Notes:} Shares are calculated as each partner's goods imports divided by total for all economies. Partner shares computed for China, Mexico, Canada, Japan, Vietnam, Taiwan Province of China, South Korea, India, and Thailand. Data for 2026 refer to the period January-April. \par
  \smallskip
\textbf{Sources:} U.S. Census Bureau (\url{https://www.census.gov/foreign-trade/balance/index.html}).
\end{minipage}
\end{figure}

From 2021 through 2024, this tariff structure was largely maintained with
only selective and sector-specific adjustments. In 2025, however, the
conflict entered a much more volatile phase. A new wave of tariff actions
associated with President Trump's April 2, 2025 ``Liberation Day'' tariffs
and related executive orders pushed effective tariff rates sharply higher
(including brief peaks at triple-digit levels) before partial rollbacks, as
seen in Figure\ \ref{fig:us_china_effective_tariff} %
\citep{Bown2025USChinaTariffsChart}. On May 12, 2025, a joint U.S.--China
statement issued after talks in Geneva committed the United States to remove
many of the duties imposed in early April, but as of November of 2025, the
trade-weighted average U.S. tariff on imports from China remained at 47.5\% 
\citep[see Figure \ref{fig:us_china_effective_tariff},
and][]{Bown2025USChinaTariffsChart}.\footnote{%
The ratio of \emph{collected} customs duties to U.S. imports from China has
settled at a slightly lower level of around 40\%, as reported by Mike
Waugh's ``Trade War Tracker'' https://www.tradewartracker.com/.} The scale
and framing of the 2025 escalation were widely read as evidence that the
U.S.--China relationship is being redefined on geopolitical grounds, making
the shock appear long-lasting (or even permanent) rather than transitory.
With that interpretation in mind, we now document how international trade
flows have adjusted.

\subsubsection*{The Great Reallocation}

Against that backdrop, the central adjustment has been a reorganization of
trade relationships -- a `Great Reallocation'\ of U.S. sourcing patterns --
with broader implications for third countries. Figure \ref{Fig:Mkt_Shares}
summarizes the evolution of U.S. import market shares across major partners
from 2015 April of 2026 (the most recent month for which
data is available). The dominant fact is an initially gradual but recently
more pronounced decline in China's share. Overall, the Chinese share in U.S. imports fell from around 22 percent in 2017 to around 8 percent in 2026. This decline was mirrored by gains for a set of
alternative suppliers, most prominently Vietnam, Mexico and Taiwan.
Importantly, the timing of these gains is uneven: Vietnam's rise is
relatively gradual, whereas Taiwan, Thailand and Mexico's gains are much more
concentrated in the later years of the reallocation (2024--2026).

Much has been written about the slow initial decline in U.S. imports from
China. The COVID-19 pandemic was certainly a contributor: although imports
fell at the onset of the pandemic, they surged back within months as
consumers shifted toward goods and e-commerce, helping propel a strong
rebound in U.S. goods imports from China. Another key reason why the early
tariff shock did not immediately translate into a one-for-one collapse of
Chinese exports to the United States is that supply chains feature
meaningful adjustment frictions. \cite{Antras2021DeglobalisationECB}
emphasized the importance of sunk costs in generating hysteresis in
importing given the large policy uncertainty in the early 2020s, while \cite%
{AlfaroChor2025AnatomyGreatReallocation} have recently shown that the
reallocation away from China was slower in contract-intensive and
relationship-sticky goods. In Appendix \ref{App: APP}, we provide evidence
for a novel channel contributing to the gradual nature of the U.S.--China
decoupling. More specifically, building on a measure of the duration of
production developed by \cite{antras2025measuring} -- the so-called average
period of production or $\mathcal{APP}$ -- we show that reallocation of Chinese 
exports away from the U.S. was faster in products with a shorter $\mathcal{APP}$,
while products with longer production cycles adjusted more slowly.



\subsubsection*{Trump 2.0 and Liberation Day}

Returning to the more drastic reallocation observed since the start of
President Trump's second administration, Figure \ref{fig:Liberation_Day_US}
uses monthly data to zoom in on the impact on U.S. imports for selected
trading partners during 2025. The figure illustrates the remarkable drop in
U.S. imports from China coupled with a reallocation of those import flows to
Taiwan, Vietnam and to a lesser extent Mexico. Imports from Canada fell
during this period. Interestingly, with the exception of China, there
appears to be a \emph{positive} relationship between import growth and the
tariffs imposed on these partners on Liberation Day, since Mexico and Canada
were imposed a 10 percent tariff, while Taiwan Province of China and Vietnam were given tariff
rates of 32 and 46 percent, respectively. Beyond these five trading
partners, there does not seem to exist a strong systematic relationship
between changes in U.S. imports and the tariff rates faced by partner
countries. Figure \ref{fig:Liberation_Day_Scatter}, inspired by Figure 9 in 
\cite{AlfaroChor2025AnatomyGreatReallocation}, may be interpreted as
suggesting that countries facing higher new tariff rates tended, on average,
to see weaker growth in U.S. imports. However, this pattern is not robust:
it is largely driven by China, and it fades when the relationship is
estimated via ordinary least squares (rather than when weighting
observations by import value, as in %
\citeauthor{AlfaroChor2025AnatomyGreatReallocation}, %
\citeyear{AlfaroChor2025AnatomyGreatReallocation}).\footnote{%
The legend in Figure \ref{fig:Liberation_Day_Scatter} includes the
regression results. The weighted least squares coefficient excluding China
is negative but much smaller and significant only at the 10 percent level.
The ordinary least squares coefficient (including China) is essentially 0.}

\begin{figure}[htbp]
\caption{The Great Reallocation during Trump's Second Administration}
\label{fig:Liberation_Day_US}
\begin{center}
{\includegraphics[width=.65\textwidth]{Figures/US_market_shares_wide_post_trump.png}%
}
\end{center}
\par
\begin{minipage}{.9\linewidth}
{\footnotesize \textbf{Notes:} Series show U.S. goods imports from each partner, indexed to Nov. 2024 = 100. Values are a 3-month backward moving average (month t equals the average of t-2, t-1, and t). Vertical dotted line marks Nov. 2024.
\par
\smallskip
\textbf{Source:} U.S. Census Bureau (\url{https://www.census.gov/foreign-trade/balance/index.html}).}
\end{minipage}
\end{figure}

\begin{figure}[htbp]
\caption{Liberation Day and Sources of U.S. Imports}
\label{fig:Liberation_Day_Scatter}
\begin{center}
{\includegraphics[width=.65\textwidth]{Figures/Fig_AC_Updated.png}}
\end{center}
\par
\begin{minipage}{.95\linewidth}
{\footnotesize \textbf{Notes:} Each dot is a U.S. import partner. The y-axis is the change in the partner's share of total U.S. imports from March to October 2025, demeaned by the partner's average March-to-October change over 2022-2024. Dot size is proportional to the partner's 2024 total imports. Lines are linear fits: import-weighted (weights = 2024 imports), import-weighted excluding China, and unweighted OLS; significance stars denote p<0.10, *p<0.05, **p<0.01.
\par

\textbf{Source:} U.S. Census Bureau import totals by partner; Liberation Day country tariff rates from the April 2, 2025 Executive Order, Annex I.}

\end{minipage}
\end{figure}

Liberation Day also led to a remarkable reorientation of Chinese exports.
Figure\ \ref{fig:trump2_reallocation} tracks China's exports to the United
States relative to its exports to other destinations, normalized to November
2024. After the November 2024 election, Chinese exports to the U.S. declined
sharply, while exports to Europe and to other Asian economies rose,
suggesting a rapid diversion of Chinese sales away from the U.S. market and
toward alternative destinations. This prospect has heightened anxiety in the
recipient markets: in Europe, policymakers explicitly warned about trade
diversion and moved to strengthen import-surveillance capacity to detect and
respond to sudden surges, amid concerns about being \textquotedblleft
flooded\textquotedblright\ by redirected Chinese goods. In parts of Asia,
similar fears about displacement and import surges have been reflected in a
more defensive stance -- e.g., new or extended anti-dumping actions against
specific Chinese products -- alongside broader public concern about the
competitive pressure created by low-priced Chinese exports being redirected
into regional markets. [\textbf{[ADD REFERENCES; CHINA 2.0]}]

\begin{figure}[htbp]
\caption{The Reallocation of Chinese Exports during Trump's Second
Administration}
\label{fig:trump2_reallocation}
\begin{center}
{\includegraphics[width=.65\textwidth]{Figures/China_exports_us_eu_asia.png}}
\end{center}
\par
\begin{minipage}{.9\linewidth}
{\footnotesize \textbf{Notes:} Exports are 3-month backward-looking moving
averages. The vertical line is drawn in correspondence with the November 2024
U.S. presidential election. The Asia aggregate includes Bangladesh, Cambodia, Hong Kong SAR, India, Indonesia, Malaysia, Pakistan, the Philippines, Singapore, Sri Lanka, Taiwan Province of China, Thailand, and Vietnam. EU = European Union.
\par

\textbf{Sources:} Trade Data Monitor.}
\end{minipage}
\end{figure}

The reorientation of U.S. imports and Chinese exports is mirrored in broader
shifts in global trade patterns. Figure \ref{Fig: Heat Map} plots the change
in bilateral import shares (in panel (a)) and export market shares (in panel
(b)) for various countries and regions when comparing the year after the
U.S. election (Nov. 2024 -- Feb. 2026) to the previous period Jun. 2023 -- Oct. 2024. At one level, the heat map reiterates the same basic narrative in Figures \ref%
{fig:Liberation_Day_US} and \ref{fig:trump2_reallocation}: the share of
China's exports flowing to the U.S. market contracts sharply (-3.3 p.p.),
but it increases in Asia (+2.4 p.p.) and the rest of the world (+0.9 p.p.).
On the flip side, the share of U.S. imports originating in China falls by
3.5 p.p., with large increases from Asia (+1.5 p.p.) and the rest of the world. But
Figure \ref{Fig: Heat Map} also makes clear that the post-election
reallocation is not confined to the U.S.--China dyad: there are sizeable market-share shifts in third markets that do not directly involve either country. For example, the America bloc (Canada--Mexico) reoriented their exports towards China, the EU market, and rest-of-the-world destinations (+0.2 p.p.). At the same time, countries in the European Union recorded gains in import shares across China (+0.6 p.p.), developing Asia (+0.6 p.p.), Canada and Mexico (+0.2 p.p.), while reducing imports from the rest of the world region (-1.6 p.p.).

\begin{figure}[htbp]
\caption{A Worldwide Reallocation during Trump's Second Administration}
\label{Fig: Heat Map}%
\begin{subfigure}[t]{0.45\linewidth}
    \centering
    \includegraphics[width=\linewidth]{Figures/heatmap_X_ms.png}
    \caption{Change in export shares}
   \hspace{0.25\textwidth}
  \end{subfigure}\hfill 
\begin{subfigure}[t]{0.45\linewidth}
    \centering
    \includegraphics[width=\linewidth]{Figures/heatmap_M_ms.png}
    \caption{Change in import shares}
   
  \end{subfigure}
\par
\bigskip
\par
\begin{minipage}{.90\linewidth}
{\footnotesize \textbf{Notes:} The two charts report the change in export market shares (panel a) and import market shares (panel b) in percentage points, computed around the 2024 U.S. election, comparing Nov. 2024--Feb. 2026 to Jun. 2023--Oct. 2024. The Asia aggregate includes Bangladesh, Cambodia, Hong Kong SAR, India, Indonesia, Malaysia, Pakistan, the Philippines, Singapore, Sri Lanka, Taiwan Province of China, Thailand, and Vietnam. EU = European Union. ROW = Rest of the world. Changes in market shares sum to zero by row in panel (a) and by column in panel (b).
\par
\textbf{Sources:} Trade Data Monitor.}
\end{minipage}
\end{figure}

\subsubsection*{Taking Stock}

Taken together, Figures\ \ref{fig:Liberation_Day_US}--\ref{Fig: Heat
Map} suggest that Liberation Day\ may have been interpreted less as a
temporary escalation and more as a permanent shock to the U.S.--China trade
relationship -- an inflection point toward deeper, sustained decoupling.
While other countries were also hit with higher U.S. tariffs, those measures
were surrounded by reversals and uncertainty, so it is unsurprising that we
see no clean, systematic trade response. By contrast, the sharp curtailment
of China's ability to sell into the U.S. market generated sizable spillovers
elsewhere: Chinese exports were redirected into other destinations, and
competitive pressures intensified in third markets where China and other
exporters meet head-to-head. We next turn to a more systematic analysis of
these spillover effects to third markets.

\section{Caught in the Middle:\ Third-Country Spillovers from the
U.S.--China Tariff War \label{Sec: Third Countries}}

When two giants fight, smaller economies do not simply \textquotedblleft sit
out\textquotedblright\ the conflict. The new phase of the U.S.--China tariff
war has curtailed China's sales into the U.S.\ market and, in turn, altered
the competitive environment faced by firms around the world. This section
examines these third-country spillovers systematically.

A natural starting point is trade diversion in the narrow sense: third
countries may gain U.S.\ market share as Chinese suppliers are displaced, or
they may gain market share in China if U.S. sellers are displaced. The
patterns since Liberation Day, however, offer limited support for this
narrow diversion narrative. The cross-partner relationship between U.S.\
import growth and the new tariff schedule appears weak (Figure\ \ref%
{fig:Liberation_Day_Scatter}), and the relatively small scale of U.S.\ goods
exports to China (compared with U.S.\ imports from China) implies that, for
many countries, direct bilateral exposure to China through that channel is
second-order. These observations motivate a different perspective in which
the main action occurs in markets \emph{other} than the U.S.\ itself.

Our emphasis is therefore on spillovers generated by China's impaired access
to the U.S.\ market but realized in third markets. We decompose exposure
into four conceptually distinct margins. First, \emph{export competition}:
as Chinese exporters redirect shipments, they displace other suppliers in
shared destination markets. Second, \emph{import competition}: redirected
Chinese goods raise import penetration and competitive pressure in domestic
markets. Third, \emph{consumer} opportunities: those same redirected
Chinese goods will tend to be sold at lower prices than the purchases they
displace, which will reduce the price level for consumers in those import
markets, and perhaps also increase the variety of goods available to them.
Fourth, \emph{sourcing opportunities}: lower-priced (or more readily
available) Chinese intermediates can reduce input costs and expand
downstream activity. A special case of this latter channel reflects a 
\emph{rerouting opportunity}, and captures how cheaper Chinese imports enhance
exports specifically to the U.S. market, which can be interpreted as Chinese
firms circumventing U.S.\ barriers by shifting assembly, routing trade
through intermediaries, or relocating stages of production to countries
with better market access to the United States.

We next operationalize this taxonomy with tractable diagnostics and then use
them to characterize which countries are most exposed along each margin.
Before we dive into the specifics of our indices, we offer some broad
observations applying to all these indices.

\subsection{General Considerations}

Our goal is to develop measures of the extent to which countries are exposed
to an increase in Chinese exports to countries other than the U.S.,
resulting from the constrained ability of these same Chinese exporters to
sell to the United States. At a very general level, it is clear that
countries are more likely to face export competition from China if their
export patterns are \emph{similar} to those of China. At the same time, the
other channels of exposure should be more salient when a country's reliance
on imports from China is \emph{large} and \emph{similar} in structure to
that of the U.S., since this would indicate that U.S. policies would divert
relatively more Chinese exports toward that exposed country.

A\ common approach to capture similarity in exports and imports follows from
the influential work of \cite{FingerKreinin1979ExportSimilarity}. Let us
review the basics of their index for the case of exports. Suppose that one
is interested in assessing how similar is the export structure of two
countries, $i$ and $j$. By export structure, we mean a vector $\mathbf{X}%
_{o}=\left\{ X_{o11},...,X_{odp}\right\} $ where $X_{odp}$ refers to
(origin) country $o$'s exports of product $p$ in (destination) market $d$. 
\cite{FingerKreinin1979ExportSimilarity} proposed the following similarity
index between the exports of country $i$ and $j$ in destination market $d$:%
\begin{equation}
ESI_{ijd}=\sum_{p\in \mathcal{P}}\min \left( \frac{X_{idp}}{%
\sum\nolimits_{p^{\prime }\in \mathcal{P}}X_{idp^{\prime }}},\frac{X_{jdp}}{%
\sum\nolimits_{p^{\prime }\in \mathcal{P}}X_{jdp}}\right) \times 100.
\label{eq: FingerKreinin}
\end{equation}

To see how this measures captures export similarity, notice that if the
countries $i$ and $j$ feature the same \emph{share} of exports in all goods $%
p$, then $ESI_{ijd}=100$ because the two countries are 100 percent similar
in their export structure in the destination market $d$. Conversely, if
there is \emph{no} product $p$ for which \emph{both} countries feature
positive exports, then their similarity is zero.\footnote{%
As shown by \cite{sun2000measurement}, the \cite%
{FingerKreinin1979ExportSimilarity} index has good axiomatic properties, in
the sense that it is invariant with respect to proportional
sub-classifications.}

The \cite{FingerKreinin1979ExportSimilarity} index is an insightful index of
export similarity, but it is much less clear that it constitutes a good
index of export competition pressures. As pointed out by \cite%
{jenkins2008measuring}, the fact that the index $ESI_{ijd}$ in equation (\ref%
{eq: FingerKreinin}) is symmetric is problematic: when interpreted as
capturing competitive pressures, it would imply that China is as impacted by
competition from Honduras as Honduras is impacted by competition from
China. Furthermore, the index is also symmetric across products and thus does not
take into account that competitive pressures might be felt more intensively
in some products than others.

With that in mind, below we will develop alternative indices of exposure
that are theoretically grounded, and we will compare them with the
corresponding similarity measures. More specifically, in Appendix \ref{App: MicroFound},
we develop an Armington model with product-level CES demand: in each
destination market, expenditure on product $p$ is allocated across exporter
varieties according to a constant elasticity. A China-specific U.S. policy
shock reduces Chinese exports to the U.S. and, via an upward sloping Chinese
export supply function, leads to lower prices in all other destinations. The
model then traces the first-order implications of this shock on the exports of a
given country $j$ to all destination markets other than the U.S. and China.

We construct all our exposure measures
based on exports and imports data. We use bilateral trade data at the product
level sourced from Trade Data Monitor (TDM), which compiles monthly trade
data from government offices and national agencies. Data are available with
a lag of just a few months and cover monthly goods trade (imports and
exports) from 104 reporting economies to 195 partner economies, and are
disaggregated at the six-digit Harmonized System (HS-6) level, covering
7,352 unique products. We construct a balanced panel of trade flows (in
nominal USD) between 195 countries using imports from country A to country B
to measure exports from B to A, when the latter is not available because B
is not a reporting country. 

\subsection{Diagnosing Export Exposure}

With this background in mind, we begin by proposing indices capturing the
extent to which third countries may face increased export competition from
China in countries other than the US (and China) as a result of the
inability of Chinese exporters to sell to the United States. We begin
with the standard approach of producing measures of export similarity,
but then propose our theoretically-grounded index of export competition.

\subsubsection{Export Similarity}

As argued before, export competition from China might be expected 
to be fiercer for countries with an export structure relatively more similar 
to that of China. To measure the export similarity between a given country $j$ and
China, we first compute the following variant of the \cite%
{FingerKreinin1979ExportSimilarity} export similarity index ($ESI$):%
\begin{equation}
ESI_{j}= 100 \times \sum_{d\in \mathcal{W}_{-}}\sum_{p\in \mathcal{P}}\min \left( \frac{%
X_{jdp}}{X_{j}},\frac{X_{cdp}}{X_{c}}\right) ,
\label{eq: ESI Formula}
\end{equation}%
where $c$ denotes China, $X_{jdp}$ are country $j$'s exports of product $p$
in destination $d$ (in some baseline year), $\mathcal{P}$ is the set of all
products, $\mathcal{W}_{-}$ is the set of all countries in the world ($%
\mathcal{W}$) excluding the U.S.\ and China, and where $X_{j}$ and $X_{c}$
are aggregate exports of $j$ and China in all countries other than U.S. and
China (in the baseline year). To avoid anticipation effects in the run-up of
the U.S. presidential election, we choose 2023 as our baseline year.\footnote{
We have also computed a variant of the export similarity index that 
aggregates across destinations as follows:
\begin{equation}
ESI_{j}= 100 \times \sum_{d\in \mathcal{W}_{-}}\frac{%
X_{jd}}{X_{j}}\sum_{p\in \mathcal{P}}\min \left( \frac{%
X_{jdp}}{X_{jd}},\frac{X_{cdp}}{X_{cd}}\right) ,
\end{equation}
which amounts to an export-weighted average of the destination-specific
export similarity indices with China. The correlation.}

In robustness tests, we also compute coarser indices of export similarity
that only consider similarity in the types of products exported (regardless
of where they are sold) and in the destination of sales (regardless of which
goods are sold in those destinations). The product-level index, when
computed in some baseline year is given by%
\begin{equation}
ESI_{j}^{prod}= 100 \times \sum_{p\in \mathcal{P}}\min \left( \frac{\sum_{d\in \mathcal{%
W}_{-}}X_{jdp}}{X_{j}},\frac{\sum_{d\in \mathcal{W}_{-}}X_{cdp}}{X_{c}}%
\right) ,  \label{eq: ESI Formula Product}
\end{equation}%
and it may capture better than equation (\ref{eq: ESI Formula}) situations
in which the shock (e.g., Liberation Day) leads to a disproportionate
increase in Chinese exports to destinations in which it sold relatively less
in the baseline year. On the other hand, the coarser similarity index based
on the destination of sales is given by 
\begin{equation}
ESI_{j}^{dest}= 100 \times \sum_{d\in \mathcal{W}_{-}}\min \left( \frac{\sum_{p\in 
\mathcal{P}}X_{jdp}}{X_{j}},\frac{\sum_{p\in \mathcal{P}}X_{cdp}}{X_{c}}%
\right),  \label{eq: ESI Formula Destination}
\end{equation}%
and it may better capture concerns about China disproportionately increasing
the sales of \emph{new} types of goods in markets in which both $j$ and China had
sizeable market shares in the baseline year.

Table \ref{tab:top20_esi_dp} lists the 20 self-governed economies that, according to
export similarity at the product and destination levels (see equation (\ref%
{eq: ESI Formula})), face the highest export exposure to China.\footnote{We
exclude Hong Kong SAR, which has the highest export similarity with China at 31.01. In
later tables we also exclude Macao SAR and other non-self-governed territories.
Values for all countries and territories are provided on our support website. \textbf{[ADD 
LINK IN DUE COURSE]}} Most of
these most exposed economies are in Asia, although the table also includes a
number of European economies, such as Germany or Italy.

\begin{table}[htbp]
\caption{Top 20 economies by Export Similarity ($ESI_{j}$)}
\label{tab:top20_esi_dp}\centering
\include{Figures/esi_top20.tex}
\end{table}

Figure \ref{fig:esi} plots our benchmark export similarity measure against
the two previously mentioned coarser measures that allow for more flexible channels of export
competition. Panel (a) plots export similarity when computed at the product
level. The correlation between the two indices is very high (0.87) but, at
the product level, export similarity with China is always larger. The
figure further indicates that European countries tend to be much more
impacted by this measure, reflecting the fact that these countries export
the same products as China does, but not necessarily to the same
destinations. This further implies that if Chinese exporters were to
reorient their sales to these destinations currently favored by European
firms, the shock could be much more intensively felt in European countries.

Panel (b) is analogous but plots export similarity when computed at the
destination level. This captures threats stemming from China ramping up
their exports of new or low-selling products in the destinations in which
the potentially impacted countries concentrate their sales in the baseline
year. European countries seem to be more immune from this shock, but many Asian
economies -- as well as Canada and Brazil -- would face much fiercer enhanced 
export competition on this account.

\begin{figure}[htbp]
\centering
\par
\begin{subfigure}[t]{0.48\linewidth}
    \centering
    \includegraphics[width=\linewidth]{Figures/esi_product_vs_product_destination_by_continent_updated.png}
    \caption{Export Similarity at the Product Level}
  \end{subfigure}\hfill 
\begin{subfigure}[t]{0.48\linewidth}
    \centering
    \includegraphics[width=\linewidth]{Figures/esi_destination_vs_product_destination_by_continent_updated.png}
    \caption{Export Similarity at the Destination Level}
  \end{subfigure}
\caption{Alternative Measures of Export Similarity}
\label{fig:esi}
\end{figure}

\subsubsection{An Export Competition Index}

We next turn to our theoretically-grounded index of export competition. We
omit all the technical details in the main text, but the expressions below
are derived from an Armington model in Appendix \ref{App: MicroFound}. The 
percentage response of country $j$'s exports is shaped by three key
empirical objects. 

\paragraph{a) Size of product-level shock}
A first consideration is the size of the U.S.-China tariff shock for 
different products $p$. More specifically, we seek to measure by how much will 
(free-on-board) Chinese prices fall in third markets as a result of the tariffs
set by the U.S. on China, and we denote that Chinese price decline by $\Phi _{p}$. For 
uniform U.S. tariff changes and uniform demand and supply elasticities
across sectors, the relative size of the shock across sectors turns out to be 
governed solely by the share of overall Chinese production of good $p$ that was 
shipped to the U.S. (relative to other destinations) in the baseline year, or%
\begin{equation}
\phi _{p}=\frac{X_{cup}}{\sum\nolimits_{d\in \mathcal{W-}\left\{ c\right\}
}X_{cdp}}.  \label{eq: Relative Shock}
\end{equation}%
Intuitively, the larger is the share of Chinese output that was sold in the
U.S., the larger will be the amount of exports that will need to be redirected
to other markets, and with an upward sloping Chinese export supply, the larger 
will be the necessary price decline in other markets.

\begin{table}[htbp]
\centering
\small
\hspace*{-0.5cm}
\include{Figures/phi_cup_top20.tex}
\caption{Top 20 HS6 Products by $\phi_{p}$}
\label{tab:top20_phi}

\vspace{0.5em}

\begin{minipage}{.98\textwidth}
\footnotesize 
\textit{Note:} HS codes with missing elasticity estimates or tariff data are excluded. 
\end{minipage}
\end{table}

Table~\ref{tab:top20_phi} reports the twenty HS6 products with the highest values of $\phi _{p}$ in 2023, excluding special classification codes.\footnote{Table \ref{tab: price_shocks} reports the twenty HS6 products with the highest values of $\Phi_{p}$, and decomposes it into the components of (\ref{Phi_p}).} The products are concentrated in a relatively narrow set of primary commodities (e.g., seafood and cobalt ores), specialized chemical inputs, and selected high-value metals such as rhodium and beryllium. High values of $\phi_{p}$ indicate that, for Chinese exporters of these goods, the U.S. market represents a particularly important destination. Exports of these products from China are therefore disproportionately oriented toward the United States.\footnote{Importantly, these results should \emph{not} be interpreted as evidence that China is a critical supplier of these goods for the United States. A high $\phi _{p}$ reflects the dependence of Chinese exporters on the U.S. market, not the dependence of the U.S. on China as a source country.}

As the model in Appendix \ref{App: MicroFound} shows, the size of the 
product-level tariff shocks is not only shaped by shares of Chinese exports 
going to the U.S., but also by the tariff changes and the demand and supply 
elasticities in those products. More specifically, the implied decline
in the Chinese (f.o.b) price of good $p$ in other destinations is given by

\begin{equation}
\Phi_{p}=-d\log p_{cp}=\phi_{cup}\times\frac{\varepsilon_{cup}}{\bar{\varepsilon}_{cp}+\eta_{cp}}\times\Delta\log\left(1+\tau_{cup}\right).
\label{Phi_p}
\end{equation} 

In this expression, $\varepsilon_{cup}$ is the U.S. demand elasticity
for Chinese product $p$, $\bar{\varepsilon}_{cp}$ is the average 
import-demand elasticity faced by China across export markets, $\eta_{cp}$
is the elasticity of supply of Chinese exports of good $p$, and
$\Delta\log\left(1+\tau_{cup}\right)$ is the change in US tariffs
on imports from China of good $p$. 

We construct a proxy for $\Phi_{p}$ based on $\phi_{p}$, demand and supply
elasticities from \cite{soderbery2018trade}, and tariff changes from
the Global Tariff Database in \cite{rodriguez2025}
\citep[based on the methodology developed in][]{teti2024missing}. The details
are in Appendix \ref{app:Index_Empirics}. The 
correlation between $\Phi_{p}$ and $\phi_{p}$ is very high (0.822),
suggesting that most of the variation in the shock is coming from the 
transparent and observable Chinese share of exports destined to the U.S.,
rather than the more debatable elasticity estimates. Indeed, a variance
decomposition of equation (\ref{Phi_p}) indicates that more than
90\% of the variation in $\Phi_{p}$ is driven by $\phi_{p}$,
with elasticities accounting for around 7\% and tariff changes for the
remaining 3\%.\footnote{We compute tariff changes from January 2025 
(before Trump's second inauguration) to  September 2025. Later, we will 
consider robustness to a `long' variant with  export shares computed from 
2018 data and tariff differences from January 2018 to September 2025. This
has no bearing for the decomposition above: $\phi_{p}$ continues to account
for more than 90\% of the variation in $\Phi_{p}$, while the contribution
of tariff changes actually drops to 2\%.}

\paragraph{b) Incidence of product-level shock in destination market $d$}

The second key variable shaping exposure is the extent to which Chinese
exports of good $p$ diverted to a given destination $d$ will crowd out other
sellers of good $p$ in that market. Our Armington-CES model indicates that
this force can be approximated (up to a first order) by the share of
expenditure in country $d$ that is accounted for by Chinese exports, or 
\begin{equation}
\chi _{cdp}=\frac{X_{cdp}}{\sum\nolimits_{o\in \mathcal{W}}X_{odp}}.
\label{eq: China Absorption Share}
\end{equation}
The term $\chi_{cdp}$ captures the incidence of the shock in destination market $d$. Intuitively, a fall in the price of Chinese exports of product $p$ generates larger competitive pressure in markets where China is already an important supplier of that product. If China accounts for only a negligible share of purchases of product $p$ in destination $d$, then even a sizable fall in the Chinese price has little first-order effect on the revenues of other exporters in that market. By contrast, when $\chi_{cdp}$ is large, Chinese exporters are already close competitors in that product--destination cell, so a reduction in their relative price produces a larger displacement of other suppliers.

\paragraph{c) Relevance of the shock for origin country $o$}

Finally, the effect of this destination--product shock on the origin country $o$ depends on how important that same product--destination cell is for country $o$'s own exports. For a given destination $d$, define
\begin{equation}
s_{odp}=\frac{X_{odp}}{X_{od}}, \label{eq: o Export Share}
\end{equation}
where $X_{od}=\sum_{p\in\mathcal{P}}X_{odp}$. Thus, $s_{odp}$ is the share of country $o$'s exports to destination $d$ accounted for by product $p$. This term captures the relevance of the shock for the origin country: even if Chinese competition in product $p$ is intense in destination $d$, country $o$ is exposed only to the extent that it also sells product $p$ in that destination. Conversely, if product $p$ represents a large share of $o$'s exports to $d$, then displacement in that product--destination cell has a larger effect on $o$'s exports to that market.

\paragraph{The Export Competition Index}
Finally, aggregating across destinations -- weighting each destination 
by its importance in country $j$'s overall export basket -- the exposure 
of country $j$'s overall exports to the U.S.--China trade shock is 
given by
\begin{equation}
 \boldsymbol{ECI}_{j}=100\times\sum_{d\in\mathcal{W}_{-}}\frac{X_{jd}}{X_{j}}\sum_{p\in\mathcal{P}}\left(\sigma_{jp}-1\right)\times s_{jdp}\times\chi_{cpd}\times \Phi_{p}
\label{eq: ECI Baseline}
\end{equation}
Here $\mathcal{W}_{-}$ is the set of all countries in the world excluding the U.S. and China. The index therefore combines three forces: the size of the China price shock in product $p$, $\Phi_{p}$; the intensity of Chinese competition in destination market $d$, $\chi_{cdp}$; and the importance of that product for origin country $j$'s exports to destination $d$ in country $j$'s total exports.


The term $\sigma_{dp}-1$ in (\ref{eq: ECI Baseline}) appears because the relevant mechanism is substitution across country-of-origin varieties within a destination--product market. A U.S. tariff reduces Chinese export demand in the U.S.; with upward-sloping Chinese export supply, this lowers the Chinese f.o.b. price and therefore the delivered price of Chinese varieties in third markets. The revenue loss for other exporters then depends on how strongly buyers in destination $d$ substitute toward the now-cheaper Chinese variety. When the elasticity of substitution is high, a given fall in China's relative price induces a larger reallocation of expenditure away from other origins and toward China. When demand is less elastic, the same price change generates less displacement. This is why the model weights the China price shock by the demand elasticity term.

There are three main differences between the export competition index in
equation (\ref{eq: ECI Baseline}) and the Finger-Kreinin similarity index in
equation (\ref{eq: ESI Formula}). First, although our index includes
information on both country $j$'s export shares $s_{jdp}$ and China's market
shares $\chi _{cpd}$, these shares are conceptually different from each
other. In particular, for China what matters is the relative importance of
Chinese exports in total absorption of good $p$ in destination $d$, rather
than how important those Chinese exports are relative to overall Chinese
exports (including exports of other goods and to other destinations).
Second, the model indicates that exposure is proportional to the product of
these shares $\chi _{cdp}$ and $s_{jdp}$ rather than the minimum of the two,
as in the similarity measure (\ref{eq: ESI Formula}). Third, equation (\ref%
{eq: ECI Baseline}) also captures the fact that the shock may be more
disproportionately felt in some sectors or products than others, and it
shows that this heterogeneity can be proxied by the share of Chinese exports
of good $p$ that were destined to the U.S. market relative to other markets
in the baseline year.

We compute our export competition index with the same data used to compute
the Finger-Kreinin export similarity index in (\ref{eq: ESI Formula}).
Because we do not have data on domestic production disaggregated at the HS-6
product level, the denominator of $\chi _{cdp}$ in
equation (\ref{eq: China Absorption Share}) does not include country $d$'s
domestic sales of product $p$.\footnote{For similar reasons,
Chinese domestic sales are not included in the denominator of 
$\Phi _{p}$ in equation (\ref{eq: Relative Shock}), although this is not 
necessarily an issue if one interprets the Chinese export supply elasticity
to be determined independently of Chinese domestic supply (see Appendix \ref{App: MicroFound}).}

Table \ref{tab:top20_eci} reports the top-20 economies by the Export Competition Index with China ($\boldsymbol{ECI}_{j}$). Remember that the figures in the table correspond to the estimated
percentage declines of these countries' exports due to increased Chinese competition. 
The ranking is headed by a set of relatively small and/or specialized economies, but some
notable medium-sized exporters, such as Vietnam or Hungary, emerge. Taiwan is also part of the 
top 10 most affected countries. Overall, Table \ref{tab:top20_eci} points to exposure that is concentrated in smaller, export-specialized economies in Asia and a handful of small island or low-income countries. 

\begin{table}[htbp]
\caption{Top 20 economies by Export Competition Index ($\boldsymbol{ECI}_{j}$)}
\label{tab:top20_eci}\centering
\include{Figures/eci_top20.tex}
\end{table}

Relative to the export similarity in Table \ref{tab:top20_esi_dp}, the composition shifts markedly away from large diversified manufacturing exporters (e.g., South Korea, Japan, Germany, and several large European economies) and toward smaller economies. There is some overlap: Vietnam and Taiwan rank near the top in both tables, and a few additional economies (Thailand, the Philippines, and Poland) appear in both lists, but many of the prominent Table \ref{tab:top20_esi_dp} countries do not appear among the top-20 in Table \ref{tab:top20_eci}. Having said this, our index still predicts non-negligible drops in exports for some of the largest economies in the world. For instance, the exports of Mexico, Turkey, Japan and Indoesia are all predicted to drop between 0.80\% and 0.90\%.

The differences between $ESI_{j}$ and $\boldsymbol{ECI}_{j}$ stem from the fact that the two measures capture different objects. $ESI_{j}$ is an \emph{overlap} index: it is high when a country's export shares across product--destination cells look similar to China's. This naturally elevates large, diversified exporters that ``look like China'' in terms of what they sell and where they sell it. By contrast, $\boldsymbol{ECI}_{j}$ is an \emph{exposure-to-competition} measure: it loads heavily on where the country actually sells and how intense Chinese competition is in those markets. In particular, a country can have a high $\boldsymbol{ECI}_{j}$ if it is specialized in products and destinations where China has high baseline penetration in absorption (i.e., China sells a lot in those destination markets), even if the country's overall export structure is not especially similar to China's. This mechanism is especially relevant for smaller economies: because their exports are often concentrated in a narrow set of product--destination cells, they can be disproportionately affected by a China shock precisely when those cells coincide with markets in which China is a major supplier. Importantly, this can occur even when those product--destination cells are not large in the aggregate composition of Chinese exports; what matters for competitive pressure is China's \emph{market share in the destination} and the country's \emph{own concentration} in that market, not whether the market represents a large fraction of China's total exports.


Panel (a) of Figure \ref{fig:eci_analog} plots our export competition index ($\boldsymbol{ECI}$) against the export similarity index ($ESI$). Despite the differences flagged above, the two indices feature a positive correlation of 0.495.


\begin{figure}[htbp]
\centering
\par
\begin{subfigure}[t]{0.48\linewidth}
    \centering
    \includegraphics[width=\linewidth]{Figures/eci_vs_esi_july_2026.png}
    \caption{$\boldsymbol{ECI}_{j}$ versus $ESI_{j}$}
  \end{subfigure}\hfill 
\begin{subfigure}[t]{0.48\linewidth}
    \centering
    \includegraphics[width=\linewidth]{Figures/updated_eci_product_vs_product_destination_by_continent.png}
    \caption{$\boldsymbol{ECI}_{j}$ versus $\boldsymbol{ECI}_{j}^{prod}$}
  \end{subfigure}
\caption{Export Competition Index: Comparisons}
\label{fig:eci_analog}
\end{figure}



In addition to the baseline index in equation (\ref{eq: ECI Baseline}), we 
also compute a coarser
version of our export competition index at the product level (aggregating
the data first across destinations), which is given by 
\begin{equation}
\boldsymbol{ECI}_{j}^{prod}= 100 \times \sum_{p\in \mathcal{P}}\left(\tilde{\sigma}_{jp}-1\right)\times
\frac{\sum_{d\in \mathcal{W}_{-}}X_{jdp}}{%
X_{j}}\times \frac{%
\sum_{d\in \mathcal{W}_{-}}X_{cdp}}{\sum_{o\in \mathcal{W}}\sum_{d\in 
\mathcal{W}_{-}}X_{odp}}\times\Phi _{p},  \label{eq: ECI_prod}
\end{equation}%
where $\tilde{\sigma}_{jp}$ is the (export-weighted) average demand elasticity
for good $p$ faced by country $j$. This alternative index retains 
the structure of our preferred,
theoretically-grounded index, but it captures better situations in which a
large enough shock (e.g., Liberation Day) may lead to large increases in
Chinese exports to destinations in which it sold relatively less in the
baseline year, something that is inconsistent with the effects of a
first-order shock in our model.

Panel (b) of Figure \ref{fig:eci_analog} plots this alternative export competition
index against our baseline one.\footnote{This alternative index can be downloaded
from our website.} The rank correlation between the two indices is 
high (0.862), but there are some notable differences. In particular, Bangladesh 
and many Eastern European economies (including the Czech Republic, Poland and Hungary)
appear much more exposed. These results are in line with those obtained
under the export similarity measure (see panel (a) of Figure \ref{fig:esi}).

On our web interface, we also explore the robustness of our export competition index
to the use of a more basic measure of the price shock ($\phi_{p}$ instead of 
$\Phi_{p}$) and to the `long' variant using market shares computed from 
2018 data and tariff differences from January 2018 to September 2025. We label
these as $\boldsymbol{ECI}^{basic}$ and $\boldsymbol{ECI}^{long}$. The rank correlation
between $\boldsymbol{ECI}$ and $\boldsymbol{ECI}^{basic}$ is 0.940, while that 
between $\boldsymbol{ECI}$ and $\boldsymbol{ECI}^{long}$ is 0.971.


\subsection{Diagnosing Import Exposure\label{Sec: Import Competition}}

Bystanders of the current U.S.--China trade war may view the constraints
faced by Chinese exporters in selling to the United States not only as a
competitive threat in export markets, but also as a competitive threat in
their \emph{own} market. This reflects a more standard import competition
shock, which in this case is triggered by developments in the United States,
not China. Of course, the flip side of these import competition forces is
the fact that cheaper Chinese imports will reduce the prices faced by
consumers, which is a positive aspect of that trade deflection. For this
type of import threats or opportunities, neither export similarity between a
given country $j$ and China nor our export competition index are
particularly informative.

\subsubsection*{Import Similarity}

We begin by exploring an index of the similarity between the products that
both country $j$ \emph{and} the U.S. import from China in a baseline year.
The idea is that if a given country $j$ imports the same types of goods from
China as the U.S., it is more likely that the U.S--China trade war will
result in a surge of Chinese exports to that market $j$. More specifically,
we define this import similarity index ($ISI$) as 
\begin{equation}
ISI_{j}= 100 \times \sum_{p\in \mathcal{P}}\min \left\{ \frac{M_{jcp}}{%
M_{jc}},\frac{M_{ucp}}{M_{uc}}\right\},  \label{eq: ISI Formula}
\end{equation}%
where $M_{jcp}$ are country $j$'s imports of good $p$ from China, $M_{ucp}$
are U.S. imports of good $p$ from China, and $M_{jc}$ and $M_{uc}$, are
aggregate imports from China of country $j$ and of the U.S., respectively. We
again use 2023 as the reference year for all these flows.

A limitation of the above index is that it may not properly reflect exposure
faced by a country that imports certain goods (though not necessarily from
China), and these goods happen to be goods that the U.S. was
disproportionately importing from China before Liberation Day. The following
alternative import similarity index may better capture these situations:%
\begin{equation}
ISI^{prod}_{j}= 100 \times \sum_{p\in \mathcal{P}}\min \left\{ 
\frac{\sum_{o\in \mathcal{W}}M_{jop}}{M_{j}},\frac{M_{ucp}}{M_{uc}}\right\}
.  \label{eq: ISI Alternative}
\end{equation}%
The key difference is that product $p$'s market share for country $j$ is now
computed aggregating across origin countries (and not just based on imports
from China). The idea is that, if the goal is to predict a surge of imports
into country $j$, it may be better to rely on information on the type of
products in which China is more likely to suffer from excess supply,
regardless of whether China was already exporting those goods to country $j$
in the baseline year.

A\ remaining drawback of the above two measures is that they may not capture
exposure stemming from Chinese exporters redirecting (or \textquotedblleft
dumping\textquotedblright ) into third-country markets certain products with
limited exports to the United States in the benchmark year. This can occur
in sectors where Chinese capacity has expanded rapidly but direct exports to
the United States have been constrained by U.S. tariffs and other regulatory
barriers since the onset of the trade war (if not earlier). Electric
vehicles (EVs) illustrate the point: even though Chinese EV sales in the
U.S. remain limited, they have sparked significant concern in Europe,
despite the consumer gains that would be associated with cheaper EVs. To
some extent, these concerns are linked to the U.S.--China trade war: had
trade between the two economies remained more liberalized, U.S. demand could
have absorbed part of China's surging EV supply, reducing the incentive to
push those exports into other markets. To capture these situations, we
suggest the following index:%
\begin{equation}
ISI^{broad}_{j}= 100 \times \sum_{p\in \mathcal{P}}\min \left\{ \frac{%
M_{jcp}}{M_{j}},\frac{X_{cp}}{X_{c}}\right\},
\label{eq: ISI Alternative 2}
\end{equation}%
which captures similarity between the set of products country $j$ imports
from China and the set of goods China exports to \emph{all} destinations
(not just the U.S.) in the baseline year.\footnote{%
Similarly, one could also envision an even more agnostic measure in which
the first element of the $\min $ operator is replaced with $\sum_{o\in 
\mathcal{W}}M_{jop}/M_{j}$ as in equation (\ref{eq: ISI Alternative}).}

Table \ref{tab:top20_isi_p_china} lists the 20 economies that, according to
our import similarity index in equation (\ref{eq: ISI Formula}), may face
the highest import exposure because they import from China the type of goods
the U.S. imported from China in the baseline year. This list is quite
different from that in Table \ref{tab:top20_esi_dp} and it is dominated by
European countries; it includes only one Asian economy (Japan), while some
of the most impacted countries, such as Canada or Finland, were less
affected by export competition.

\begin{table}[htbp]
\caption{Top 20 economies by Import Similarity ($ISI_{j})
}
\label{tab:top20_isi_p_china}\centering
\include{Figures/isi_top20.tex}
\end{table}

Figure \ref{fig:isi} plots our benchmark import similarity measure against
the two alternative measures developed in equations (\ref{eq: ISI
Alternative}) and (\ref{eq: ISI Alternative 2}). In panel (a), we consider
an import similarity index that compares the products China exported to the
US in the baseline year with the products imported by country $j$ from all
countries (and not just China). The correlation between the two series is
high (0.81), but most countries have a higher import similarity vis {\`a}
vis the U.S. when focusing on imports originating from China, indicating
that the possibility that China starts exporting to those countries (or
ramps its exports of) products it previously sold to the U.S. does not
constitute a notable additional source of exposure for most countries.

When exploring our second alternative import similarity index in equation (%
\ref{eq: ISI Alternative 2}), we find similar results. The correlation
between this alternative index and our baseline one is very high (0.94), so
the set of most impacted countries isted in Table \ref%
{tab:top20_isi_p_china} is largely unaffected. We do observe, however, that
for the vast majority of countries, our baseline import similarity measure
might underestimate exposure related to China potentially increasing their
exports of certain goods that did not feature prominently in U.S. imports
(such as EVs) but are significant for the imports of those economies. The
differences, however, are modest.

\begin{figure}[tbph]
\centering
\par
\begin{subfigure}[t]{0.48\linewidth}
    \centering
    \includegraphics[width=\linewidth]{Figures/isi_p_china_vs_isi_p_all.png}
    \caption{Import Similarity with U.S. (all economies)}
  \end{subfigure}\hfill 
\begin{subfigure}[t]{0.48\linewidth}
    \centering
    \includegraphics[width=\linewidth]{Figures/isi_p_china_vs_isi_p_x_m.png}
    \caption{Import Similarity with Chinese Exports}
  \end{subfigure}
\caption{Alternative Measures of Import Similarity}
\label{fig:isi}
\end{figure}

\subsubsection*{An Import Competition Index}

The above import similarity measures are informative, but they have at least
two limitations. First, they are at best measures of `exposure', so it is
hard to disentangle whether they capture import competition forces -- which
reduce the sales and employment of domestic firms -- or they reflect
potential consumer opportunities derived from access to cheaper imported goods.
Second, as with the Finger-Kreinin export similarity index, these indices
are not theoretically grounded, so it is not clear that the import shares in
equation (\ref{eq: ISI Formula}), as well as the $\min $ operator in the
index, are the relevant objects for assessing import competition threats or
cheaper import opportunities. In Appendix \ref{App: MicroFound}, we use our Armington-CES\
model to shed light on these issues.

We begin by constructing an index of import competition. In particular, we
seek a measure of the extent to which the increased Chinese exports diverted
into a market $j$ will reduce the size of the domestic,
import-competing sector in that country $j$. Naturally, there is a clear 
connection between this
question and our derivations above related to how Chinese exports were
crowding out third-market exporters in export markets. In this case, the
destination market under study coincides with the `origin' country $j$ under
consideration, and the relevant reduced sales $X_{jjp}$ correspond to
domestic sales of good $p$ by producers in $j$, which we simply denote by $S_{jp}$.
The resulting import competition index, when expressed relative to aggregate
domestic sales (i.e., $\sum_{p\in \mathcal{P}}S_{jp}$) reduces to%
\begin{equation}
\boldsymbol{ICI}_{j}= \sum_{p\in \mathcal{P}}\left(\sigma_{op}-1\right)\times s_{jp}\times \chi
_{cjp}\times\Phi _{p},  \label{eq: ICI Baseline}
\end{equation}%
where $s_{jp}=S_{jp}/\sum_{p^{\prime }\in \mathcal{P}}S_{jp^{\prime }}$ is
the share of domestic sales accounted for by sector $p$ domestic producers.

The intuition for this expression is analogous to that behind our export
competition index. The terms $\Phi _{p}$ and $\chi _{cjp}$ reflect the
relative size of the shock in product $p$, and how it will be felt in
country $j$, given how much that country relies on Chinese imports of good $p
$ in the baseline year. Import growth of product $p$ will be higher when the
product $\Phi _{p}\times \chi _{cjp}$ is larger. The term $s_{jp}$ then
captures how this import surge will affect domestic sales, and corresponds
to the relative size of sector $p$ in the overall domestic sales in country $j$.
The predicted import growth may be large, but if there are no (or close to
no) domestic producers of that good, the resulting decline in domestic sales
will be negligible and does not pose a significant threat for the country.

A\ difficulty in implementing the import competition index in (\ref{eq: ICI
Baseline}) is that it requires information on domestic sales at the product
level, which are not readily available to researchers, at least in a unified
dataset covering many countries. With that in mind, we compute $\boldsymbol{%
ICI}_{j}$ proxying the domestic share $s_{jp}=S_{jp}/\sum_{p^{\prime }\in 
\mathcal{P}}S_{jp^{\prime }}$ with the share of overall exports of country $j$
that are accounted by exports of good $p$, or%
\begin{equation}
s_{jp}\simeq s_{jp}^{X}=\frac{\sum_{d\in \mathcal{W}}X_{jdp}}{\sum_{d\in 
\mathcal{W}}\sum_{p^{\prime }\in \mathcal{P}}X_{jdp^{\prime }}}.
\label{eq: X_share}
\end{equation}%
The idea is that countries will tend to feature higher exports in sectors
that are disproportionately large in the industrial structure of a country,
and vice versa. Naturally, variation in the tradability of different
products introduces noise in our proxy, so our results below should be
treated with a grain of salt. Researchers interested in assessing import
competition pressures in specific countries should definitely use those 
product-level domestic sales when that data is available.
For similar data limitation issues, and as in the export competition index, we continue to approximate $\chi_{cop}$ with the share of China in country $o$'s imports of good $p$.

We compute our import competition index $\boldsymbol{ICI}_{j}$ in equation (%
\ref{eq: ICI Baseline}) for all countries in our sample for the year 2023.
%As in the case of
%our export competition index, we also present results with a coarser,
%product-level index of import competition, which replaces $\chi _{cdp}$ in (%
%\ref{eq: ICI Baseline}) with a measure of the share of Chinese total exports
%of good $p$ in worldwide imports of good $p$ (leaving out U.S. imports).%
%\footnote{%
%More specifically, $\chi _{cdp}=X_{cdp}/\sum_{o\in \mathcal{W}}X_{odp}$ is
%replaced with 
%\begin{equation*}
%\mathcal{D}}\sum_{d\in \mathcal{W}_{-}}X_{odp}}\text{.}
%\end{equation*}%
%} The thought behind this robustness test is to consider situations in which
%China may ramp up exports of certain goods in destination markets in which
%it did not sell significant amounts of those goods in the baseline year.
Table \ref{tab:top20_ici} ranks countries by this index $\boldsymbol{ICI}_{j}$. The top of the distribution is again dominated by smaller and/or more specialized economies (e.g., Bangladesh, Haiti, Vietnam, Cambodia, Myanmar, and Pakistan), followed by a mix of Asian economies and a few small countries. This pattern is consistent with $\boldsymbol{ICI}$ being an exposure-to-competition measure: it is high when (i) Chinese import penetration is large in the products affected by the shock and (ii) those products represent a meaningful share of the destination country's domestic, import-competing sales.


\begin{table}[htbp]
\caption{Top 20 economies by Import Competition Index ($\boldsymbol{ICI}_{j}$)}
\label{tab:top20_ici}\centering
\include{Figures/ici_top20.tex}
\end{table}

This ranking differs sharply from an import similarity index ($ISI$), which is an overlap measure based on the composition of a country's imports (typically relative to the set of goods the U.S. imported from China in the baseline year). $ISI$ therefore tends to rank large, high-income import markets whose import baskets are ``U.S.-like'' in product composition. By contrast, $\boldsymbol{ICI}$ is designed to capture the \emph{incidence on domestic producers}: it weights product-level import pressure by the importance of each sector in domestic sales. As a result, countries can have high $\boldsymbol{ICI}$ even if they are not particularly similar to the U.S. in what they import from China; what matters is that domestic producers are concentrated in sectors where Chinese import pressure rises. Consistent with the different objects being measured, the overlap between the top-20 lists for $\boldsymbol{ICI}$ and $ISI$ is very limited (in our data, only the Netherlands appears in both top-20 lists), and the raw cross-country correlation between $ICI$ and the baseline ISI measure is close to zero (0.004), as illustrated in panel (a) of Figure \ref{fig:ici_analog}.

\begin{figure}[htbp]
\centering
\par
\begin{subfigure}[t]{0.48\linewidth}
    \centering
    \includegraphics[width=\linewidth]{Figures/fig_ici_on_isi.pdf}
    \caption{$\boldsymbol{ICI}_{j}$ versus $ISI_{j}$}
  \end{subfigure}\hfill 
\begin{subfigure}[t]{0.48\linewidth}
    \centering
    \includegraphics[width=\linewidth]{Figures/fig_ici_on_ipdi.pdf}
    \caption{$\boldsymbol{ICI}_{j}$ versus $\boldsymbol{CGI}_{j}$}
  \end{subfigure}
\caption{Import Similarity, Import Competition and Import Price Decline Indices}
\label{fig:ici_analog}
\end{figure}



\subsubsection*{A Consumer Gains Index}

As shown above, the Armington-CES\ model can be used to
assess the implications of the U.S.--China trade war for firm sales.
Interestingly, it can also be employed to evaluate its consequences for
consumer prices. In Appendix \ref{App: MicroFound}, we show that the percentage decline in import
prices in a given market $j$ can be approximated, up to a first
order, with the following expression%
\begin{equation}
\boldsymbol{CGI}_{j}= 100 \times \sum_{p\in \mathcal{P}} s_{jp}^{M}\times \chi
_{cjp}^{M}\times \Phi _{p},
\label{eq:CGI}
\end{equation}%
where $\Phi _{p}$ is identical to that in the expressions for $\boldsymbol{%
ECI}_{j}$ and $\boldsymbol{ICI}_{j}$ in equations (\ref{eq: ECI Baseline})
and (\ref{eq: ICI Baseline}), and where $\chi _{cjp}^{M}$ is the Chinese
share in total imports of good $p$ in country $j$ (which was already our
proxy for $\chi _{cjp}$ when empirically implementing our export and import competition
indices). The term $s_{jp}^{M}$ in the above equation for $\boldsymbol{CGI}%
_{jc}$ is the share of country $j$'s imports accounted for by sector $p$.
Intuitively, up to scalars, the predicted price decline of imports of good $p
$ is proportional to the surge in imports of that good (which is captured
by the product $\Phi _{p}\times \chi _{cjp}^{M}$), and then the spending shares that are
relevant to compute the aggregate consumer gains associated with those lower
import prices are the import shares $s_{jp}^{M}$.

Table \ref{tab:top20_ipdi} ranks countries by the predicted import-price decline from a China-specific supply shock, $\boldsymbol{CGI}_{j}$. To reiterate, this measures mechanically puts weight on (i) how large China is in a country's import basket in the affected products ($\chi^{M}_{cjp}$) and (ii) how important those products are in the country's overall import spending ($s^{M}_{jp}$). 


\begin{table}[htbp]
\caption{Top 20 economies by Consumer Gains Index ($\boldsymbol{CGI}_{j}$)}
\label{tab:top20_ipdi}\centering
\include{Figures/cgi_top20.tex}
\end{table}

Compared to Table \ref{tab:top20_ici}, which captures producer-side exposure and therefore loads on the importance of the affected sectors in domestic, import-competing activity rather than in import expenditures), the list in Table \ref{tab:top20_ipdi} tilts toward economies that are highly import-dependent on China in the relevant sectors, so it is not surprising to see several smaller, import-reliant economies (and a few larger ones with sizable China import shares in the impacted goods, such as Australia) rise to the top, even if their domestic producers are not especially concentrated in those sectors. By contrast, Table~\ref{tab:top20_ici} features economies where the shock maps strongly into domestic producer displacement because the affected sectors represent a meaningful share of domestic sales. The limited overlap between the two tables (e.g., Hong Kong SAR and a few Asian economies) is thus exactly what the theory implies: consumer price gains are governed by import-expenditure weights, whereas producer pressures are governed by the domestic sectoral footprint of the same goods.\textbf{[EDIT WHEN FINAL TABLES ARE AVAILABLE]}

In panel (b) of Figure \ref{fig:ici_analog} we plot $\boldsymbol{ICI}_{j}$ and $\boldsymbol{CGI}_{j}$ against each other for all countries in our sample. The correlation between the two series is positive but very small (0.133).{[EDIT WHEN FINAL TABLES ARE AVAILABLE]}


\begin{table}[htbp]
\caption{Top 20 economies by Consumer Gains Index (Downstream)($\boldsymbol{CGI}_{j}^{downstream}$)}
\label{tab:top20_ipdi_down}\centering
\include{Figures/cgi_downstream_top20.tex}
\end{table}


\textbf{[ADD CGI FOR FINAL GOODS]}. 
Operationally, we distinguish inputs from final goods using the value-chain
position of products in the spirit of \cite{antras2012measuring}: products
that are more \emph{upstream} (i.e., typically used as inputs into further
production stages before reaching final demand) are classified as
intermediates, whereas more \emph{downstream} products are closer to final
absorption. We label product $p$ as a final good if
its upstreamness measure is lower than 2. We define the set of final products by $\mathcal{P}%
^{F}\subseteq \mathcal{P}$.\footnote{%
In the implementation, this requires mapping the upstreamness measure to our
HS product classification and defining the input set $\mathcal{P}^{F}\subseteq \mathcal{P}$ at the HS6 level.}

\subsection{Diagnosing Sourcing Exposure}

A distinctive feature of the current U.S.--China trade war is that the
excess Chinese supply displaced from the U.S. market need not be
concentrated in final or consumption goods. A substantial share of the
affected products are \emph{intermediate inputs} -- parts, components, and
materials used by firms in subsequent stages of production \citep[cf.,][]{bown2019measuring}. When Chinese exports of such inputs are redirected toward third markets, the resulting
shock has a fundamentally different interpretation than in the case of final
goods: it may intensify import competition for domestic input producers, but
it can also relax sourcing constraints and reduce costs for domestic
downstream users of those inputs, thereby strengthening their
competitiveness.

It is straightforward to develop a version of our import competition index in 
equation \eqref{eq: ICI Baseline} that restricts attention to intermediate inputs.
We do so in Appendix \ref{App: MicroFound} and calculate it for all countries
in Appendix XYZ. In the main text, we focus on developing an index of the 
sourcing gains associated with cheaper Chinese intermediate inputs. This is 
not simply a version of our consumer gains index restricted to intermediate
inputs because properly accounting for sourcing gains requires the use of 
information on input-output links between sectors.

\subsubsection*{A Sourcing Gains Index}

To account for the gains from cheaper inputs, we develop a variant of our 
Armington-CES model that incorporates input-output links. In particular, 
production of each good $p$ now uses inputs of product-type $z$ with a 
cost share $a_{pz}$. Cheaper intermediate inputs generate heterogeneous
effects on downstream sectors depending on the drop in prices of the 
particular inputs that each downstream sector uses. To aggregate these 
effects, we trace the implications of these lower costs for the exports
of a country's exports to all destinations (other than the U.S. and
China). In Appendix \ref{App: MicroFound}, we then show that 
the percentage response of country $o$'s exports to Chinese cheaper 
intermediate inputs is given by

\begin{equation}
 \boldsymbol{SGI}_{j}=100\times\sum_{d\in\mathcal{W}_{-}}\frac{X_{jd}}{X_{j}}\sum_{p\in\mathcal{P}}\left(\sigma_{dp}-1\right)\times s_{jdp}\times\sum_{z\in \mathcal{P}^{I}}a_{pz}\times  \chi _{cjz}^{M}\times \Phi _{z}.
\label{eq: SGI}
\end{equation}

\textbf{[DONE UP TO HERE]}

\begin{equation}
\boldsymbol{RPI}_{ciu}\equiv \sum_{p\in \mathcal{P}}s_{iup}\left( \sum_{z\in 
\mathcal{P}^{I}}a_{pz}\times \Phi _{z}\times \chi _{ciz}^{M}\right) .
\label{eq: RPI Appendix}
\end{equation}


show that the percentage fall in imported input prices in a given
destination market $d$ can be approximated, up to a first order, with the
following expression%
\begin{equation}
\boldsymbol{SPDI}_{dc}=100 \times \sum_{p\in \mathcal{P}^{I}}\Phi _{p}\times \chi
_{cdp}^{M}\times s_{dp}^{M,I},
\label{eq: SPDI index}
\end{equation}%
where $\Phi _{p}$ and $\chi _{cdp}^{M}$ are (again) identical to those employed
in the previous indices, and where $%
s_{dp}^{M,I}$ is the share of country $d$'s imports of intermediate inputs
accounted for by sector $p$. In the model in Appendix \ref{App: MicroFound}, the term $%
\boldsymbol{SPDI}_{dc}$ captures (up to a scalar common across countries) the
percentage decline in intermediate input import prices associated with the
U.S. policy that has rendered Chinese exports of inputs cheaper in other markets.

Table \ref{tab:top20_spdi} reports the Sourcing Price Decline Index ($\boldsymbol{SPDI}_{dc}$), led by Tonga, the Kyrgyz Republic, Russia, Tajikistan, and the Solomon Islands (with Zimbabwe, Cuba, and Montenegro also near the top). Compared with the sourcing competition index $\boldsymbol{SCI}$ in Table \ref{tab:top20_sci} -- whose top countries are Bangladesh, Vietnam, Myanmar, Cambodia, and Tonga -- $\boldsymbol{SPDI}$ shifts attention from \emph{producer competition} to \emph{import-cost exposure}. Economically, $\boldsymbol{SPDI}$ is high where a country's imported intermediate-input basket is concentrated in products for which China's predicted price declines (and/or China's supply presence) are particularly large, so the main effect is cheaper inputs for downstream firms rather than displacement of domestic intermediate producers. That is why $\boldsymbol{SPDI}$ highlights several small, import-dependent economies and a few countries with concentrated intermediate import profiles, while $\boldsymbol{SCI}$ highlights places where domestic upstream producers are most likely to be challenged. The overlap for Tonga (top in $\boldsymbol{SPDI}$ and also top-5 in $\boldsymbol{SCI}$) and Cambodia (top-10 in both) is consistent with some economies facing both forces simultaneously: tougher upstream competition alongside cheaper imported inputs for downstream production. Panel (b) of Figure \ref{fig:sci_analog} plots $\boldsymbol{SCI}_{dc}$ against $\boldsymbol{SPDI}_{dc}$ and shows that there is virtually no correlation between the two series (-0.005).


\begin{table}[htbp]
\caption{Top 20 economies by $\boldsymbol{SGI}_{dc}$}
\label{tab:top20_spdi}\centering
\include{Figures/sgi_top20.tex}
\end{table}



\subsection{Diagnosing Rerouting}

Our final \emph{rerouting} channel captures a distinct margin of adjustment to
U.S.--China trade frictions: production (or, more mechanically, recorded
trade flows) shifts away from China and toward third countries that become
platforms for serving the U.S. market. In its ``real'' form, this reflects
relocation of downstream stages -- often final assembly -- to a third
country, while upstream stages and component production remain in China. In
its ``statistical'' form, it reflects re-exporting/transshipment or limited
transformation, whereby goods enter the U.S. after passing through an
intermediary location.

This mechanism does not simply reflect \emph{sourcing opportunities},
according to which cheaper Chinese intermediates can improve competitiveness
broadly, potentially raising exports across \emph{all} destinations.
Instead, under rerouting, the salient prediction is \emph{%
destination-specific}: countries that expand exports to the U.S. in sectors
where they are simultaneously deepening linkages to Chinese supply chains
(whether through inputs or even the same HS product) are those most exposed
to rerouting dynamics.

\subsubsection*{Import-Export Similarity}

To operationalize this idea, we first construct an Import--Export Similarity
Index (IESI) that compares (i) the composition of country $i$'s exports to
the U.S. with (ii) the composition of country $i$'s imports from China. When
computed at a fine product level (e.g., HS-6), the index is informative
about rerouting of very similar products (including potential re-exporting).
When computed at a more aggregated level (e.g., say HS-4 or HS-2 levels), it
is better suited to capture relocation patterns in which China continues to
supply parts and components while final-stage production is performed
elsewhere.\footnote{%
When available, direct information on input-output links would naturally
lead to more precise indices of sourcing and rerouting opportunities.}
Formally, we define:%
\begin{equation}
IESI_{ciu}=100 \times \sum_{p\in \mathcal{P}}\min \left\{ \frac{X_{iup}}{X_{iu}},\frac{%
M_{icp}}{M_{ic}}\right\} ,  \label{eq: Imp Exp Similarity}
\end{equation}%
where $X_{iu}$ are country $i$'s total exports to the U.S. and $M_{ic}$ are
country $i$'s total imports from China. As a baseline, we compute this Import--Export similarity index at the HS-4
level (with 2023 data), but we will explore its robustness to alternative
disaggregation levels.

Table \ref{tab:top20_iesi} lists the 20 economies with the highest amount of
rerouting according to our similarity index in equation (\ref{eq: Imp Exp
Similarity}). The list includes a few economies -- Hong Kong SAR, Taiwan Province of China, or
Mexico -- that have been mentioned in recent years as ``connectors'' between
China and the U.S., but it also includes a number of Asian and European
countries that, according to our index, have the potential to serve as
``intermediaries'' of this bilateral relationship in future years.

\begin{table}[htbp]
\caption{Top 20 economies by Import-Export China-U.S. Similarity ($%
IESI_{ciu}$)}
\label{tab:top20_iesi}\centering
\include{Figures/iesi_top20.tex}
\end{table}

Figure \ref{fig:iesi} plots our benchmark HS-4 similarity index against two
alternative indices computed at a higher disaggregation level (HS-6 in panel
(a)) and at a lower disaggregation level (HS-2 in panel (b)). The results
are qualitatively similar in both cases, with correlations of 0.89 and 0.80
in panel (a) and (b), respectively. As a result, the top countries in terms
of rerouting potential are the same ones according to all three measures.

\begin{figure}[tbph]
\centering
\par
\begin{subfigure}[t]{0.48\linewidth}
    \centering
    \includegraphics[width=\linewidth]{Figures/iesi_hs6_vs_iesi_hs4_v2.png}
    \caption{Rerouting at HS-4 vs. HS-6 level}
  \end{subfigure}\hfill 
\begin{subfigure}[t]{0.48\linewidth}
    \centering
    \includegraphics[width=\linewidth]{Figures/iesi_hs2_vs_iesi_hs4_v2.png}
    \caption{Rerouting at HS-4 vs. HS-2 level}
  \end{subfigure}
\caption{Rerouting at Alternative Disaggregation Levels}
\label{fig:iesi}
\end{figure}

\subsubsection*{Rerouting Potential Index}

In addition to the Import--Export Similarity Index, we construct a
theoretically-motivated rerouting index that links cheaper Chinese
intermediate inputs to the potential expansion of country $i$'s exports to
the U.S.\ in the same sectors. The key idea is that rerouting is a
\emph{destination-specific} adjustment margin: countries are more likely to
expand exports to the U.S.\ in sectors in which their marginal costs fall
because they source intermediate inputs from China.

In principle, the marginal-cost reduction for an exported product $p$ depends
on the full input--output (IO) structure. In Appendix \ref{App: MicroFound}, we show that, under
Armington CES demand in the U.S.\ and a China-specific price shock, the
percentage change in exports of product $p$ from $i$ to the U.S.\ is
proportional (up to common elasticities and pass-through parameters) to the
fall in the intermediate input price index relevant for producing $p$ in $i$.
This yields the IO-based rerouting potential index
\begin{equation}
\boldsymbol{RPI}_{icu}\equiv 100 \times\sum_{p\in\mathcal{P}}s_{iup}\left(
\sum_{z\in\mathcal{P}^{I}}a_{pz}\times \Phi_{z}\times \chi_{ciz}^{M}\right),
\label{eq: RPI Baseline}
\end{equation}%
where $s_{iup}\equiv X_{iup}/X_{iu}$ is the baseline share of product $p$ in
country $i$'s exports to the U.S., $\chi_{ciz}^{M}$ is China's share in
country $i$'s imports of intermediate input $z$, $\Phi_{z}$ is the
China-only shock term defined above, and $a_{pz}$ are IO weights mapping
output $p$ into its intermediate input requirements.

In the absence of a harmonized IO table at the HS level, we propose a set of
trade-based approximations that impose structure on $a_{pz}$ using HS nesting.
For $\ell\in\{6,4,2\}$, define the HS-$\ell$ neighborhood of output good $p$ as
$\mathcal{P}^{I,(\ell)}(p)\equiv\{z\in\mathcal{P}^{I}:HS_{\ell}(z)=HS_{\ell}(p)\}$,
and proxy $a_{pz}$ by country $i$'s baseline intermediate import shares within
that neighborhood:
\begin{equation*}
s_{iz\mid p}^{M,I,(\ell)}\equiv \frac{M_{iz}^{I}}{\sum_{z^{\prime}\in\mathcal{P}^{I,(\ell)}(p)}M_{iz^{\prime}}^{I}},
\qquad q\in\mathcal{P}^{I,(\ell)}(p).
\end{equation*}%
This yields a product-specific sourcing decline index,
\begin{equation*}
\boldsymbol{SSDI}_{icp}^{(\ell)}\equiv 100 \times\sum_{q\in\mathcal{P}^{I,(\ell)}(p)}
\Phi_{q}\times \chi_{ciq}^{M}\times s_{iq\mid p}^{M,I,(\ell)},
\end{equation*}%
and an empirically implementable rerouting potential index,
\begin{equation}
\boldsymbol{RPI}_{icu}^{(\ell)}\equiv \sum_{p\in\mathcal{P}} s_{iup}\times
\boldsymbol{SSDI}_{icp}^{(\ell)}.
\label{eq: RPI HS}
\end{equation}%
The three levels of $\ell$ correspond to progressively broader assumptions
about input linkages: HS-6 captures ``own-product'' rerouting (including
re-exporting), while HS-4 and HS-2 allow for parts-and-components linkages
within a narrower or broader industry. Finally, if all outputs were to use a
common economy-wide input bundle, then $\boldsymbol{SSDI}_{icp}^{(\ell)}$
would collapse to an aggregate sourcing price decline measure and the export
composition weights $s_{iup}$ would no longer matter for the predicted
percentage change in $X_{iu}$.\footnote{We have also computed a version of the $\boldsymbol{RPI}$ index that includes \emph{all} goods and not just intermediate inputs in constructing the price decline component of the index. Although there
are some differences in the top countries under both measures, the correlation between the two indices is high (0.56).}

Table \ref{tab:top20_rpi_hs4} reports the top-20 countries in terms of our Rerouting Potential Index ($\boldsymbol{RPI}$). The ranking is led by Laos, Nepal, Haiti and Bangladesh (with Cambodia and Vietnam also near the top). The limited overlap between this list and that for the Import-Export Similarity Index in Table \ref{tab:top20_iesi} shows is intuitive: IESI is a \emph{pure composition-overlap} statistic between what a country exports to the U.S. and what it imports from China, whereas RPI further scales that overlap by where the model predicts the China-to-U.S. sourcing reallocation/price pressure is most intense across HS4 sectors. As a result, countries with highly concentrated export baskets and strong reliance on Chinese inputs in a few ``high-rerouting'' sectors can rise sharply in the $\boldsymbol{RPI}$ ranking even if their overall overlap is not among the very highest in $IESI$. Where the lists \emph{do} intersect (e.g., Hong Kong SAR, Taiwan Province of China), it flags economies that are simultaneously tightly linked to Chinese inputs and have U.S.-oriented exports in the same industries, making them natural candidates for trade rerouting effects.


\begin{table}[htbp]
\caption{Top 20 economies by $\boldsymbol{RPI}_{ciu}$ (HS4)}
\label{tab:top20_rpi_hs4}\centering
\include{Figures/rpi_top20.tex}
\end{table}

Panel (a) and (b) of Figure \ref{fig:rpi_analog} plot our baseline  $\boldsymbol{RPI}_{ciu}$ index computed at the HS-4 level against the import-export similarity index $IESI_{ciu}$ and the $\boldsymbol{RPI}_{ciu}$ index computed at the HS-6 level, respectively. Panel (a) shows close to no correlation (-0.040) between $\boldsymbol{RPI}_{ciu}$ and $IESI_{ciu}$, while the level of disaggregation of the $\boldsymbol{RPI}_{ciu}$ seems less material, as the correlation between the two series in panel (b) is very high (0.949).


\begin{figure}[htbp]
\centering
\par
\begin{subfigure}[t]{0.48\linewidth}
    \centering
    \includegraphics[width=\linewidth]{Figures/fig_rpi_hs4_input_on_iesi_hs4.pdf}
    \caption{$\boldsymbol{RPI}_{ciu}$ (HS-4) versus $IESI_{ciu}$}
  \end{subfigure}\hfill 
\begin{subfigure}[t]{0.48\linewidth}
    \centering
    \includegraphics[width=\linewidth]{Figures/fig_rpi_hs4_input_on_rpi_hs6.pdf}
    \caption{$\boldsymbol{RPI}_{ciu}$ (HS-4) versus $\boldsymbol{RPI}_{ciu}$ (HS-6)}
  \end{subfigure}
\caption{Import-Export Similarity and Rerouting Potential}
\label{fig:rpi_analog}
\end{figure}




\begin{comment}

\subsection{Summary}

We finally propose combining our four measures of exposure to shed light on
the extent to which countries are exposed to either competitive threats or
to opportunities resulting from the ongoing U.S.--China trade war.

One way to do so is via a Venn diagram which assigns countries to categories
-- export competition, import exposure, sourcing exposure, rerouting
potential -- if that country is in the top quartile of that category's
baseline similarity measure. With 193 countries, this corresponds to being
in the top 48 of a measure's ranking. The resulting graph in Figure \ref%
{fig:venn} illustrates some countries -- such as United Arab Emirates, Saudi
Arabia or Russia -- that face both significant export \emph{and} import
exposure, while ranking much lower on sourcing exposure and rerouting
potential. Conversely, a number of countries -- such as Serbia, Morocco or
Egypt -- face significant sourcing exposure or rerouting potential, while
appearing more insulated from export and import exposures.

\begin{figure}[htbp]
\caption{Venn Diagram of Third-Market Spillovers}
\label{fig:venn}
\begin{center}
{\includegraphics[width=.85\textwidth]{Figures/Venn_Diagram_Pol.pdf}}
\end{center}
\end{figure}

We next attempt to project the independent information of the four indices
into a lower number of dimensions. First, we combine the export and import
competition indices and create an ``average competitive threat'' index. To
do so, we first standardize our two indices and then take a simple average
of the two. This is obviously somewhat arbitrary, but without the guidance
of a quantitative model, it is hard to assess which of these two channels
should impact an economy more severely. We then construct an analogous
``average opportunity index'' by averaging a standardized version of our
sourcing and rerouting opportunity indices.

Figure \ref{fig:threat_opp} plots this two average indices against each
other. Consistent with the message of the Venn diagram in Figure \ref%
{fig:venn}, there exists an interesting heterogeneity in how economies are
impacted by threats and opportunities. Countries below the 45-degree line
appear to be relatively negatively impacted, while those above that line are
relatively positively affected. Interesting, this variation persists even
when zooming on income groups (advanced economies, emerging markets, and
developing economies). In Figure \ref{fig:threat_opp_regions} of Appendix %
\ref{App: MultiCountry} we show that a similar heterogeneity remains within
regional areas (Europe, Asia, Africa, America and Oceania).

\begin{figure}[htbp]
\caption{Average Competitive Threats vs. Opportunities}
\label{fig:threat_opp}\centering
\par
\begin{center}
{\includegraphics[width=.85\textwidth]{Figures/threat_opp_index.png}}
\end{center}
\end{figure}

We finally go one step further and aggregate all four indices into a single
measure, which we label the ``net threat'' index, and is constructed by
subtracting the average opportunity index from the average threat index.
Figure \ref{fig:net_threat} depicts this index for the top 100 exporting
economies.\footnote{%
In Figure \ref{fig:net_threat_all} of Appendix \ref{App: Robustness Indices}%
, we plot the net threat index for all 193 countries in our sample.} We
necessarily take this index with a (big) grain of salt because, as we have
mentioned above, the cardinal aspects of our indices are not directly
income- or welfare-relevant. Having said this, the message that emerges from
the figure is in line with those in Figures \ref{fig:venn} and \ref%
{fig:threat_opp}. Some countries, most notably Australia, United Arab
Emirates, Saudi Arabia and Russia, are significantly exposed to export or
import competition shock, while our similarity indices indicate that they
are not likely to benefit from significant sourcing or rerouting
opportunities. Conversely, a number of European economies and Mexico appear
to be much better positioned to benefit from the worldwide reallocations
triggered by the continuing U.S--China trade war.

\begin{figure}[htbp]
\caption{Net Threat for Top 100 Exporters}
\label{fig:net_threat}
\begin{center}
{\includegraphics[width=.9\textwidth]{Figures/Net_Threat_Top_100.png}}
\end{center}
\end{figure}
\end{comment}


% \section{The Age of Geoeconomics}

% \label{sec:geoeconomics}

% So far our analysis has taken the observed policy changes in recent months
% and years as given. As anticipated in the Introduction, many aspects of the
% post--2018 (and especially post--2023) reallocation are not only about costs
% and efficiency, but increasingly about \emph{policy risk} and \emph{strategic%
% } considerations. This shift is tightly connected to a broader
% transformation in the way governments and firms think about cross-border
% linkages: trade and supply chains are now routinely evaluated through the
% lens of economic security, resilience, and geopolitical exposure, rather
% than exclusively through the lens of comparative advantage and efficiency %
% \citep{wto2023world,Aiyar2023Geo,alfaro_trade_2025,alekseev2025trade}.

% These developments have helped revive and reshape an old idea: that
% international trade is not merely an economic relationship, but also a
% source of geopolitical leverage. Classic contributions in political economy
% and international relations emphasized that asymmetric trade dependence can
% be weaponized %
% \citep{hirschman1945national,Baldwin1985EconomicStatecraft,krasner1976state}%
% . Yet, the 1980s, 1990s and early 2000s brought about a remarkable increase
% in global economic interdependence, and in recent times, we have also
% observed a notable expansion in the breadth of policy instruments used to
% pursue geopolitical objectives -- from targeted tariffs and export controls
% to sanctions, investment screening, and financial restrictions %
% \citep{FarrellNewman2019WeaponizedInterdependence,ClaytonMaggioriSchreger2023FrameworkGeoeconomics,BeckoOConnor2025StrategicDisIntegration,BeckoGrossmanHelpman2025OptimalTariffsGeopoliticalAlignment}%
% . A rapidly growing literature now studies these issues under the umbrella
% term \emph{Geoeconomics}: the strategic use of economic relationships and
% economic policy tools to influence outcomes in the global economy %
% \citep{mohr-trebesch2025,ClaytonMaggioriSchreger2025NBERMacro}.

% \subsection{Geopolitical Alignment and International Trade}

% A central empirical implication of the geoeconomic turn is that geopolitical 
% \emph{alignment} should become increasingly predictive of trade outcomes.
% Recent work by \cite{gopinath2025changing} documents precisely this pattern:
% global linkages have increasingly reorganized along geopolitical lines, with
% rising ``friend-shoring'' and the growing importance of connector economies
% that intermediate between rival blocs.

% Figure \ref{Fig: Geo Alignment} provides two complementary perspectives on
% these developments. Panel (a), borrowed directly from \cite%
% {gopinath2025changing}, shows that, in 2022-24, trade growth slowed
% disproportionately more among countries that belong to politically distant
% blocs. Panel (b) illustrates the increasing connection between geopolitics
% and trade using a standard gravity-style approach based on political
% distance. Specifically, we estimate a gravity model in which bilateral goods
% imports are regressed on geographic distance and a measure of political
% distance -- the ideal-point distance in U.N. voting -- interacted with year
% dummies, with a demanding set of fixed effects.\footnote{This approach -- and the results -- mimic what is done on FDI by \citet{Aiyar-etal2024}, and deliver results similar to those obtained by \cite{bosone2024beyond} using a more sophisticated structural gravity framework.} The resulting time series of
% coefficients shows that political distance has become progressively more
% predictive of bilateral trade: trade has become more \emph{club-like}, in
% the sense that geopolitically aligned countries trade disproportionately
% more with each other, while trade between politically distant countries is
% comparatively depressed. 

% \begin{figure}[htbp]
% \caption{Geopolitical Alignment and Trade Flows}
% \label{Fig: Geo Alignment}\centering
% \begin{subfigure}[t]{0.40\linewidth}
%     \centering
%     \includegraphics[width=\linewidth]{Figures/GopinathPresbitero.pdf}
%     \caption{Differential Trade Growth}
   
%   \end{subfigure}\hfill 
% \begin{subfigure}[t]{0.53\linewidth}
%     \centering
%     \includegraphics[width=\linewidth]{Figures/trade_ipd.png}
%     \caption{Ideal Point Distance and Trade Flows}
   
%   \end{subfigure}
% \par
% \bigskip 
% \begin{minipage}{.95\linewidth}
% {\footnotesize \textbf{Notes:} Panel (a) plots average trade growth during
% 2022Q2-2024Q1 minus the average trade growth during
% 2017Q1-2022Q1 within and between blocs. Bloc
% definition is based on a hypothetical Western bloc
% centered around the U.S. and Europe and a hypothetical
% Eastern bloc centered around China and Russia. Panel (b) plots the coefficient of the ideal
% point distance (IPD) \citep{Bailey-etal_2017} estimated in a gravity model
% where bilateral goods imports are regressed against geographical distance,
% the IPD interacted with yearly dummies, destination $\times$ year, origin $%
% \times$ year and destination $\times$ origin fixed effects. The shaded area
% represents the 90 percent confidence band. Standard errors are clustered at
% the origin-destination level. \par
% \textbf{Sources:} \cite{gopinath2025changing}, \cite{moyer2024relative}; \cite{Bailey-etal_2017}; IMF, \emph{International trade in goods by partner country dataset} (IMTS).}
% \end{minipage}
% \end{figure}

% This alignment channel is increasingly prominent in the empirical trade
% literature. For instance, \cite{kleinmanliuRedding2024friends} show that
% economic dependence and political alignment co-move: countries that are more
% dependent on a partner tend to realign toward that partner. Relatedly, \cite%
% {nerilaine2023sovereign} documents a sizable military-alliance effect on
% trade, while earlier work in political science and economics linked
% political regimes and international relations to trade integration %
% \citep{mansfieldmilnerrosendorff2000freetotrade,martinmayertheonig2008maketradenotwar}%
% . In recent work, \cite{broner2025hegemonic} and \cite{broner2025hegemony}
% emphasize that global power structures (a world with a hegemon versus a
% multipolar world) shape alignment, and that alignment (proxied by
% treaty-making) is a meaningful predictor of subsequent trade integration.

% Although this \textquotedblleft alignment channel\textquotedblright\ is
% predominant in the recent literature, it may have weakened in the most
% recent policy environment: since Trump's reelection, tariff actions have
% increasingly looked broad-based rather than narrowly targeted at
% geopolitical adversaries. Figure \ref{fig: effective tariffs} confirms this
% by plotting the recent evolution of U.S. effective tariff rates across broad
% partner groups in 2025. The graph shows a pronounced upward shift after
% March 2025. Yet, the increase is \emph{not} concentrated on geopolitical
% adversaries: tariff escalation is broad-based and, in several episodes,
% appears to fall at least as heavily -- if not more heavily -- on nonaligned
% partners and on close allies. One reading of this is that U.S. trade policy
% has entered a phase that is less organized around coalition-building and
% more unilateral and transactional, with tariffs used as a general instrument
% of leverage rather than primarily as a tool to discipline geopolitical
% rivals.

% \begin{figure}[htbp]
% \caption{U.S. Effective Tariff Rates in Recent Months}
% \label{fig: effective tariffs}
% \begin{center}
% {\includegraphics[width=.5\textwidth]{Figures/blocs_etr_weighted_noCHN.png}}
% \end{center}
% \par
% \bigskip 
% \begin{minipage}{.80\linewidth}
% {\footnotesize \textbf{Notes:} The chart plots the effective tariff rates, defined as collected duties over imports, imposed by the U.S. on three groups of countries. China is excluded in the China-leaning bloc. \par
% \textbf{Sources:} IMF World Economic Outlook and World Bank World Development
% Indicators.}
% \end{minipage}
% \end{figure}

% \subsection{Measuring Geoeconomic Power}

% The alignment patterns above naturally raise a measurement question: if
% geoeconomics is about the strategic use of economic relationships, what is
% the relevant notion of \emph{geoeconomic power}, and how can it be
% quantified in practice?

% A first approach -- closer to the international relations tradition -- is to
% measure a country's overall capability to project influence in global
% affairs using broad indices of national power. This is the basis of the
% \textquotedblleft Global Power Index\textquotedblright\ (\emph{GPI})
% developed by \cite{moyer2024relative}, which attempts to measure the \emph{%
% relative} national power of countries. Their measure constitutes a weighted
% sum of the relative dominance of a country in various sub-indicators, which
% themselves are grouped into five broad categories: (i) military, (ii)
% population, (iii) economy, (iv) technology, and (v) diplomatic networks.
% Panel (a) of Figure \ref{Fig: Both Powers} illustrates a marked shift in the
% global distribution of power: China accounts for a remarkably larger share
% of global power today than it did in 1960, largely at the expense of the
% United States, which remains the most powerful country in the world, though
% it has arguably lost its ``hegemonic'' dominance.\footnote{%
% The GPI index is available on a much longer-run basis, going back to 1816.
% As emphasized by \cite{broner2025hegemonic}, the early post-WWII era
% featured a sharp rise in U.S. dominance, which widened further after the end
% of the Cold War and then began to narrow again as China's relative power
% accelerated in the 1990s and 2000s.}

% \begin{figure}[htbp]
% \caption{Two Measures of Geoeconomic Power}
% \label{Fig: Both Powers}\centering
% \begin{subfigure}[t]{0.50\linewidth}
%     \centering
%     \includegraphics[width=\linewidth]{Figures/GPI_bars.png}
%     \caption{Global Power Index of \cite{moyer2024relative}}
   
%   \end{subfigure}\hfill 
% \begin{subfigure}[t]{0.50\linewidth}
%     \centering
%     \includegraphics[width=\linewidth]{Figures/IP_bars.png}
%     \caption{International Power of \cite{liu2025international}}
   
%   \end{subfigure}
% \par
% \bigskip 
% \begin{minipage}{.95\linewidth}
% {\footnotesize \textbf{Notes:} Panels (a) and (b) plot the value of the Global Power Index and of the measure of International Power, respectively, for specific countries or regions. Europe includes France, Germany, Italy, Spain, and United Kingdom. Other EMs includes the following emerging markets: Brazil, Egypt, Ethiopia, India, Indonesia, Iran, Russia, Saudi Arabia, South Africa, and  United Arab Emirates. \par
% \textbf{Sources:} \cite{moyer2024relative}; \cite{liu2025international}.}
% \end{minipage}
% \end{figure}

% A second approach -- closer to the geoeconomics tradition %
% \citep[cf.][]{hirschman1945national} -- is to focus directly on the leverage
% created by \emph{asymmetric} trade dependence. Building on this idea, \cite%
% {liu2025international} relate international power to the net leverage a
% country has over another country stemming from asymmetric trade dependence.
% The measure first quantifies how much a country (say country $A$) depends on
% imports from another country (say country $B$), giving more weight to
% sectors that are harder to substitute away from (i.e., where trade is less
% elastic). Power of $B$ over $A$ is then defined as $A$'s dependence on $B$
% minus $B$'s dependence on $A$, capturing who would be hurt more -- and thus
% who has more leverage -- if trade were disrupted. To reduce this bilateral
% variable to a country-level one, in panel (b) of Figure \ref{Fig: Both
% Powers} we construct a weighted average international power of a \emph{%
% coercer} country as a weighted sum of the bilateral power of this coercer
% over all other \emph{target} countries, where the weights are proportional
% to the GDPs of the target countries. As is clear from the picture, the
% resulting patterns are very similar to those produced by the Global Power
% index of \cite{moyer2024relative}, though the rise of China and the drop of
% the U.S. is much more pronounced in the International Power measure of \cite%
% {liu2025international}. In fact, according to our weighted average of their
% bilateral power index, China is currently much more powerful than any other
% country in the world.

% \subsection{A Rising Correlation between Size and Power}

% The changing distribution of power is often described as a transition from a
% U.S.-led unipolar system toward a more multipolar (or even bipolar) world.
% While that description captures an important part of the story, the data
% also reveal a deeper and less appreciated development: economic size and
% geopolitical power have become more tightly linked over time, despite the
% relative decline of the United States in recent decades.

% To document this, we combine the two empirical measures of power discussed
% above with cross-country measures of GDP. First, we use the Global Power
% Index (GPI) of \cite{moyer2024relative}, but to avoid a mechanical
% correlation, we strip out the GPI components that directly measure economic
% size, and reassign their weights proportionately to other sub-indicators.%
% \footnote{%
% This is feasible to do with the spreadsheet \cite{moyer2024relative} made
% available at
% https://dataverse.harvard.edu/dataset.xhtml?persistentId=doi:10.7910/DVN/QRGCWI.%
% } The left panel of Figure \ref{fig: Correlation Data} shows the
% cross-country correlation between the GPI and GDP (at market prices) from
% 1995 until 2020. It is clear that even when removing economic variables as
% inputs to the GPI, the correlation between these two series remains \emph{%
% very} high, well above 0.95. More interestingly, the data indicates a
% non-negligible increase in the correlation between these two variables in
% the period 1995-2020. One may worry that this pattern is driven by
% individual countries (i.e., Chinese growth in both dimensions during this
% period), but the pattern remains when removing China from the correlation
% and also when computing a Spearman rank correlation (see Appendix \ref{App:
% Corr Robust}).

% Using instead the trade-leverage measure of \cite{liu2025international},
% which is designed to be less mechanically tied to GDP, we also observe a
% clear upward trend in the GDP-power correlation since around 2000, even
% though the level correlation remains much lower than for the GPI. More
% specifically, the right panel of Figure \ref{fig: Correlation Data} shows
% this correlation in the period 1995-2020, and it is always lower than 0.4.
% Interestingly, and after an initial drop between 1995 and 2000, this
% correlation has been growing over time, and peaked at the onset of the
% U.S.--China trade tensions. As we document in Appendix \ref{App: Corr Robust}%
% , the Spearman rank correlation between this notion of power and economic
% size has also steadily increased since the year 2000.

% \begin{figure}[htbp]
% \centering
% \par
% \begin{subfigure}[t]{0.48\linewidth}
%     \centering
%     \includegraphics[width=\linewidth]{Figures/corr_gpi_gdp_1995_2020_pearson.pdf}
%     \caption{Correlation between Global Power Index \citep{moyer2024relative} and GDP}
%     \label{fig: GPI}
%   \end{subfigure}\hfill 
% \begin{subfigure}[t]{0.48\linewidth}
%     \centering
%     \includegraphics[width=\linewidth]{Figures/corr_weightedpower_gdpshare_pearson.pdf}
%     \caption{Correlation between International Power \citep{liu2025international} and GDP }
%     \label{fig: LiuYang}
%   \end{subfigure}
% \caption{Correlation between Geopolitical Power and Economic Size
% (1995-2020) }
% \label{fig: Correlation Data}
% \end{figure}

% This rising size-power correlation is conceptually distinct from a simple
% \textquotedblleft one hegemon versus two great powers\textquotedblright\
% narrative. Even holding fixed how many major powers the world has, a tighter
% mapping from economic size to coercive capacity changes the strategic
% environment for trade policy: powerful countries become more likely to be
% the same ones with greater ability to influence global prices and
% allocations, and thus to have stronger incentives (and greater room) to act
% unilaterally. In the next section, we develop an analytical framework to
% formally trace the implications of this development for the sustainability
% of a cooperative world trade system.

% \section{Geopolitical Power and Trade Wars: Analytical Notes\label%
% {Sec:Theory}}

% In this section, we develop a simple framework that clarifies why global
% free trade becomes harder to sustain when geopolitical influence and
% economic size are tightly aligned. We begin with a two-country
% terms-of-trade model in which larger markets have stronger unilateral
% incentives to use tariffs, making cooperation fragile when transfers are
% unavailable. We then enrich the setup by allowing governments to exert
% geopolitical influence abroad, captured by a parameter that makes foreign
% policymakers partially internalize the welfare of powerful states. The key
% implication is stark: free trade is most viable when economic size and
% geopolitical power are \emph{not} strongly positively correlated---because,
% when the largest countries are also the most influential, they can often
% secure favorable outcomes even without cooperation and therefore have less
% to gain from binding themselves to free-trade commitments.

% \subsection{Baseline Model}

% \subsubsection*{Economic Environment}

% We begin with a simple model with two countries, Home and Foreign, along the
% lines of \cite{antras2011foreign}. Each country is populated by a continuum
% of measure one of individuals with identical preferences:%
% \begin{equation}
% u^{j}=c_{0}^{j}+\sum_{s=1}^{2}u_{s}^{j}\left( c_{s}^{j}\right) \text{, \ \ }%
% j=H,F \text{,}  \label{U}
% \end{equation}%
% where $u_{s}^{j}\left( \cdot \right) $ is increasing and strictly concave.
% All individuals supply inelastically one unit of labor. Good $0$ serves as
% the num{\'e}raire, is costlessly traded and not subject to tariffs. Its
% world and domestic price is normalized to $1$. It is produced one to one
% with labor everywhere in the world, which pins down the wage rate to $1$ in
% all countries. The other goods can also be traded internationally, but may
% be subject to taxes. We shall also assume that good 1 is a \textquotedblleft
% natural export\textquotedblright\ of Home, while good 2 is a
% \textquotedblleft natural export\textquotedblright\ of Foreign. More
% precisely, we assume that trade policy cannot revert \textquotedblleft
% natural\textquotedblright\ comparative advantage patterns.

% We will focus on a world in which countries may tax trade flows to shift the
% terms of trade in their favor, but for simplicity, we will focus on the case
% in which countries only use import tariffs (and not export taxes) to achieve
% that goal. Let $p_{s}^{W}$ denote the world untaxed price of good $s$. This
% corresponds to the price paid by consumers in the exporting country, since
% there are no taxes nor transport costs involved in that transaction. On the
% other hand, the domestic price in the importing country $j$ will be given by 
% $\left( 1+t_{s}^{j}\right) p_{s}^{W}$, where $t_{s}^{j}$ denotes the import
% tariff set by country $j$.

% Non-num{\'e}raire goods are produced combining labor and sector-specific
% capital according to a constant returns to scale technology. Let $\Pi
% _{s}^{j}$ be the aggregate rent accruing to sector $s$ specific factor in
% country $j$. Capital is evenly distributed among the measure $1$ of workers
% in each country.

% A convenient property of the quasilinear representation of preferences in (%
% \ref{U}) is that aggregate welfare in country $j$ can be written as $%
% v^{j}\left( \mathbf{p}\right) =I^{j}\left( \mathbf{p}\right) +S^{j}\left( 
% \mathbf{p}\right) ,$ where $I^{j}\left( \mathbf{p}\right) $ denotes
% aggregate income in country $j$, $S^{j}\left( \mathbf{p}\right) $ denotes
% consumer surplus, and $\mathbf{p}$ is the vector of domestic prices $\mathbf{%
% p\equiv }\left( 1,p_{1}^{j},p_{2}^{j}\right) $. Given our assumptions, we
% can further write aggregate income in country $j$ as $I^{j}=1+\Pi
% _{1}^{j}\left( p_{1}^{j}\right) +\Pi _{2}^{j}\left( p_{2}^{j}\right)
% +R^{j}\left( \mathbf{t},\mathbf{p}^{W}\right) $, where $R^{j}\left( \mathbf{t%
% },\mathbf{p}^{W}\right) $ is tariff revenue in country $j$.\footnote{%
% This tariff revenue is given by:%
% \begin{equation*}
% R^{j}\left( \mathbf{t},\mathbf{p}^{W}\right) =\left\{ 
% \begin{tabular}{lll}
% $t_{2}^{H}p_{2}^{W}\left( c_{2}^{H}\left( p_{2}^{H}\right) -y_{2}^{H}\left(
% p_{2}^{H}\right) \right) $ & $\text{if}$ & $j=H$ \\ 
% $t_{1}^{F}p_{1}^{W}\left( c_{1}^{F}\left( p_{1}^{F}\right) -y_{1}^{F}\left(
% p_{1}^{F}\right) \right) $ & if & $j=F$%
% \end{tabular}
% \text{.} \ \ \ \ \right.
% \end{equation*}%
% } Note also that consumer surplus is simply given by $S^{j}\left( \mathbf{p}%
% \right) =\sum_{s=1}^{2}\left[ u_{s}^{j}\left( c_{s}^{j}\left(
% p_{s}^{j}\right) \right) -p_{s}^{j}c_{s}^{j}\left( p_{s}^{j}\right) \right] $%
% .

% Given quasilinear preferences, we can study trade policy sector by sector.
% We can focus on the problem of a single country setting tariffs on the good
% that is a natural import for that country. In doing so, it is important to
% remember that the world price $p_{s}^{W}\,$is endogenous and must satisfy
% market clearing, or%
% \begin{equation}
% M_{s}^{j}\left( p_{s}^{j}\right) \equiv c_{s}^{j}\left( p_{s}^{j}\right)
% -y_{s}^{j}\left( p_{s}^{j}\right) =y_{s}^{k}\left( p_{s}^{W}\right)
% -c_{s}^{k}\left( p_{s}^{W}\right) \equiv X_{s}^{k}\left( p_{s}^{W}\right) 
% \text{ for }j\neq k\text{.}  \label{MktClearing}
% \end{equation}

% \subsubsection*{Optimal Tariffs: The Benevolent Social Planner Benchmark}

% Ignoring irrelevant terms, we can write the problem of the government in
% country $j$ as:%
% \begin{equation*}
% \max_{t_{s}^{j}}\text{ }\Pi _{s}^{j}\left( p_{s}^{j}\right)
% +t_{s}^{j}p_{s}^{W}\left( c_{s}^{j}\left( p_{s}^{j}\right) -y_{s}^{j}\left(
% p_{s}^{j}\right) \right) +u_{s}^{j}\left( c_{s}^{j}\left( p_{s}^{j}\right)
% \right) -p_{s}^{j}c_{s}^{j}\left( p_{s}^{j}\right)
% \end{equation*}%
% subject to $p_{s}^{j}=\left( 1+t_{s}^{j}\right) p_{s}^{W}$ and $%
% c_{s}^{j}\left( p_{s}^{j}\right) -y_{s}^{j}\left( p_{s}^{j}\right)
% =y_{s}^{k}\left( p_{s}^{W}\right) -c_{s}^{k}\left( p_{s}^{W}\right) $, for $%
% j\neq k$. This program delivers the following standard optimal tariff
% formula: 
% \begin{equation}
% \hat{t}_{s}^{j}=\frac{1}{\xi _{s}^{k}}\equiv \frac{X_{s}^{k}\left(
% p_{s}^{W}\right) }{p_{s}^{W}X_{s}^{k\prime }\left( p_{s}^{W}\right) }.
% \label{eq: Optimal Tariff}
% \end{equation}%
% In words, a country's (country $j$'s) optimal import tariff is equal to\ the
% inverse of the export supply elasticity it faces.

% For most parameterizations, the larger is a country (say Home) relative to
% the other country (say Foreign), the lower is the export supply elasticity
% it will face, and thus the higher will its optimal tariff be. To see this,
% note that the export supply elasticity of country $j$ faces from $k$ can be
% expressed as 
% \begin{equation*}
% \xi _{s}^{k}=\frac{\eta _{s}^{S,k}+\eta _{s}^{D,k}\gamma _{s}^{k}}{1-\gamma
% _{s}^{k}},
% \end{equation*}%
% where $\eta _{s}^{S,k}$ is country $k$'s supply elasticity in sector $s$, $%
% \eta _{s}^{D,k}$ is country $k$'s demand elasticity in sector $s$, and $%
% \gamma _{s}^{k}$ is country $k$'s own trade share in sector $s$, or $\gamma
% _{s}^{k}=c_{s}^{k}\left( p_{s}^{W}\right) /y_{s}^{k}\left( p_{s}^{W}\right) $%
% . Provided that demand and supply elasticities are non-increasing in the
% size of a country (a weak assumption), and provided that the own trade share
% is lower for countries that are small \emph{relative} to their trading
% partners (also a weak assumption), then it is clear that $\xi _{s}^{k}$ will
% be decreasing in the size of country $k$ and increasing in the size of
% country $j$.

% To illustrate this, consider the parametric example developed by \cite%
% {antras2011foreign} with linear demand and supply functions. In particular,
% assume that the utility functions $u_{s}^{j}$ in (\ref{U}) are quadratic, so
% that demand functions are linear and given by%
% \begin{equation}
% c_{s}^{H}\left( p_{s}^{H}\right) =\lambda ^{H}\left( \alpha _{s}^{H}-\beta
% p_{s}^{H}\right) \text{; \ \ \ }c_{s}^{F}\left( p_{s}^{F}\right) =\lambda
% ^{F}\left( \alpha _{s}^{F}-\beta p_{s}^{F}\right) \text{, \ \ for }s=1,2%
% \text{,}  \label{DemandLinear}
% \end{equation}%
% where $\alpha _{2}^{H}=\alpha _{1}^{F}=\bar{\alpha}>\underline{\alpha }%
% =\alpha _{1}^{H}=\alpha _{2}^{F}$. Furthermore, the rent functions $\Pi
% _{s}^{j}$ are also assumed to be quadratic, thus leading to linear supply
% functions in each country:%
% \begin{equation}
% y_{s}^{H}\left( p_{s}^{H}\right) =\lambda ^{H}\left( a+bp_{s}^{H}\right) ;%
% \text{ \ \ \ }y_{s}^{F}\left( p_{s}^{F}\right) =\lambda ^{F}\left(
% a+bp_{s}^{F}\right) \text{, \ \ for }s=1,2\text{.}  \label{SupplyLinear}
% \end{equation}%
% Notice that both countries share similar demand and supply functions, but
% Home demand is disproportionately large in sector 2, while Foreign demand is
% disproportionately large in sector 1. As a result, Home is a natural
% importer\ in sector 2, while Foreign is a natural importer in sector 1. In
% order for equilibrium autarky prices to be positive, we need to assume that $%
% \bar{\alpha}>\underline{\alpha }>a$. Note that the parameters $\lambda ^{H}$
% and $\lambda ^{F}$ capture the market sizes of the Home and Foreign
% countries, respectively.

% For this case, we find that the optimal tariff of country $j$ is given by%
% \begin{equation*}
% \hat{t}^{j}=\frac{\bar{\alpha}-\underline{\alpha }}{\left( \bar{\alpha}%
% -a\right) +\left( \underline{\alpha }-a\right) \left( \lambda ^{k}/\lambda
% ^{j}+1\right) },
% \end{equation*}%
% and clearly $\hat{t}^{j}$ is increasing in $\lambda ^{j}$ and decreasing in $%
% \lambda ^{k}$.

% We can also consider the incentives of countries to seek a free trade
% agreement that avoids this non-cooperative Nash equilibrium in tariffs.
% Because free trade maximizes world welfare, if countries could include
% transfers of the outside good in the negotiations, such a free trade
% agreement would always be reached. Nevertheless, when these transfers are
% not allowed (as in the bulk of the international trade policy literature),
% there are situations in which a disproportionately large country may want to 
% \emph{opt out} of a free trade agreement. To see this, consider the limit in
% which Foreign becomes arbitrarily small. In such a case, Foreign will face
% an infinitely large Home export supply elasticity, which from equation (\ref%
% {eq: Optimal Tariff}) will lead to $\hat{t}_{1}^{F}=0$.\footnote{%
% In our linear example above, if $\lambda ^{H}/\lambda ^{F}\rightarrow \infty 
% $, then $\hat{t}_{2}^{H}=\left( \bar{\alpha}-\underline{\alpha }\right)
% /\left( \bar{\alpha}+\underline{\alpha }-2a\right) $, while $\hat{t}%
% _{1}^{F}=0$.} Because Home is already obtaining free trade under the
% non-cooperative equilibrium, why would it give up its ability to set its
% positive optimal import tariff? More generally, and as first pointed out by 
% \cite{johnson1953optimum}, if one of the two countries is disproportionately
% larger than the other, it will tend to oppose a free trade agreement, or it
% will renege on an existing one. For our parametric example with linear
% demand and supply in equations (\ref{DemandLinear}) and (\ref{SupplyLinear}%
% ), one can show that a free trade agreement will only be signed if $\lambda
% ^{H}/\lambda ^{F}$ is sufficiently close to $1$, or $\Lambda \geq \lambda
% ^{H}/\lambda ^{F}\geq 1/\Lambda $, where $\Lambda =1.377$ regardless of
% other parameter values. This range is plotted in Figure \ref{Fig: Size}.

% \begin{figure}[htbp]
% \caption{Country Size and the Feasibility of a Free Trade Agreement}
% \label{Fig: Size}\centering%
% \includegraphics[scale=0.55,page=1]{Figures/GraphsAntrasPresbitero.pdf} %
% \subcaption*{\footnotesize {\textbf{\emph{Notes:}} The figure plots the
% feasibility of a free trade agreement between Home and Foreign for various
% values of the size of Home and Foreign when demand and supply are given by
% equations (\ref{DemandLinear}) and (\ref{SupplyLinear}).}}
% \end{figure}

% \subsubsection*{Geoeconomic Power}

% We next expand the analysis to consider scenarios in which countries might
% wield their geopolitical power to shape the policies of other countries to
% their own advantage. Following the approach in \cite{antras2011foreign} , we
% model \emph{foreign influence} as involving costly actions (by a government)
% that probabilistically affect election outcomes in other countries.\footnote{%
% See also \cite{bubeckmarinov2017} and \cite{antras2025exporting}, and the
% empirical evidence in \cite{bubeckmarinov2019}.} This approach delivers a
% very simple and intuitive (subgame-perfect) \emph{equilibrium}. In
% particular, each country implements policies which end up maximizing a
% weighted sum of domestic welfare and of foreign welfare. This occurs because
% political parties in a given country (say Foreign) tilt their platforms in
% favor of the Home country in order to forestall a potential Home
% intervention in support of their domestic opponents. The influence power of
% countries (i.e., the extent to which they can have foreign governments
% internalize the effects of their policies on them) depends on parameters
% shaping the marginal cost and the marginal benefit of exerting influence. To
% simplify matters, in this paper we will treat this weight $\mu ^{j}$ as
% fixed, and we will interpret it as a measure of the \emph{geopolitical power}
% of country $j$.

% We now return to the non-cooperative determination of import tariffs.
% Ignoring irrelevant terms, we can now write the Home government problem as:%
% \begin{eqnarray*}
% &&\max_{t_{s}^{j}}\text{ }\Pi _{s}^{j}\left( p_{s}^{j}\right)
% +t_{s}^{j}p_{s}^{W}\left( c_{s}^{j}\left( p_{s}^{H}\right) -y_{s}^{j}\left(
% p_{s}^{j}\right) \right) +u_{s}^{j}\left( c_{s}^{j}\left( p_{s}^{j}\right)
% \right) -p_{s}^{j}c_{s}^{j}\left( p_{s}^{j}\right) \\
% &&\qquad +\mu ^{k}\left[ \Pi _{s}^{k}\left( p_{s}^{W}\right)
% +u_{s}^{k}\left( c_{s}^{k}\left( p_{s}^{W}\right) \right)
% -p_{s}^{W}c_{s}^{k}\left( p_{s}^{W}\right) \right] \text{.}
% \end{eqnarray*}%
% This program delivers the following general solution for the optimal tariff
% of country $j$ on country $k$:%
% \begin{equation}
% \tilde{t}_{s}^{j}=\left( 1-\mu ^{k}\right) \frac{1}{\xi _{s}^{k}}\equiv
% \left( 1-\mu ^{k}\right) \frac{X_{s}^{k}\left( p_{s}^{W}\right) }{%
% p_{s}^{W}X_{s}^{k\prime }\left( p_{s}^{W}\right) }.
% \label{eq: Import Tariff Power}
% \end{equation}%
% Note that when $\mu ^{k}=0$, country $k$ does not exert any influence on $j $%
% , and naturally we obtain the standard optimal tariff formula in equation (%
% \ref{eq: Optimal Tariff}). Conversely, when $\mu ^{k}=1$, $k$'s influence is
% so powerful that it precludes any terms-of-trade manipulation on the part of
% country $j$. In such a case, we have that foreign influence leads to free
% trade in country $k$'s import sector. This is not surprising because, in
% such a case, country $j$ would be choosing $\hat{t}_{s}^{j}$ to maximize
% aggregate world welfare, and this is achieved with free trade.
% Interestingly, in cases in which $\mu ^{k}>1$, the model predicts that
% country $j$ will adopt an import \emph{subsidy} instead of an import tax.

% Outside these limit cases, characterizing how geopolitical power shapes the
% non-cooperative equilibrium is complicated by the fact that unless export
% supply elasticities are constant, they will be shaped by equilibrium trade
% flows, which are in turn also affected by geopolitical power. In the linear
% example in equations (\ref{DemandLinear}) and (\ref{SupplyLinear}), however,
% we obtain particularly sharp results. In such a case, the Home and Foreign
% optimal tariffs are given by:%
% \begin{equation*}
% \tilde{t}^{j}=\frac{\left( 1-\mu ^{k}\right) \left( \overline{\alpha }-%
% \underline{\alpha }\right) }{\left( \overline{\alpha }-a\right) +\left( 
% \underline{\alpha }-a\right) \left( \lambda ^{k}/\lambda ^{j}+1-\mu
% ^{k}\right) }.
% \end{equation*}%
% Simple differentiation indicates that each country $j$'s tariff is
% increasing in its relative size $\lambda ^{j}/\lambda ^{k}$ and decreasing
% in the power of the other country $\mu ^{k}$.

% \subsubsection*{Geopolitical Power and the Sustainability of Free Trade}

% We next explore how the worldwide distribution of geopolitical power shapes
% the viability of a free trade agreement in the absence of means to transfer
% utility. Consider first a situation in which the two countries are of equal
% size $\lambda ^{H}=\lambda ^{F}$ but they differ in their relative
% geopolitical power. When Home (say $\mu ^{H}\rightarrow 1$) is very powerful
% and Foreign is very weak (say $\mu ^{F}\rightarrow 0$), then Home will
% clearly opt out of a free trade agreement because it already obtains free
% trade under the non-cooperative equilibrium, while giving up its
% unilaterally optimal tariff would be costly. More generally, and as pointed
% out by \cite{antras2011foreign}, if one of the two countries is
% disproportionately more geopolitically powerful than the other, it will tend
% to oppose a free trade agreement, or it will renege on an existing one. For
% our parametric example with linear demand and supply, one can show that a
% free trade agreement will only be signed if $\mu ^{H}/\mu ^{F}$ is
% sufficiently close to $1$, or $\bar{\Psi}\geq \mu ^{H}/\mu ^{F}\geq 
% \underline{\Psi },$ where these bounds do not depend on other parameter
% values if $\lambda^j$ is the same in all countries 
% \citep[see][for
% details]{antras2011foreign}. Figure \ref{Fig: Miu Symmetry} plots this
% region of the parameter space whenever $\left( \mu ^{H},\mu ^{F}\right) \in %
% \left[ 0,1\right] \times \left[ 0,1\right] $. The figure clearly illustrates
% that the existence of power imbalances across countries may lead the
% powerful country to block a socially efficient free trade agreement. This
% result embodies a strong intuition: if, absent an agreement, weak countries
% are already forced to acquiesce with the interests of powerful countries,
% the latter have little to gain from concerted moves to world welfare
% maximizing policies.

% \begin{figure}[htbp]
% \caption{Power Asymmetries and the Feasibility of a Free Trade Agreement}
% \label{Fig: Miu Symmetry}\centering%
% \includegraphics[scale=0.55,page=2]{Figures/GraphsAntrasPresbitero.pdf} %
% \subcaption*{\footnotesize {\textbf{\emph{Notes:}} The figure plots the
% feasibility of a free trade agreement between Home and Foreign for various
% values of the geopolitical power of Home and Foreign, assuming both
% countries have the same size, and when demand and supply are given by
% equations (\ref{DemandLinear}) and (\ref{SupplyLinear}).}}
% \end{figure}

% We can also use an analogous graph to formally study the interaction of
% economic size and geopolitical power in affecting the viability of free
% trade agreements. In particular, consider the case in which $\lambda
% ^{H}>\lambda ^{F}$, so Home is a relatively large country. In such a case,
% and as illustrated above in Figure \ref{Fig: Size}, free trade might not be
% achieved even when influence power is balanced (e.g., when $\mu ^{H}=\mu
% ^{F}=0$) because Home will block it. In those situations, free trade will
% only be achievable whenever, despite its large size, Home's power is small
% relative to Foreign's so that, in the absence of agreement, Home is forced
% to impose a small tariff on Foreign. Hence, when $\lambda ^{H}/\lambda ^{F}$
% is large, a free trade agreement will only be signed when $\mu ^{F}\ $is
% high relative to $\mu ^{H}$. This is illustrated in Figure \ref{Fig: Miu
% Asymmetry} for the case $\lambda ^{H}/\lambda ^{F}=1.5$. The same figure
% depicts the case in which\ Home is small relative to Foreign ($\lambda
% ^{H}/\lambda ^{F}=2/3$), for which the feasibility of a free trade agreement
% again rests on the relatively small country being relatively powerful.

% \begin{figure}[htbp]
% \caption{Country Size, Power Asymmetries, and the Feasibility of a Free
% Trade Agreement}
% \label{Fig: Miu Asymmetry}\centering%
% \includegraphics[scale=0.55,page=3]{Figures/GraphsAntrasPresbitero.pdf} %
% \subcaption*{\footnotesize {\textbf{\emph{Notes:}} The figure plots the
% feasibility of a free trade agreement between Home and Foreign for various
% values of the size of Home and Foreign and for two relative sizes of Home
% and Foreign when demand and supply are given by equations (\ref{DemandLinear}%
% ) and (\ref{SupplyLinear}).}}
% \end{figure}

% In sum, the feasibility of a free trade requires a \emph{negative}
% correlation between size and influence power. For the case of linear demand
% and supply functions in equations (\ref{DemandLinear}) and (\ref%
% {SupplyLinear}), \cite{antras2011foreign} showed that this result can be
% summarized more formally as follows (see their Proposition 5):

% \begin{proposition}
% \label{Prop: Agreements}For each country $j=H,F$, there exist a threshold $%
% \widetilde{\mu }^{j}\in \left[ 0,1\right] $ such that if $\mu ^{j}>%
% \widetilde{\mu }^{j}$, then country $j$ will find it beneficial to opt out
% of a free trade agreement. Furthermore, the threshold $\widetilde{\mu }^{j}$
% is necessarily decreasing in the relative size of country $j$ (i.e., $%
% \partial \widetilde{\mu }^{j}/\partial \left( \lambda ^{j}/\lambda
% ^{k}\right) <0$, implying that the feasibility of a free trade agreement
% requires a negative correlation between size and geopolitical power.
% \end{proposition}

% \subsubsection*{An Extension to Many Countries}

% We finally extend our analysis to a world with many countries that can all
% trade bilaterally as well as exert geopolitical pressures on each other. We
% first do so in the simplest possible way to illustrate the robustness of the
% negative correlation result in Proposition \ref{Prop: Agreements}. With that
% in mind, we consider a world with $N$ countries in which, in each country,
% individuals consume a num{\'e}raire good and $2\times \left( N-1\right) \ $%
% other goods. More specifically, preferences in country $j$ are identical to
% those in equation (\ref{U}) except that we replace $2$ with $2\times \left(
% N-1\right) $ in the upper limit of the summation sign in the utility
% function. The key simplifying assumption is that we assume that each of
% these $2\times \left( N-1\right) $ goods relevant for welfare in country $j$
% is only consumed and produced in one of the foreign countries $k\neq j$.%
% \footnote{%
% Overall, the world economy produces $\left( N-1\right) \times N+1$ goods,
% including the num{\'e}raire good.} Together with the quasilinearity of
% preferences, this ensures that (i) welfare in country $j$ as additively
% separable across the welfare associated with exchanging each of the $2\times
% \left( N-1\right) $ with the $N-1$ foreign countries; (ii) we can study
% optimal trade policy not only sector by sector, but also country by country.

% In the presence of geopolitical pressures, the optimal tariff of country $j$
% on country $k$ is again given by equation (\ref{eq: Import Tariff Power}),
% so $\tilde{t}^{j\rightarrow k}=\left( 1-\mu ^{k}\right) /\xi ^{k\rightarrow
% j}$, where $\mu ^{k}$ is the geopolitical power of $k$ and $\xi
% ^{k\rightarrow j}$ is the export supply elasticity associated with country $%
% k $'s good that is only consumed in $j$ (other than in $k$). In our
% recurring linear example, we further obtain%
% \begin{equation*}
% \tilde{t}^{j\rightarrow k}=\frac{\left( 1-\mu ^{k}\right) \left( \overline{%
% \alpha }-\underline{\alpha }\right) }{\left( \overline{\alpha }-a\right)
% +\left( \underline{\alpha }-a\right) \left( \lambda ^{k}/\lambda ^{j}+1-\mu
% ^{k}\right) }.
% \end{equation*}%
% Again, relatively larger countries impose larger tariffs, while relatively
% larger or more powerful countries face relatively lower tariffs.

% In Appendix \ref{App: MultiCountry}, we further consider the welfare
% implications of moving from a global tariff war in which each country sets
% their optimal tariffs vis {\`a} vis all other countries in the world to a
% world of global free trade. As in our two-country model, relatively large
% and powerful countries will be tempted to opt out of such a global free
% trade agreement.\footnote{%
% More precisely, there continues to exist a threshold $\widetilde{\mu }%
% ^{j}\in \left[ 0,1\right] $ such that if $\mu ^{j}>\widetilde{\mu }^{j}$,
% then country $j$ will find it beneficial to opt out of a free trade
% agreement with $k$, and the threshold $\widetilde{\mu }^{j}$ is necessarily
% decreasing in the relative size of country $j$ (i.e., $\partial \widetilde{%
% \mu }^{j}/\partial \left( \lambda ^{j}/\lambda ^{k}\right) <0$.} More
% generally, the sustainability of global cooperation is endangered by a
% positive correlation between country size and geopolitical power. To
% illustrate this, we develop numerical simulations of our linear model for a
% world of $N=20$ countries in which the parameters $\lambda ^{j}\,$and $\mu
% ^{j}$ are generated assuming varying degrees of correlation $\rho $ between $%
% \lambda ^{j}$ and $\mu ^{j}$, while bounding these parameters by $\lambda
% ^{j}\in \lbrack 0.8,1.25]$ and $\mu ^{j}\in \lbrack 0,0.5]$.\footnote{%
% The details of these numerical simulations are in Appendix \ref{App:
% MultiCountry}. For each correlation value, we run $10,000$ independent draws
% of the model. The resulting standard errors are negligibly small.} The left
% panel of Figure \ref{fig:app_twopanels} plots the estimated probability that 
% \emph{all }countries gain from a shift to global trade against the average
% realized Pearson correlation within each $\rho $ grid point, while the right
% panel of the figure plots the average share of countries that gain from such
% a global free trade agreement. As is clear from the figures, the probability
% of all countries benefiting from trade is close to $1$ when size and power
% are perfectly negatively correlated, while it goes to $0$ when they are
% perfectly positively correlated. More generally, free trade finds more
% support the lower is the correlation between size and power. 
% \begin{figure}[htbp]
% \centering
% \par
% \begin{subfigure}[t]{0.48\linewidth}
%     \centering
%     \includegraphics[width=\linewidth]{Figures/prob_all_positive_vs_corr_mu0_0p5.pdf}
%     \caption{Probability all countries gain from free trade}
%     \label{fig:app_allpos}
%   \end{subfigure}\hfill 
% \begin{subfigure}[t]{0.48\linewidth}
%     \centering
%     \includegraphics[width=\linewidth]{Figures/share_positive_vs_corr_mu0_0p5.pdf}
%     \caption{Share of countries that gain from free trade}
%     \label{fig:app_sharepos}
%   \end{subfigure}
% \caption{Correlation between Size and Power and the Feasibility of Global
% Free Trade}
% \label{fig:app_twopanels}
% \end{figure}

% Admittedly, the above results were obtained in a very special setting with
% quasilinear preferences and separability of tariff effects across sectors
% and country-pairs. Nevertheless, the main mechanism generating Figure \ref%
% {fig:app_twopanels} continues to operate in more general multi-country,
% multi-good models general equilibrium models. With many potential suppliers
% and destinations, a country's tariff on a given good affects the \emph{%
% residual} foreign export supply faced by a country, and also create
% third-market spillovers by shifting other countries' terms of trade and
% fiscal revenues. In this environment, optimal trade taxes can still be
% understood as balancing an \textquotedblleft
% opportunistic\textquotedblright\ terms-of-trade gain against the induced
% changes in other countries' real incomes, weighted by the extent to which
% the government internalizes foreign welfare, as emphasized by \cite%
% {adao2024world} in a general neoclassical setting. In those more general
% environments, powerful countries will tend to face non-cooperative tariff
% equilibria that are not too damaging to their welfare, and if these
% countries are also large, they will find it costly to give up their
% incentive to set sizable tariffs to shift their terms of trade in their
% favor.

% \subsubsection*{Taking Stock}

% In sum, the theory and the evidence point to the same vulnerability of the
% world trading system. In the model, the sustainability of free trade hinges
% on a \emph{negative} relationship between size and influence: as the
% correlation between economic size and geopolitical power rises, it becomes
% increasingly likely that at least one large, powerful country prefers to opt
% out of (or renege on) cooperation, and in the multi-country extension the
% probability that all countries support global free trade falls sharply as
% this correlation turns positive. Consistent with this mechanism, the
% evidence discussed in Section \ref{sec:geoeconomics} show that the
% correlation between GDP and two distinct measures of geoeconomic power has
% trended upward since the mid-1990s---even when stripping economic inputs
% from a broad \textquotedblleft power index,\textquotedblright\ and even
% using a trade-leverage concept of power designed to be less mechanically
% tied to size. Together, these results suggest that the recent drift toward a
% world in which economic size and geopolitical power increasingly coincide
% is, in itself, a force undermining the durability of global free trade.

\section{Conclusions \label{Sec: Conclusion}}

The post-2018 period is best understood not as a generalized reversal of
globalization, but as a reorganization of trade and production linkages in
response to policy shocks and strategic rivalry. Using up-to-date
information on the evolving tariff regime and recent trade flows, we have
documented that the U.S.--China trade war has generated a sharp reallocation
of U.S. import sourcing away from China and a corresponding reorientation of
Chinese exports toward other destinations, with an acceleration in
adjustment during the most recent escalation following Liberation Day. These
shifts have reshaped competitive conditions globally, even when aggregate
world trade remains broadly stable.

A central message of the paper is that third-country outcomes cannot be read
solely through the lens of simple trade diversion. To discipline an
interpretation of these third-market effects, we have proposed a taxonomy
that distinguishes six (non-mutually exclusive) margins -- export competition,
import competition, import price declines, sourcing competition, sourcing price
declines, and rerouting potential -- and operationalized them with diagnostic indices.
This approach clarifies why some economies appear to benefit on one margin
while facing intensified pressure on another.

We also highlighted that the current episode is unfolding in an ``age of
geoeconomics,'' where trade policy tools are increasingly used not just to
pursue economic objectives but to generate leverage and manage geopolitical
risk. Empirically, geopolitical alignment has become more predictive of
bilateral trade outcomes, consistent with the rise of ``friend-shoring''; at
the same time, the most recent U.S. tariff actions appear broad-based,
suggesting a more unilateral and transactional deployment of tariffs as a
general instrument of leverage. We then brought these developments to the
question of whether cooperative trade outcomes can be sustained when policy
incentives are shaped by both market power and geopolitical influence.

The analytical framework in Section \ref{Sec:Theory} makes precise a
mechanism linking geopolitical influence to the stability of low-tariff
cooperation. When influential states can secure relatively favorable
outcomes outside agreements -- because other governments partially
internalize their welfare -- those states have less to gain from binding
commitments and stronger incentives to act unilaterally. In that
environment, the durability of a cooperative trading system becomes
especially sensitive to how geopolitical power is distributed and how
tightly it aligns with economic size. Consistent with this mechanism, the
evidence reviewed in Section \ref{sec:geoeconomics} indicates that the
empirical alignment between economic size and geoeconomic power has
strengthened over time, reinforcing the paper's broader interpretation of
why strains on rules-based cooperation have become more acute.

Looking ahead, our results suggest that restoring a stable, rules-based
trading order will require more than marginal tariff rollbacks. Even when
world aggregate remains broadly stable, we are likely entering a phase of 
\emph{fragmented interdependence}, in which the structure of integration is
likely to become more segmented, with greater weight on geopolitically
aligned partners and on ``connector'' economies that facilitate reallocation
and rerouting.

\newpage

\bibliographystyle{chicago}
\bibliography{AP_library}

\clearpage

\renewcommand\thefigure{\Alph{section}.\arabic{figure}} \setcounter{table}{0}
\renewcommand\thetable{\Alph{section}.\arabic{table}}\setcounter{equation}{0}
\numberwithin{equation}{section} \numberwithin{figure}{section} %
\numberwithin{table}{section}

\appendix

\section{Appendix}


\subsection{The Average Period of Production and Trump 1.0 tariffsl \label{App: APP}}

Figure \ref{fig:app_trade_war} provides evidence
for a novel channel contributing to the gradual nature of the U.S.--China
decoupling. More specifically, building on a measure of the duration of
production developed by \cite{antras2025measuring} -- the so-called average
period of production or $\mathcal{APP}$ -- Figure \ref{fig:app_trade_war}
shows that reallocation of Chinese exports away from the U.S. was faster in
products with a shorter $\mathcal{APP}$, while products with longer
production cycles adjusted more slowly. Indeed, the figure shows that the
initial jump in Chinese exports to countries other than the U.S. after the
introduction of the first U.S. tariffs in January 2018 was entirely driven
by ``fast to produce'' goods.

\begin{figure}[htbp]
\caption{The Great Reallocation and the Temporal Dimension of Production}
\label{fig:app_trade_war}
\begin{center}
{\includegraphics[width=.7\textwidth]{Figures/Trade_War_Fast_Slow_1.pdf}}
\end{center}
\par
\begin{minipage}{.85\linewidth}
{\footnotesize \textbf{Notes:} Monthly exports are computed as backward-looking 6-month moving averages. We define goods as `fast to produce' if their $\mathcal{APP}$ is in the top quartile of the distribution and we label the others as ``slow to produce''.\par

\textbf{Sources:} Trade Data Monitor, \cite{Bailey-etal_2017}, \cite{antras2025measuring}. }
\end{minipage}
\end{figure}

\subsection{Microfoundation for our Indices: An Armington-CES\ Model \label{App: MicroFound}}

\paragraph{Environment}

The world consists of a finite set of countries indexed by $o,d\in \mathcal{W%
}$ and a finite set of products indexed by $p\in \mathcal{P}$. We interpret $%
d$ as a destination (importing) market, while $o$ denotes an origin
(exporting) country. Countries may act as both importers and exporters. At
times we will refer explicitly to the U.S. (indexed by $u$) and China
(indexed by $c$).

For each destination $d$ and product $p$, total nominal expenditure on
product $p$ in market $d$ is denoted $E_{dp}$. Throughout, $E_{dp}$ is
treated as a constant in our comparative-static exercises below.

Each origin $o$ supplies a single Armington variety of each product $p$ to
each destination $d$ at a delivered price $p_{odp}$. These delivered prices
can be interpreted as incorporating production costs, trade costs, and
markups.

\paragraph{Preferences and Demand}

Fix a destination market $d$ and product $p$. Consumers in $d$ have CES
preferences over country-of-origin varieties of product $p$: 
\begin{equation}
Q_{dp}=\left( \sum_{i\in \mathcal{W}}q_{odp}^{\frac{\sigma _{p}-1}{\sigma
_{p}}}\right) ^{\frac{\sigma _{p}}{\sigma _{p}-1}},\qquad \sigma _{p}>1,
\label{eq: CES aggregator}
\end{equation}%
where $q_{odp}$ is quantity of the $i$-variety of product $p$ consumed in
market $m$, and $\sigma _{p}$ is the elasticity of substitution across
origins within product $p$. The CES price index for product $p$ in
destination $s$ is 
\begin{equation*}
P_{dp}=\left( \sum_{i\in \mathcal{W}}p_{odp}^{1-\sigma _{p}}\right) ^{\frac{1%
}{1-\sigma _{p}}},
\end{equation*}%
where remember that $p_{odp}$ is the delivered price of variety $o$ of good $%
p$ sold in market $d$.

Let $X_{opd}$ denote nominal expenditure (revenue) in destination $d$ on
product $p$ sourced from origin $o$: 
\begin{equation*}
X_{opd}\equiv p_{opd}q_{opd}.
\end{equation*}%
CES demand implies that the expenditure share of origin $o$ in destination $%
d $ and product $p$ is 
\begin{equation}
\chi _{opd}\equiv \frac{X_{opd}}{E_{dp}}=\frac{p_{opd}^{1-\sigma _{p}}}{%
\sum_{o^{\prime }\in \mathcal{W}}p_{o^{\prime }pd}^{1-\sigma _{p}}},\qquad
\sum_{i\in \mathcal{W}}\chi _{opd}=1.
\end{equation}%
Thus, bilateral revenues satisfy $X_{opd}=\chi _{opd}E_{pd}$.

\paragraph{A Price Shock and Third-Country Export Losses}

We are interested in how an exporter-specific price change in a
`destination--product' market reallocates expenditure across origins, and
thereby affects third-country export revenues.

Fix a particular origin $c$ (e.g.,\ China) and consider a small proportional
decline in its delivered price in a given destination $d$ and product $p$: 
\begin{equation}
d\ln p_{cdp}=-\Phi _{p},\qquad \Phi _{p}>0,
\label{eq:shock}
\end{equation}%
holding fixed (i) total expenditure $E_{dp}$ and (ii) the delivered prices
of all other origins, $p_{opm}$ for $o\neq c$. The shock may be interpreted
as a cost decrease, a relative wage decline, or an outward supply shift for
exporter $c$ in market $(p,d)$. We shall return to the latter interpretation
in later subsection. Because we only consider this bilateral shock, from the
point of view of third markets, we can safely refer to $\varepsilon _{cdp}$
as product-level shock, which we denote by $\Phi _{p}$.

Let $o$ denote a third-country exporter potentially displaced by $c$ in
market $(p,d)$. Since $X_{odp}=E_{dp}\chi _{odp}$, changes in $X_{opd}$
arise through changes in $\chi _{opd}$ induced by \eqref{eq:shock}.
Differentiating the CES share formula yields the standard identity: 
\begin{equation}
dX_{odp}=-(\sigma _{p}-1)\times \Phi _{p}\times E_{dp}\times \chi
_{odp}\times \chi _{cpd}.  \label{eq:key_identity}
\end{equation}

Equation \eqref{eq:key_identity} captures the crowding out effect of the
increased Chinese exports, i.e., the revenue loss of exporter $o$ in a given
market. This loss is proportional to (i) the substitution elasticity $\sigma
_{p}-1$, (ii) the size of the shock $\Phi _{p}$ (ii) the size of the market $%
E_{dp}$, and (iii) the product of the two competitors' initial expenditure
shares, $\chi _{odp}\times \chi _{cpd}$. Holding fixed the elasticity and
market size, third-country losses are largest in destination--product
markets where (a) exporter $o$ is important (high $\chi _{opd}$) and (b)
exporter $c$ is also important (high $\chi _{cpd}$). If either exporter has
negligible market share, there is little (first-order) displacement in that
market.

\paragraph{Aggregation Across Products and Destinations}

Define total exports of country $o$ as the sum of revenues across all
destinations and products: 
\begin{equation*}
X_{o}\equiv \sum_{d\in \mathcal{W}}\sum_{p\in \mathcal{P}}X_{opd}.
\end{equation*}%
Summing \eqref{eq:key_identity} across $(p,d)$ yields the first-order change
in total exports: 
\begin{equation}
dX_{o}=-\sum_{d\in \mathcal{W}}\sum_{p\in \mathcal{P}}(\sigma _{p}-1)\times
\Phi _{p}\times E_{dp}\times \chi _{cpd}\times \chi _{odp}.
\label{eq:total_change}
\end{equation}

To express \eqref{eq:total_change} in terms of observable trade shares,
define: 
\begin{align}
X_{od}& \equiv \sum_{p\in \mathcal{P}}X_{opd}\quad & & \text{(exports of $o$
to destination d)}, \\
s_{od}& \equiv \frac{X_{od}}{X_{o}}\quad & & \text{(destination share in $o$%
's total exports)}, \\
s_{odp}& \equiv \frac{X_{odp}}{X_{od}}\quad & & \text{(product share within $%
o$'s exports to }d\text{)}.
\end{align}%
Note that $X_{opd}=X_{o}\cdot s_{od}\cdot s_{odp}$, and that $\chi
_{cdp}=X_{cdp}/E_{dp}$ is the share of destination $d$'s expenditure on
product $p$ sourced from China.

\paragraph{The Export Competition Index}

Dividing the change in country $o$'s exports by the initial level of
exports, we obtain%
\begin{equation*}
\frac{dX_{o}}{X_{o}}=-\sum_{d\in \mathcal{W}}\frac{X_{od}}{X_{o}}\sum_{p\in 
\mathcal{P}}(\sigma _{p}-1)\times \Phi _{p}\times \chi _{cpd}\times s_{odp}.
\end{equation*}%
If the elasticity of demand does not vary across products, then notice that 
\begin{equation}
\boldsymbol{ECI}_{oc}=\sum_{d\in \mathcal{W}}\frac{X_{od}}{X_{o}}\sum_{p\in 
\mathcal{P}}\Phi _{p}\times \chi _{cpd}\times s_{odp},
\label{eq: ECI Appendix}
\end{equation}%
which corresponds to equation (\ref{eq: ECI Baseline}) in the main text,
except for the set $\mathcal{W}$ in the first summation. A\ key aspect of $%
\boldsymbol{ECI}_{oc}$ is that it satisfies:%
\begin{equation*}
\frac{dX_{o}/X_{o}}{dX_{o^{\prime }}/X_{o^{\prime }}}=\frac{\boldsymbol{ECI}%
_{oc}}{\boldsymbol{ECI}_{o^{\prime }c}},
\end{equation*}%
so the $\boldsymbol{ECI}_{oc}$ is a useful index of the \emph{relative}
impact of the shock on the aggregate exports across countries.

\paragraph{Empirical Implementation for a Given $\Phi _{p}$}

Suppose we have a proxy for the size of the price shock in sector $p$, which
remember reduces Chinese export prices to all destinations by a percentage $%
\Phi _{p}$. We can then compute an index of the loss of export sales in all
destinations other than U.S. and China by implementing the formula above for
the case in which $\mathcal{W=W}_{-}$, where $\mathcal{W}_{-}$ is the set of
countries in the world except for the U.S. and China.

Notice that because the formula includes exports only, the summation also
omits country $o$'s itself. Relatedly, $X_{o}$ in the formula is the sum of
all of country $o$'s exports to all destinations other than the U.S. and
China. As for $\chi _{cpd}$, in principle this should be computed as the
share of total absorption of good $p$ in destination market $d$ accounted
for by China, but we shall approximate it by the share of overall imports of
good $p$ by country $d$ that originate in China. We do so because we do not
have access to information on the domestic sales of producers of good $p$ in 
$d$. A\ way to theoretically rationalize this empirical strategy would be to
stipulate preferences as being separable in domestic varieties and imported
ones, so that all the relevant CES\ aggregators sum only over foreign
varieties. In such a case, the model would dictate that $\chi _{cpd}$ should
indeed be computed as the ratio of $d$'s imports of good $p$ from China over 
$d$'s aggregate imports of good $p$.

\paragraph{Opening the Black Box of $\Phi _{p}$}

We next provide a structural interpretation of the reduced-form China export
price shock $\Phi _{p}$ introduced in the main analysis. In particular, we
show how a U.S. ad-valorem tariff on Chinese imports of a given price $p$
leads to a common decline in China's export price across non-U.S.
destinations when Chinese export supply is upward sloping.

To begin, consider a single product $p$ exported by China to a set of
destination countries $d\in \mathcal{W}$, including the United States ($d=U$%
). China charges a common producer (FOB) price $p_{cp}$ across destinations,
while delivered prices are given by 
\begin{equation*}
\tilde{p}_{cdp}=(1+\tau _{cdp})p_{cp},
\end{equation*}%
where (for simplicity) only the United States imposes a tariff: 
\begin{equation*}
\tau _{cup}>0,\qquad \tau _{cdp}=0\text{ for }d\neq U.
\end{equation*}

Demand in each destination follows the Armington-CES structure described
above. As in the rest of the Appendix, we treat destination expenditure on
product $p$, $E_{dp}$, as locally fixed. Under Armington-CES demand in
destination $d$ for product $p$, China's expenditure share is 
\begin{equation*}
\chi _{cdp}\equiv \frac{\tilde{p}_{cdp}q_{cdp}}{E_{dp}},
\end{equation*}%
and the exact own-price elasticity of demand for the Chinese variety
(holding competitor's prices fixed) is 
\begin{equation*}
\varepsilon _{cdp}\equiv -\frac{d\ln q_{cdp}}{d\ln \tilde{p}_{cdp}}=\sigma
_{p}-(\sigma _{p}-1)\chi _{cdp}.
\end{equation*}%
We thus have that 
\begin{equation*}
d\ln q_{cmp}=-\varepsilon _{cdp}d\ln \tilde{p}_{cmp}.
\end{equation*}%
Aggregating Chinese export demand across destinations, we have that 
\begin{equation*}
d\ln Q_{cp}=\sum_{d\in \mathcal{W}-\{c\}}\lambda _{dp}d\ln
q_{cdp}=-\sum_{d\in \mathcal{W}-\{c\}}\phi _{cdp}\varepsilon _{cdp}d\ln 
\tilde{p}_{cmp},
\end{equation*}%
where 
\begin{equation*}
Q_{cp}\equiv \sum_{d\in \mathcal{W}-\left\{ c\right\} }q_{cdp},
\end{equation*}%
and destination quantity shares 
\begin{equation*}
\phi _{cdp}\equiv \frac{q_{cdp}}{Q_{cp}},\qquad \sum_{d\in \mathcal{W}%
-\left\{ c\right\} }\phi _{cdp}=1.
\end{equation*}%
Since $\tilde{p}_{cdp}=(1+\tau _{cdp})p_{cp}$, we have 
\begin{equation*}
d\ln \tilde{p}_{cdp}=d\ln p_{cp}+d\ln \left( 1+\tau _{cdp}\right) .
\end{equation*}%
Because only the U.S.\ tariff changes, we can write 
\begin{equation*}
d\ln Q_{cp}=-\bar{\varepsilon}_{p}d\ln p_{cp}-\phi _{cup}\varepsilon
_{cup}\,d\ln (1+\tau _{cup}),
\end{equation*}%
where 
\begin{equation*}
\bar{\varepsilon}_{cp}\equiv \sum_{d\in \mathcal{W}-\left\{ c\right\} }\phi
_{cdp}\varepsilon _{cdp}\quad \text{and}\quad \varepsilon _{cup}=\sigma
_{p}-(\sigma _{p}-1)\chi _{cup}.
\end{equation*}

Assume now that China faces an upward-sloping export supply curve for
product $p$, summarized by 
\begin{equation*}
d\ln p_{cp}=\frac{1}{\eta _{p}}d\ln Q_{cp},
\end{equation*}%
where $\eta _{p}>0$ is the export supply elasticity, which in principle can
vary across products.

Substituting the export-demand aggregation into the supply relation and
solving yields 
\begin{equation}
d\ln p_{cp}=-\phi _{cup}\frac{\varepsilon _{cup}}{\eta _{p}+\bar{\varepsilon}%
_{cp}}d\ln (1+\tau _{cup}).  \label{eq: Price Change App}
\end{equation}%
How FOB\ Chinese prices react to the U.S. tariff depend on five key objects:
(i) the share of Chinese exports (in quantity terms) that is shipped to the
United States, $\phi _{cup}$; (ii) the elasticity of demand faced by Chinese
exporters in the U.S., $\varepsilon _{cup};$ (iii) the elasticity of supply
of Chinese exports, $\eta _{p};$ (iv) a weighted-average elasticity of
demand faced by Chinese exporters worldwide, $\bar{\varepsilon}_{cp}$; and
(v) the size of the U.S. tariff.

\paragraph{Empirical Implementation}

To simplify matters, in the main text we implement equation (\ref{eq: Price
Change App}) for the special case in which the U.S. policy shock is common
across sectors ($\tau _{cup}=\tau $) demand parameters $\sigma _{p}$ and
supply elasticities $\eta _{p}$ are common across sectors, and China
constitutes a small enough share of absorption in all countries ($\chi
_{cdp}\simeq 0$), to approximate the demand elasticity faced by Chinese
exporters as given by $\sigma $. Under these assumptions, equation (\ref{eq:
Price Change App}) reduces to 
\begin{equation*}
d\ln p_{cp}=-\phi _{cup}\frac{\varepsilon }{\eta +\varepsilon }d\ln (1+\tau
),
\end{equation*}%
and thus the relative size of the shock in different sectors depends only on 
$\phi _{cup}=q_{cup}/Q_{cp}$. Finally, if unit values are approximately
common across destinations (FOB), $\phi _{cup}$ is well approximated by the
corresponding export share: 
\begin{equation*}
\frac{q_{cup}}{Q_{cp}}\simeq \frac{X_{cup}}{\sum_{\mathcal{W}-\left\{
c\right\} }X_{cdp}}\equiv \Phi _{p},
\end{equation*}%
which corresponds to equation (\ref{eq: Relative Shock}) in the main text.

\paragraph{An Import Competition Index}

We now consider the impact of the same shock on the domestic producers of a
destination market $d$. As stated in the main text, there is a clear
connection between this question and our derivations above, in the sense
that we are considering a special case of our shock in which the destination
market $d$ under consideration is identical to the origin country $o$. These
considerations have no bearing on the size of the shock $\Phi _{p}$
originating in the U.S., and they are neither relevant, \emph{in theory},
for how the shock impacts market $d$ in general, as proxied by the Chinese
absorption share $\chi _{cpd}$. The main difference is that the shock is
focused on the domestic market $d$ and that the relevant sales share is the
share of domestic sales accounted by sector $p$, or $s_{dp}=S_{dp}/\sum_{p^{%
\prime }\in \mathcal{P}}S_{dp^{\prime }}$. Overall, the change in domestic
sales, referred to as $S_{d}$, is given by%
\begin{equation*}
\frac{dS_{d}}{S_{d}}=-\sum_{p\in \mathcal{P}}(\sigma _{p}-1)\times \Phi
_{p}\times \chi _{cpd}\times s_{dp}.
\end{equation*}%
Ignoring cross-section heterogeneity in demand elasticities ($\sigma
_{p}\simeq \sigma $ for all $p$), we obtain our import competition index in
equation (\ref{eq: ICI Baseline}) of the main text, or%
\begin{equation*}
\boldsymbol{ICI}_{dc}=\sum_{p\in \mathcal{P}}\Phi _{p}\times \chi
_{cdp}\times s_{dp},
\end{equation*}%
which satisfies 
\begin{equation*}
\frac{dS_{d}/S_{d}}{dS_{d^{\prime }}/S_{d^{\prime }}}=\frac{\boldsymbol{ICI}%
_{dc}}{\boldsymbol{ICI}_{d^{\prime }c}}.
\end{equation*}

As argued in the main text, when implementing this measure, researchers
cannot typically rely on exact measures of $s_{dp}$, so in our empirical
analysis, we proxy $s_{dp}\,$with the share of overall exports of country $d$
that are accounted by exports of good $p$, as indicated in equation (\ref%
{eq: X_share}). For similar reasons, and as in the export competition index,
we continue to approximate $\chi _{cpd}$ with the share of China in country $%
d$'s imports of good $p$, as in equation (\ref{eq: X_share}) in the main
text.

\paragraph{An Import Price Decline Index}

In the main text, we stated that the percentage fall in import prices in a
given destination market $d$ can be approximated, up to a first order, with
the following expression%
\begin{equation*}
\boldsymbol{IPDI}_{dc}=\sum_{p\in \mathcal{P}}\Phi _{p}\times \chi
_{cdp}^{M}\times s_{dp}^{M},
\end{equation*}%
where $\Phi _{p}$ is identical to that in the expressions for $\boldsymbol{%
ECI}_{oc}$ and $\boldsymbol{ICI}_{dc}$ in equations (\ref{eq: ECI Baseline})
and (\ref{eq: ICI Baseline}), and where $\chi _{cdp}^{M}$ is the Chinese
share in total imports of good $p$ in destination $d$ (which was already our
proxy for $\chi _{cdp}$ when empirically implementing our export competition
index). The term $s_{dp}^{M}$ in the above equation for $\boldsymbol{IPDI}%
_{dc}$ is the share of country $d$'s imports accounted for by sector $p$.

We next offer a simple proof of this claim. Define the CES import price
index for product $p$ in destination $d$ as 
\begin{equation*}
P_{dp}^{M}\equiv \left( \sum_{o\in \mathcal{W-}\left\{ d\right\}
}p_{odp}^{1-\sigma _{p}}\right) ^{\frac{1}{1-\sigma _{p}}},
\end{equation*}%
so that $P_{dp}=\left( p_{ddp}^{1-\sigma _{p}}+\left( P_{dp}^{M}\right)
^{1-\sigma _{p}}\right) ^{1/\left( 1-\sigma _{p}\right) }$ is the overall
price index for good $p$ in country $d$.

The associated import expenditure shares are 
\begin{equation*}
\chi _{odp}^{M}\equiv \frac{p_{odp}^{1-\sigma _{p}}}{\sum_{i\in \mathcal{W-}%
\left\{ d\right\} }p_{idp}^{1-\sigma _{p}}},\qquad \sum_{o\in \mathcal{W-}%
\left\{ d\right\} }\chi _{odp}^{M}=1.
\end{equation*}%
A first-order log-differentiation of $P_{dp}^{M}$ yields the standard
Divisia result: 
\begin{equation*}
d\ln P_{dp}^{M}=\sum_{o\in \mathcal{W-}\left\{ d\right\} }\chi _{odp}d\ln
p_{odp}.
\end{equation*}

Consider now the bilateral shock in (\ref{eq:shock}): only China's delivered
price changes for product $p$ in market $d$, so $d\ln p_{cdp}=-\Phi _{p}$
and $d\ln p_{odp}=0$ for all $o\neq c$. It follows immediately that 
\begin{equation*}
d\ln P_{dp}^{M}=-\chi _{cdp}^{M}\Phi _{p}.
\end{equation*}

Next, let the aggregate import price index in destination $d$ be a
Cobb--Douglas aggregator over product-level import price indices, 
\begin{equation*}
P_{d}^{M}\equiv \prod_{p\in \mathcal{P}}\left( P_{dp}^{M}\right)
^{s_{dp}^{M}},\qquad \text{with}\qquad s_{dp}^{M}\equiv \frac{M_{dp}}{%
\sum_{p^{\prime }\in \mathcal{P}}M_{dp^{\prime }}},
\end{equation*}%
where $M_{dp}$ denotes total imports of product $p$ by destination $d$.%
\footnote{%
More generally, any smooth aggregator with expenditure weights $s_{dp}^{M}$
delivers the same first-order log change around the initial allocation.}
Log-differentiating gives 
\begin{equation*}
d\ln P_{d}^{M}=\sum_{p\in \mathcal{P}}s_{dp}^{M}d\ln P_{dp}^{M}.
\end{equation*}%
Substituting the product-level result then implies 
\begin{equation*}
d\ln P_{d}^{M}=-\sum_{p\in \mathcal{P}}s_{dp}^{M}\chi _{cdp}^{M}\Phi _{p}.
\end{equation*}%
Therefore, the (positive) percentage fall in the aggregate import price
index is, up to a first-order approximation, 
\begin{equation*}
-\frac{dP_{d}^{M}}{P_{d}^{M}}\approx -d\ln P_{d}^{M}=\sum_{p\in \mathcal{P}%
}\Phi _{p}\times \chi _{cdp}^{M}\times s_{dp}^{M},
\end{equation*}%
which is the expression stated in the text.

\paragraph{Intermediate Inputs}

We now extend the baseline Armington--CES environment to allow a subset of
products, $p\in \mathcal{P}^{I}\subseteq \mathcal{P}$, to be used as
intermediate inputs by downstream producers in each destination market $d$.
For each input product $p\in \mathcal{P}^{I}$, downstream producers in $d$
assemble a CES composite from origin-specific varieties, 
\begin{equation*}
Q_{dp}^{I}\equiv \left( \sum_{o\in \mathcal{W}}q_{odp}^{\frac{\sigma _{p}-1}{%
\sigma _{p}}}\right) ^{\frac{\sigma _{p}}{\sigma _{p}-1}},\qquad \sigma
_{p}>1,
\end{equation*}%
with associated CES price index 
\begin{equation*}
P_{dp}^{I}\equiv \left( \sum_{o\in \mathcal{W}}p_{odp}^{1-\sigma
_{p}}\right) ^{\frac{1}{1-\sigma _{p}}}.
\end{equation*}%
Let $E_{dp}^{I}$ denote total expenditure by downstream producers in
destination $d$ on input $p$. Conditional on $E_{dp}^{I}$, the optimal
sourcing shares are identical to those derived above: 
\begin{equation*}
\chi _{odp}^{I}\equiv \frac{p_{odp}^{1-\sigma _{p}}}{\sum_{i\in \mathcal{W}%
}p_{idp}^{1-\sigma _{p}}},\qquad X_{odp}^{I}=\chi _{odp}^{I}\,E_{dp}^{I}.
\end{equation*}%
Domestic intermediate input sales in sector $p$ are therefore $S_{dp}\equiv
X_{ddp}^{I}=\chi _{ddp}^{I}E_{dp}^{I}$.

\paragraph{A Sourcing Competition Index}

Consider again the bilateral shock in \ref{eq:shock}, so that $d\ln
p_{cdp}=-\Phi _{p}$ and $d\ln p_{odp}=0$ for all $o\neq c$, and take a
first-order approximation around an initial allocation holding $E_{dp}^{I}$
fixed.\footnote{%
This \textquotedblleft local\textquotedblright\ closure obtains exactly, for
example, if downstream production uses a Cobb--Douglas aggregator over
intermediate inputs so that expenditures on each input are locally
proportional to downstream scale.} Since $\chi
_{ddp}^{I}=(p_{ddp}/P_{dp}^{I})^{1-\sigma _{p}}$, we have 
\begin{equation*}
d\ln S_{dp}=d\ln \chi _{ddp}^{I}=(\sigma _{p}-1)d\ln P_{dp}^{I}=-(\sigma
_{p}-1)\chi _{cdp}^{I}\Phi _{p}.
\end{equation*}%
Aggregating over inputs, let $S_{d}^{I}\equiv \sum_{p\in \mathcal{P}%
^{I}}S_{dp}$ and define the domestic intermediate-sales weights $%
s_{dp}^{I}\equiv S_{dp}/S_{d}^{I}$. A first-order log change in aggregate
domestic intermediate input sales is then 
\begin{equation*}
\frac{dS_{d}^{I}}{S_{d}^{I}}=-\sum_{p\in \mathcal{P}^{I}}(\sigma
_{p}-1)\times \Phi _{p}\times \chi _{cdp}^{I}\times s_{dp}^{I}.
\end{equation*}%
Ignoring cross-sectional heterogeneity in $\sigma _{p}$ and using the fact
that good $p$ is either a final good or an input so we can simply write $%
\chi _{cdp}^{I}=\chi _{cdp}$, we end uo with the sourcing competition index $%
\boldsymbol{SCI}_{dc}$ in (\ref{eq: SCI Formula}), such that 
\begin{equation*}
\frac{dS_{d}^{I}/S_{d}^{I}}{dS_{d^{\prime }}^{I}/S_{d^{\prime }}^{I}}=\frac{%
\boldsymbol{SCI}_{dc}}{\boldsymbol{SCI}_{d^{\prime }c}}\text{.}
\end{equation*}

Because we do not have data on domestic sales at a disaggregated product
level, when empirically implementing the index $\boldsymbol{SCI}_{dc}$, we
proxy the domestic share $s_{dp}^{I}$with the share of overall intermediate
input exports of country $d$ that are accounted by exports of input $p$, ,
as in equation (\ref{eq: X_share}) in the main text.

\paragraph{A Sourcing Price Decline Index}

Next, define the CES import price index for intermediate input $p\in 
\mathcal{P}^{I}$ in destination $d$ as 
\begin{equation*}
P_{dp}^{M,I}\equiv \left( \sum_{o\in \mathcal{W}\setminus
\{d\}}p_{odp}^{1-\sigma _{p}}\right) ^{\frac{1}{1-\sigma _{p}}}.
\end{equation*}%
Under the same China-only shock, $d\ln P_{dp}^{M,I}=-\chi _{cdp}^{M,I}\Phi
_{p}$, where $\chi _{cdp}^{M,I}$ is China's expenditure share in destination 
$d$'s imports of input $p$. Aggregating across intermediate inputs with a
Cobb--Douglas index, and assuming no cross-sectional variation in demand
elasticities,  
\begin{equation*}
P_{d}^{M,I}\equiv \prod_{p\in \mathcal{P}^{I}}\left( P_{dp}^{M,I}\right)
^{s_{dp}^{M,I}},\qquad s_{dp}^{M,I}\equiv \frac{M_{dp}^{I}}{\sum_{p^{\prime
}\in \mathcal{P}^{input}}M_{dp^{\prime }}^{I}},
\end{equation*}%
implies the first-order result 
\begin{equation*}
-d\ln P_{d}^{M,I}=\sum_{p\in \mathcal{P}^{I}}\Phi _{p}\times \chi
_{cdp}^{M,I}\times s_{dp}^{M,I}.
\end{equation*}%
Because a good $p$ is either a final good or an input so we can simply write 
$\chi _{cdp}^{M,I}=\chi _{cdp}^{M}$. We then have that for our empirical
proxy of $\Phi _{p}$, we have the following sourcing price decline index 
\begin{equation*}
\boldsymbol{SPDI}_{dc}=\sum_{p\in \mathcal{P}^{I}}\Phi _{p}\times \chi
_{cdp}^{M}\times s_{dp}^{M,I},
\end{equation*}%
which corresponds to equation (\ref{eq: SPDI index}) in the main text.

\paragraph{A Rerouting Potential Index}

We next construct an index that captures the extent to which country $i$'s
exports to the United States ($u$) increase in response to a reduction in
the price of intermediate inputs imported from China. Unlike the export
competition index, which operates through a demand-side shock in third
markets, here the mechanism is a supply-side cost shock: a fall in the price
of Chinese intermediate inputs lowers marginal costs in $i$ and thus reduces
the price of the varieties produced by $i$.

To allow for product-level heterogeneity in the marginal-cost reduction, let
firms producing output good $p$ in country $i$ use a Cobb--Douglas bundle of
intermediate inputs $z\in \mathcal{P}^{I}$ with cost shares $a_{pz}$, so
that the relevant intermediate input price index is 
\begin{equation*}
P_{ip}^{I}\equiv \prod_{z\in \mathcal{P}^{I}}\left( P_{iz}^{M,I}\right)
^{a_{pz}},\qquad \sum_{z\in \mathcal{P}^{I}}a_{pz}=1,
\end{equation*}%
where $P_{iz}^{M,I}$ is the CES import price index for intermediate input $z$
in country $i$ as defined above. Under the same China-only shock, we have $%
d\ln P_{iz}^{M,I}=-\chi _{ciz}^{M}\Phi _{z}$, where $\chi _{ciz}^{M}$ is
China's import expenditure share in country $i$ for input $z$. It follows
that the first-order change in the intermediate input price index relevant
for producing output $p$ is 
\begin{equation*}
-d\ln P_{ip}^{I}=\sum_{z\in \mathcal{P}^{I}}a_{pz}\times \Phi _{z}\times
\chi _{ciz}^{M}.
\end{equation*}%
Assuming that this reduction in input costs translates one-for-one into a
reduction in the producer price of the $i$-variety of product $p$, Armington
CES demand in the U.S. implies that, holding the U.S. product-level price
index fixed, 
\begin{equation*}
\frac{dX_{iup}}{X_{iup}}=-(\sigma _{p}-1)d\ln P_{ip}^{I}.
\end{equation*}%
Aggregating across products with baseline export shares to the U.S., $%
s_{iup}\equiv X_{iup}/X_{iu}$, we obtain 
\begin{equation*}
\frac{dX_{iu}}{X_{iu}}=-\sum_{p\in \mathcal{P}}s_{iup}\times (\sigma
_{p}-1)\times d\ln P_{ip}^{I}.
\end{equation*}%
If the elasticity of demand does not vary across products, then notice that
the rerouting potential index can be written as 
\begin{equation}
\boldsymbol{RPI}_{ciu}\equiv \sum_{p\in \mathcal{P}}s_{iup}\left( \sum_{z\in 
\mathcal{P}^{I}}a_{pz}\times \Phi _{z}\times \chi _{ciz}^{M}\right) .
\label{eq: RPI Appendix}
\end{equation}

If demand elasticities and pass-through are common across products and
countries, then $\boldsymbol{RPI}_{ciu}$ is proportional to the predicted
percentage change in country $i$'s exports to the U.S., and it satisfies 
\begin{equation*}
\frac{dX_{iu}/X_{iu}}{dX_{i^{\prime }u}/X_{i^{\prime }u}}=\frac{\boldsymbol{%
RPI}_{ciu}}{\boldsymbol{RPI}_{ci^{\prime }u}}\text{.}
\end{equation*}

\paragraph{Empirical implementation of the RPI without an Input--Output table%
}

In principle, one would ideally discipline the coefficients $a_{pz}$ using a
product-level input--output (IO) table that maps output goods into their
intermediate input requirements. In the absence of such data (or when the IO
table cannot be consistently harmonized with HS trade classifications), we
propose a set of trade-based approximations that impose structure on $a_{pz}$
using HS nesting. Specifically, for $\ell \in \{6,4,2\}$, define the HS-$%
\ell $ \textquotedblleft neighborhood\textquotedblright\ of output good $p$
as 
\begin{equation*}
\mathcal{P}^{I,(\ell )}(p)\equiv \left\{ z\in \mathcal{P}^{I}:\ HS_{\ell
}(z)=HS_{\ell }(p)\right\} ,
\end{equation*}%
where $HS_{6}(p)$ is the full HS-6 code (so that $\mathcal{P}^{I,(6)}(p)$
selects only $z=p$), while $HS_{4}(p)$ and $HS_{2}(p)$ denote the HS-4 and
HS-2 aggregations of $p$.

We then proxy the (unknown) IO weights $a_{pz}$ by country $i$'s baseline
intermediate import shares within the neighborhood of $p$: 
\begin{equation*}
s_{iz\mid p}^{M,I,(\ell )}\equiv \frac{M_{iz}^{I}}{\sum_{z^{\prime }\in 
\mathcal{P}^{I,(\ell )}(p)}M_{iz^{\prime }}^{I}},\qquad z\in \mathcal{P}%
^{I,(\ell )}(p).
\end{equation*}%
This yields a product-specific sourcing price decline index (SPDI) at level $%
\ell $: 
\begin{equation*}
\boldsymbol{SPDI}_{icp}^{(\ell )}\equiv \sum_{z\in \mathcal{P}^{I,(\ell
)}(p)}\Phi _{z}\times \chi _{ciz}^{M}\times s_{iz\mid p}^{M,I,(\ell )}.
\end{equation*}%
The three values of $\ell $ correspond to progressively broader assumptions
about input linkages: HS-6 restricts inputs to the same HS-6 code as the
output (an extreme \textquotedblleft own-input\textquotedblright\
assumption), HS-4 allows inputs within the same HS-4 heading, and HS-2
allows inputs within the same HS-2 chapter. These alternatives can be
interpreted as a robustness set that brackets plausible input--output
linkages when the true $a_{pz}$ matrix is unavailable.

Replacing the IO-based object in (\ref{eq: RPI Appendix}) with these proxies
gives the empirically implementable rerouting potential index 
\begin{equation*}
\boldsymbol{RPI}_{ciu}^{(\ell )}\equiv \sum_{p\in \mathcal{P}}s_{iup}\times 
\boldsymbol{SPDI}_{icp}^{(\ell )}.
\end{equation*}%
Finally, notice that if one instead assumes a common economy-wide
intermediate input bundle across all outputs (i.e., $a_{pz}$ does not vary
with $p$ and is proportional to $s_{iz}^{M,I}$), then $\boldsymbol{SPDI}%
_{icp}^{(\ell )}$ collapses to the aggregate $\boldsymbol{SPDI}_{ic}$
defined above, and product-level composition becomes irrelevant for the
predicted percentage change in $X_{iu}$.

\section{Empirical Constructions \label{app:Index_Empirics}}

\paragraph{Trade Data. }

Trade data comes from the Trade Data Monitor (TDM), which contains monthly imports and exports at the six-digit commodity level. Monthly observations are collapsed to annual totals. Country codes are standardized into ISO 3-letter codes. We substitute missing flows when possible with mirrored observations from other countries. For example, if China does not report any imports from the US, but the US reports exports to China, we use the US export data to fill in the missing Chinese imports. After mirroring, we default to using import data as the measure of trade flows. For earlier trade data, namely data from 2017 used to construct trade-weighted average elasticities and earlier baseline shares for index robustness, we use \cite{Gaulier-Zignago_2010}. 

\paragraph{Elasticities.}

Elasticity estimates come from \cite{soderbery2018trade}. Demand elasticities $\sigma_{dp}$ are reported at the destination-HS four-digit level. Inverse-supply elasticities $\eta_{odp}$ are reported as the origin-destination-HS four-digit level. Elasticities are cross-walked from their 1992 vintage to the 2017 vintage used in the trade data. When reporting demand elasticities, the author truncated estimates to 131.05, we opt to do the same to our working estimates. In addition, we apply the same truncation to supply elasticities.\footnote{\cite{soderbery2018trade} drops supply elasticity in the top and bottom 0.5\% when reporting estimates, but this still keeps many supply elasticities in the order of magnitude of $10^{6}$. We therefore opt to truncate elasticities.} We take the average of inverse-supply across destinations using export values in 2017, to get an elasticity, $\eta_{op}$, that is at the origin-HS4 level. Using an earlier period avoids endogeneity between the trade weights and tariff shocks being studied. 

In addition, \cite{soderbery2018trade} reports standard errors for estimated elasticities. Let $\theta$ denote an estimated elasticity, and $\sigma_{\theta}$ the associated standard error. We drop any elasticities where $\theta/\sigma_{\theta}<1.96$. There are estimates which can not be distinguished from zero under standard significance conventions. 

Elasticities are reported with the 1992 HS vintage, we cross-walk there to 2017 HS vintage using the Concordance R package \citep{HSconcordance}. Note that the \cite{soderbery2018trade} elasticities are estimated from trade data over the period 1993-2007, and there is certainty scope for preferences and technology to change before our baseline period of 2023. Hence, we continue to report "basic" measures as they are free of any structure elasticities and simply use observed trade shares. 

\paragraph{Tariffs.}

Tariff data comes from \cite{rodriguez2025}. This data primarily contains tariff data for the US. We consider two tariff windows. First, a short difference surrounding the second Trump administration: January 2025 to September 2025 (the final period of available data). Second, a long tariff difference beginning from before the first Trump administration's tariff policies to present, January 2018 to September 2025. 

\paragraph{Additional Data.}

The SGI requires input-output data at the HS six-digit level. We use binary input-output connections provided by \cite{AIPNET}. One limitation of this data is that no data on the extent of input usage is provided, outside of the binary classification of whether an upstream HS code is used in the production of another downstream HS code. Therefore, we give a uniform weight to each input. 

The Upstream and Downstream variants of the CGI require classifying goods as inputs or final goods. To do so we use the upstreamness measure from \cite{antras2012measuring}, and classify any goods with an upstreamness greater than 2 as intermediate goods and less than 2 as final goods. The upstreamness measure is reported for Input-Output Commodity Codes which we crosswalk to HS codes with data provided by the BEA. 

\subsection{Data Coverage}

In Table \ref{tab: Coverage} we report the fraction of products for which we observe the relevant elasticities needed to compute the shock to Chinese export prices (\ref{Phi_p}). At the very least, these products are included in our analysis. For computing subsequent measures, we may need to drop additional trade flows because of additional missing parameters. For example, the ECI requires import demand elasticities at all export markets, should these be absent a trade flow needs to be excluded from the ECI construction. In Table \ref{tab: Coverage Flows} we report the fraction of trade flows where we observe the relevant price shock and import-demand elasticity. 

Another consequence of partial coverage is then when computing indices that take an average over products, e.g. the ECI, we must re-normalize export shares so that they sum to one over the set of products for which we have the relevant data. Otherwise, indices will be correlated with coverage statistics. Since some products need to be dropped, this leads to the possibility that the indices do not accurately reflect 

\begin{table}[H]
\begin{centering}
\caption{Product Coverage\label{tab: Coverage}}
\include{Figures/coverage_table.tex}\vspace{0.25cm}
\par\end{centering}
\vspace{-0.5cm}
\raggedright{}{\footnotesize{}\textbf{Note.} For each
'component' row, 'HS6 products' reports the number of HS codes where
we observe the component, '\% products' the share of products we observe
the component, and '\% trade value' the share of trade flows accounted
for by the observed products. 'All components' refers to HS codes that contain demand elasticity, supply elasticities, and tariff changes needed to compute Chinese price shocks.}{\footnotesize\par}
\end{table}


\begin{table}[H]
\begin{centering}
\caption{Trade Flow Coverage\label{tab: Coverage Flows}}
\include{Figures/flow_coverage_table.tex}\vspace{0.25cm}
\par\end{centering}
\vspace{-0.5cm}
\raggedright{}{\footnotesize{}\textbf{Note.} For each
'component' row, '\% trade flows' reports the fraction of trade flows where
we observe the component and '\% trade value' the share of trade flows accounted
for by the observed components.}{\footnotesize\par}
\end{table}



\section{Additional Figures and Tables \label{app:Figures_Empirics}}

\begin{table}[H]
\begin{doublespace}
\begin{centering}
\caption{Price Shocks\label{tab: price_shocks}}
{\tiny{}\include{Figures/phi_top20.tex}}\vspace{0.25cm}
\par\end{centering}
\end{doublespace}
\raggedright{}{\footnotesize{}\textbf{Note.} The term $\phi_{cup}$ is the share of Chinese exports of goods $p$ going to the US, $\sigma_{up}$ is the US import demand elasticity for good $p$, $\varepsilon_{cup}=\sigma_{dp}-\left(\sigma_{dp}-1\right)\times\chi_{cup}$ where $\chi_{cup}$ is the share of US import expenditure in good $p$ on goods coming from China, $\bar{\varepsilon}_{cp}$ is the export-weighed average $\varepsilon_{cup}$ faced by China across export markets, $\Delta\tilde{\tau}_{cup}\equiv \Delta \log (1+\tau_{cup})$ is the change in tariffs between January 2025 to September 2025, and the price shock is $d\log p_{cp}=\Delta\tilde{\tau}_{cup}\times\phi_{cup}\times\varepsilon_{cup}/(\bar{\varepsilon}_{cp}+\eta_{p})$. In the baseline we use trade data from 2023 to compute shares. For long-term price shocks, $d\log p_{cp}^{Long}$, we use trade data from 2017 to compute shares and use the tariff difference between January 2018 and September 2025. }{\footnotesize\par}
\end{table}


\begin{table}[htbp]
\caption{Across-Index Rank Correlation}
\label{tab:rank_corr_across}\centering
\include{Figures/cross_index_rank_corr.tex}\vspace{0.25cm}
\raggedright{}{\footnotesize{}\textbf{Note.} Each cell reports a Spearman's rank correlation coefficient between the country-specific indices referenced in the row and column.}{\footnotesize\par}
\end{table}

\begin{table}[htbp]
\caption{Within-Index Rank Correlation}
\label{tab:rank_corr_within}\centering
\include{Figures/within_index_rank_corr.tex}\vspace{0.25cm}
\raggedright{}{\footnotesize{}\textbf{Note.} For each index (labeled in the row), each cell reports a Spearman's rank correlation coefficient between the variants (labeled in the column) of the country-specific indices. 'Full' is the baseline index, 'Basic' is to the index without any trade elasticities or tariffs but the same set of products as 'Full', 'Long' is the baseline index but with tariff differences over the 2018-2025 period, and 'Basic Full' is the same as 'Basic' but extended to include the full set of products. Product coverage differs between 'Full' and 'Basic Full' because of missing elasticity estimates as discussed in Table \ref{tab: Coverage}.}{\footnotesize\par}
\end{table}


% \subsection{Geopolitical Power and the Sustainability of Free Trade with
% Multiple Countries\label{App: MultiCountry}}

% Consider the difference in welfare for a given country $j$ of agreeing or
% abiding by to a free trade agreement with some other country $k$, or
% maintaining (or reverting) to optimal Nash tariffs. This is 
% \begin{equation*}
% \Delta W^{j,k}\equiv W^{j}(\text{free trade})-W^{j}\left( \text{Nash with }%
% t^{j,k}(\mu ^{k}),t^{k,j}(\mu ^{j})\right) .
% \end{equation*}%
% In our model with linear demand and supply, this difference reduces to:%
% \begin{equation*}
% \Delta W^{j,k}=\frac{(\bar{\alpha}-\underline{\alpha })^{2}\lambda ^{j}}{%
% 2(\beta +b)(1+\lambda ^{j}/\lambda ^{k})^{2}}\left[ \frac{(1-\mu ^{j})\left(
% 3-\mu ^{j}+2\lambda ^{j}/\lambda ^{k}\right) }{\left( 2-\mu ^{j}+\lambda
% ^{j}/\lambda ^{k}\right) ^{2}}-\frac{(1-\mu ^{k})\left( 1+\mu ^{k}+2\lambda
% ^{j}/\lambda ^{k}\right) }{\left( \lambda ^{k}/\lambda ^{j}+2-\mu
% ^{k}\right) ^{2}}\right] .
% \end{equation*}%
% Straightforward differentiation indicates that the term in brackets in $%
% \Delta W^{j,k}$ is decreasing in $\mu ^{j}$ and $\lambda ^{j}/\lambda ^{k}$
% and increasing in $\mu ^{k}$. In particular%
% \begin{eqnarray*}
% \frac{\partial \left( \frac{(1-\mu ^{j})\left( 3-\mu ^{j}+2\lambda
% ^{j}/\lambda ^{k}\right) }{\left( 2-\mu ^{j}+\lambda ^{j}/\lambda
% ^{k}\right) ^{2}}\right) }{\partial \mu ^{j}} &=&-2\frac{\left( \lambda
% ^{j}/\lambda ^{k}+1\right) ^{2}}{\left( \lambda ^{j}/\lambda ^{k}-\mu
% ^{j}+2\right) ^{3}}<0 \\
% \frac{\partial \left( \frac{(1-\mu ^{j})\left( 3-\mu ^{j}+2\lambda
% ^{j}/\lambda ^{k}\right) }{\left( 2-\mu ^{j}+\lambda ^{j}/\lambda
% ^{k}\right) ^{2}}\right) }{\partial \left( \lambda ^{j}/\lambda ^{k}\right) }
% &=&-2\left( \lambda ^{j}/\lambda ^{k}+1\right) \frac{1-\mu ^{j}}{\left(
% \lambda ^{j}/\lambda ^{k}+2-\mu ^{j}\right) ^{3}}<0 \\
% \frac{\partial \left( \frac{(1-\mu ^{k})\left( 1+\mu ^{k}+2\lambda
% ^{j}/\lambda ^{k}\right) }{\left( \lambda ^{k}/\lambda ^{j}+2-\mu
% ^{k}\right) ^{2}}\right) }{\partial \mu ^{k}} &=&-2\left( \lambda
% ^{j}/\lambda ^{k}\right) ^{2}\mu ^{k}\frac{\left( \lambda ^{j}/\lambda
% ^{k}+1\right) ^{2}}{\left( \left( 2-\mu ^{k}\right) \left( \lambda
% ^{j}/\lambda ^{k}\right) +1\right) ^{3}}<0 \\
% \frac{\partial \left( \frac{(1-\mu ^{k})\left( 1+\mu ^{k}+2\lambda
% ^{j}/\lambda ^{k}\right) }{\left( \lambda ^{k}/\lambda ^{j}+2-\mu
% ^{k}\right) ^{2}}\right) }{\partial \left( \lambda ^{j}/\lambda ^{k}\right) }
% &=&2\left( \lambda ^{j}/\lambda ^{k}\right) \left( \left( \lambda
% ^{j}/\lambda ^{k}\right) +1\right) \left( 1-\mu ^{k}\right) \frac{\left(
% 2-\mu \right) \left( \lambda ^{j}/\lambda ^{k}\right) +\mu ^{k}+1}{\left(
% \left( 2-\mu \right) \left( \lambda ^{j}/\lambda ^{k}\right) +1\right) ^{3}}%
% >0
% \end{eqnarray*}

% \bigskip

% These comparative statics ensure that, as in Proposition \ref{Prop:
% Agreements} applying to our two-country model, there continues to exist a
% threshold $\widetilde{\mu }^{j}\in \left[ 0,1\right] $ such that if $\mu
% ^{j}>\widetilde{\mu }^{j}$, then country $j$ will find it beneficial to opt
% out of a free trade agreement with $k$, and the threshold $\widetilde{\mu }%
% ^{j}$ is necessarily decreasing in the relative size of country $j$ (i.e., $%
% \partial \widetilde{\mu }^{j}/\partial \left( \lambda ^{j}/\lambda
% ^{k}\right) <0$. This result again suggests that sufficiently large and
% powerful countries will tend to opt out of free trade.

% Moving beyond bilateral agreements, country $j$ will support a global free
% trade agreement with all countries in the world whenever $\Delta W^{j}>0,$
% where%
% \begin{equation}
% \Delta W^{j}=\lambda ^{j}\sum_{k}\frac{1}{(1+\lambda ^{j}/\lambda ^{k})^{2}}%
% \left[ \frac{(1-\mu ^{j})\left( 3-\mu ^{j}+2\lambda ^{j}/\lambda ^{k}\right) 
% }{\left( 2-\mu ^{j}+\lambda ^{j}/\lambda ^{k}\right) ^{2}}-\frac{(1-\mu
% ^{k})\left( 1+\mu ^{k}+2\lambda ^{j}/\lambda ^{k}\right) }{\left( \lambda
% ^{k}/\lambda ^{j}+2-\mu ^{k}\right) ^{2}}\right]  \label{eq: Welfare Change}
% \end{equation}%
% and a multilateral free trade agreement will be sustainable if $\Delta
% W^{j}>0$ for all $j$.

% To study how the correlation between size and power shapes the
% sustainability of global free trade, we turn to a simple numerical simulation

% \paragraph{Numerical Simulations}

% We consider $N=20$ countries. Each country $j$ has parameters $(\lambda
% ^{j},\mu ^{j})$ and experiences a welfare change $\Delta W^{j}$ given by
% equation (\ref{eq: Welfare Change}), with the summation taken over all
% partners $k\neq j$. In the simulation, we impose the parameter bounds 
% \begin{equation*}
% \lambda ^{j}\in \lbrack 0.8,1.25],\qquad \mu ^{j}\in \lbrack 0,0.5].
% \end{equation*}%
% For each simulated world (one draw of $(\lambda ^{j},\mu ^{j})_{j=1}^{N}$),
% we compute $\Delta W^{j}$ for every country and record (i) whether all
% countries gain ($\Delta W^{j}>0$ for all $j$) and (ii) the share of
% countries that gain.

% To vary the cross-country correlation between $\lambda$ and $\mu$ while
% keeping their marginal distributions uniform on the above intervals, we use
% a Gaussian copula:

% \begin{enumerate}
% \item We draw correlated standard normals $(Z_{\lambda }^{j},Z_{\mu }^{j})$
% with correlation parameter $\rho $.

% \item Convert to uniforms via percentiles: $U_{\lambda}^j=\Phi(Z_{%
% \lambda}^j) $ and $U_{\mu}^j=\Phi(Z_{\mu}^j)$.

% \item Map to the desired ranges: $\lambda^j=0.8+(1.25-0.8)U_{\lambda}^j$ and 
% $\mu^j=0+(0.5-0)U_{\mu}^j$.
% \end{enumerate}

% For each value of $\rho $, we run 10{,}000 independent draws (simulated
% worlds). We report results against the \emph{average realized Pearson
% correlation} $\mathrm{Corr}(\lambda ,\mu )$ within each $\rho $ grid point.
% Figure \ref{fig:app_twopanels} plots the estimated probability that all
% countries gain and the average share of countries that gain.

% \subsection{The Correlation Between Power and Size in the Data: Robustness 
% \label{App: Corr Robust}}

% In this Appendix, we show that the increasing correlation between
% geopolitical power and economic size is robust to the computation of
% Spearman rank correlations rather than the standard (Pearson) correlations.

% \begin{figure}[htbp]
% \centering
% \par
% \begin{subfigure}[t]{0.48\linewidth}
%     \centering
%     \includegraphics[width=\linewidth]{Figures/corr_gpi_gdp_1995_2020_both.pdf}
%     \caption{Correlation between Global Power Index \citep{moyer2024relative} and GDP}
%     \label{fig: GPI}
%   \end{subfigure}\hfill 
% \begin{subfigure}[t]{0.48\linewidth}
%     \centering
%     \includegraphics[width=\linewidth]{Figures/corr_weightedpower_gdpshare_both_dualaxis.pdf}
%     \caption{Correlation between International Power \citep{liu2025international} and GDP }
%     \label{fig: LiuYang}
%   \end{subfigure}
% \caption{Correlations between Geopolitical Power and Economic Size
% (1995-2020) }
% \label{fig: Correlation Both Data}
% \end{figure}

\end{document}


